\documentclass[11pt,a4paper]{article}
\usepackage[utf8]{inputenc}
\usepackage[margin=0.75in]{geometry}
\usepackage{amsmath,amssymb,amsfonts,mathtools,extarrows,esint}
\usepackage{graphicx}
\usepackage[export]{adjustbox}
\usepackage{tcolorbox}
\usepackage{enumitem}
\usepackage{xcolor}
\usepackage{hyperref}
\usepackage{booktabs}
\usepackage{array}
\usepackage{tocloft}

\definecolor{primary}{RGB}{31, 78, 121}
\definecolor{secondary}{RGB}{47, 109, 26}
\definecolor{accent}{RGB}{204, 51, 0}
\definecolor{boxbg}{RGB}{245, 247, 250}

\hypersetup{
    colorlinks=true,
    linkcolor=primary,
    urlcolor=primary,
}

\tcbset{
    solbox/.style={
        colback=boxbg,
        colframe=primary,
        fonttitle=\bfseries,
        coltitle=white,
        boxrule=0.8pt,
        arc=4pt,
        left=8pt,
        right=8pt,
        top=6pt,
        bottom=6pt
    }
}

% Clean Table of Contents styling
\cftsetrmarg{3.5em}
\cftsetpnumwidth{2.5em}
\renewcommand{\cftsecleader}{\cftdotfill{\cftdotsep}}
\renewcommand{\cftsecfont}{\small\normalfont}
\renewcommand{\cftsecpagefont}{\small\bfseries\color{primary}}

\title{\Huge\bfseries GATE EC: Control Systems\\[0.3em] \Large Previous Year Solved Questions Archive}
\author{\textbf{Athena Project Archive}}
\date{\today}

\begin{document}
\maketitle
\tableofcontents
\newpage


\section*{Question 1: Stability Analysis \& Routh-Hurwitz (GATE EC 2026)}
\addcontentsline{toc}{section}{Question 1: Stability Analysis \& Routh-Hurwitz (GATE EC 2026)}
\noindent \textbf{\color{primary}NAT} \hfill \textbf{\color{secondary}2 Marks}


\noindent Consider the unity negative feedback control system shown in the Figure. 
The value of gain $K (>0)$ at which the given system will remain marginally stable is \underline{\hspace{1.5cm}}. (Answer in integer)

\begin{center}\includegraphics[max width=0.85\linewidth]{images/img\_78.png}\end{center}

\begin{center}\includegraphics[max width=0.85\linewidth]{images/img\_78.png}\end{center}


\begin{tcolorbox}[solbox, title=Detailed Solution -- Question 1]
\textbf{Given:}

Open-loop: $\frac{K}{s(s+7)(s+11)}$, unity negative feedback.

\textbf{Step 1: Characteristic equation}

$s(s+7)(s+11) + K = 0$

$s^3 + 18s^2 + 77s + K = 0$

\textbf{Step 2: Routh-Hurwitz array}

\[
\begin{array}{c|cc}
s^3 & 1 & 77 \
s^2 & 18 & K \
s^1 & \frac{18\times77 - K}{18} & 0 \
s^0 & K &
\end{array}
\]

\textbf{Step 3: Marginal stability - $s^1$ row = 0}

$\frac{18\times77 - K}{18} = 0$

$K = 18\times77 = 1386$
\end{tcolorbox}

\vspace{1.5em}
\hrule
\vspace{1.5em}

\section*{Question 2: Compensators \& Controllers (GATE EC 2026)}
\addcontentsline{toc}{section}{Question 2: Compensators \& Controllers (GATE EC 2026)}
\noindent \textbf{\color{primary}MCQ} \hfill \textbf{\color{secondary}2 Marks}


\noindent The state and output equations for a control system are:

\[
\dot{x} = \begin{bmatrix} -4 & -1.5 \\ 4 & 0 \end{bmatrix} x + \begin{bmatrix} 2 \\ 0 \end{bmatrix} u
\]
; 
\[
y = \begin{bmatrix} 1.5 & 0.625 \end{bmatrix} x
\]

Which of the following expressions correctly represents the transfer function $\frac{Y(s)}{U(s)}$ of the system with zero initial conditions?


\begin{enumerate}[label=(\Alph*)]
    \item $\frac{3s}{s^2+4s-6}$
    \item $\frac{3s+5}{s^2+4s-6}$
    \item \textbf{(Correct)} $\frac{3s+5}{s^2+4s+6}$
    \item $\frac{3s}{s^2+4s+6}$
\end{enumerate}

\noindent \textbf{Correct Answer:} (C)

\begin{tcolorbox}[solbox, title=Detailed Solution -- Question 2]
\textbf{Given:}

\[
\dot{x} = \begin{bmatrix} -4 & -1.5 \\ 4 & 0 \end{bmatrix} x + \begin{bmatrix} 2 \\ 0 \end{bmatrix} u
\]

\[
y = \begin{bmatrix} 1.5 & 0.625 \end{bmatrix} x
\]

\textbf{Step 1: Transfer function formula}

$\frac{Y(s)}{U(s)} = C(sI - A)^{-1}B$

\textbf{Step 2: Compute (sI - A)}

\[
sI - A = \begin{bmatrix} s+4 & 1.5 \\ -4 & s \end{bmatrix}
\]

\textbf{Step 3: Determinant}

$\det(sI-A) = s(s+4) + 6 = s^2 + 4s + 6$

\textbf{Step 4: Inverse of (sI - A)}

\[
(sI-A)^{-1} = \frac{1}{s^2+4s+6}\begin{bmatrix} s & -1.5 \\ 4 & s+4 \end{bmatrix}
\]

\textbf{Step 5: Multiply $C(sI-A)^{-1} B$}

\[
(sI-A)^{-1}B = \frac{1}{s^2+4s+6}\begin{bmatrix} s & -1.5 \\ 4 & s+4 \end{bmatrix}\begin{bmatrix}2\\0\end{bmatrix} = \frac{1}{s^2+4s+6}\begin{bmatrix}2s\\8\end{bmatrix}
\]

\[
\frac{Y(s)}{U(s)} = \begin{bmatrix}1.5 & 0.625\end{bmatrix}\frac{1}{s^2+4s+6}\begin{bmatrix}2s\\8\end{bmatrix}
\]

$= \frac{1}{s^2+4s+6}\left(1.5 \times 2s + 0.625 \times 8\right)$

$= \frac{3s + 5}{s^2 + 4s + 6}$
\end{tcolorbox}

\vspace{1.5em}
\hrule
\vspace{1.5em}

\section*{Question 3: Root Locus Technique (GATE EC 2026)}
\addcontentsline{toc}{section}{Question 3: Root Locus Technique (GATE EC 2026)}
\noindent \textbf{\color{primary}MCQ} \hfill \textbf{\color{secondary}2 Marks}


\noindent For the control system shown in the Figure, the transfer function of a plant,
 $G(s) = \frac{1}{(s+1)(s+2)}$ is connected in cascade with a compensator 
$C(s) = K(s + \alpha)$, where $K$ and $\alpha$ are positive real valued constants. 

Which of the following pairs $(K, \alpha)$ represent the correct values for the closed loop system to have poles at $(-3 \pm j\sqrt{5})$?

\begin{center}\includegraphics[max width=0.85\linewidth]{images/img\_113.png}\end{center}

\begin{center}\includegraphics[max width=0.85\linewidth]{images/img\_113.png}\end{center}


\begin{enumerate}[label=(\Alph*)]
    \item $2, 3$
    \item \textbf{(Correct)} $3, 4$
    \item $2, 4$
    \item $3, 3$
\end{enumerate}

\noindent \textbf{Correct Answer:} (B)

\begin{tcolorbox}[solbox, title=Detailed Solution -- Question 3]
\textbf{Given:}

Plant: $G(s) = \dfrac{1}{(s+1)(s+2)}$

Compensator: $C(s) = K(s + \alpha)$

Desired closed-loop poles: $s = -3 \pm j\sqrt{5}$

\textbf{Step 1: Open-loop transfer function}

$L(s) = C(s)G(s) = \frac{K(s+\alpha)}{(s+1)(s+2)}$

Characteristic equation (unity feedback):

$1 + L(s) = 0 \Rightarrow (s+1)(s+2) + K(s+\alpha) = 0$

\textbf{Step 2: Expand and match with desired polynomial}

Desired poles: $s = -3 \pm j\sqrt{5}$

Desired characteristic polynomial:

$(s+3-j\sqrt{5})(s+3+j\sqrt{5}) = (s+3)^2 + 5 = s^2 + 6s + 14$

Expanding the characteristic equation:

$s^2 + 3s + 2 + Ks + K\alpha = 0$
$s^2 + (3+K)s + (2+K\alpha) = 0$

\textbf{Step 3: Equate coefficients}

$3 + K = 6 \Rightarrow K = 3$

$2 + K\alpha = 14 \Rightarrow 2 + 3\alpha = 14 \Rightarrow \alpha = 4$
\end{tcolorbox}

\vspace{1.5em}
\hrule
\vspace{1.5em}

\section*{Question 4: Basics, Block Diagrams \& SFGs (GATE EC 2026)}
\addcontentsline{toc}{section}{Question 4: Basics, Block Diagrams \& SFGs (GATE EC 2026)}
\noindent \textbf{\color{primary}MCQ} \hfill \textbf{\color{secondary}1 Mark}


\noindent A control system is shown in the Figure.
 Which option represents the correct transfer function of the system?

\begin{center}\includegraphics[max width=0.85\linewidth]{images/img\_33.png}\end{center}

\begin{center}\includegraphics[max width=0.85\linewidth]{images/img\_33.png}\end{center}


\begin{enumerate}[label=(\Alph*)]
    \item $\frac{1}{(s+4)^2}$
    \item \textbf{(Correct)} $\frac{1}{(s+4)}$
    \item $\frac{2}{(s+4)}$
    \item $\frac{1}{(s^2+8s+17)}$
\end{enumerate}

\noindent \textbf{Correct Answer:} (B)

\begin{tcolorbox}[solbox, title=Detailed Solution -- Question 4]
\textbf{Given:}

Each block has transfer function

$G(s)=\frac{1}{s+4}$

Let the output of the left summing junction be

$E(s)$

\textbf{Step 1: Output of upper block}

The upper block output is the system output:

$C(s)=\frac{1}{s+4}E(s)$

\textbf{Step 2: Output of lower block}

The lower block also receives the same input $E(s)$, hence its output is

$\frac{1}{s+4}E(s)$

\textbf{Step 3: Right summing junction}

At the right summing junction,

- Upper input is positive
- Lower input is negative

Therefore, $\frac{1}{s+4}E(s)-\frac{1}{s+4}E(s)=0$

Hence, the feedback signal becomes zero.

\textbf{Step 4: Left summing junction}

Thus,
$E(s)=R(s)$

Substituting into the output equation,

$C(s)=\frac{1}{s+4}R(s)$
\end{tcolorbox}

\vspace{1.5em}
\hrule
\vspace{1.5em}

\section*{Question 5: Basics, Block Diagrams \& SFGs (GATE EC 2025)}
\addcontentsline{toc}{section}{Question 5: Basics, Block Diagrams \& SFGs (GATE EC 2025)}
\noindent \textbf{\color{primary}MCQ} \hfill \textbf{\color{secondary}2 Marks}


\noindent Consider a system represented by the block diagram shown below. Which of the following signal flow graphs represent(s) this system? Choose the correct option(s). 

\begin{center}\includegraphics[max width=0.85\linewidth]{images/img\_44.png}\end{center}

\begin{center}\includegraphics[max width=0.85\linewidth]{images/img\_44.png}\end{center}

\begin{center}\includegraphics[max width=0.85\linewidth]{images/img\_59.png}\end{center}

\begin{center}\includegraphics[max width=0.85\linewidth]{images/img\_59.png}\end{center}


\begin{enumerate}[label=(\Alph*)]
    \item A
    \item \textbf{(Correct)} B
    \item C
    \item D
\end{enumerate}

\noindent \textbf{Correct Answer:} (B)

\begin{tcolorbox}[solbox, title=Detailed Solution -- Question 5]
Using mason's gain formula, given block diagram

Overall T.F $=\frac{Y(s)}{R(s)}=\frac{G_{1}+G_{2}}{1+G_{1}}$

verify using options

\[
\frac{\mathrm{Y}(\mathrm{s})}{\mathrm{R}(\mathrm{s})}=\frac{\mathrm{G}_{1}+\mathrm{G}_{2}}{1+\mathrm{G}_{1}+\mathrm{G}_{2}} \rightarrow
\]
 wrong

\[
\frac{\mathrm{Y}(\mathrm{s})}{\mathrm{R}(\mathrm{s})}=\frac{\mathrm{G}_{1}+\mathrm{G}_{2}}{1+\mathrm{G}_{1}} \rightarrow
\]
 correct

\[
\frac{\mathrm{Y}(\mathrm{s})}{\mathrm{R}(\mathrm{s})}=\frac{\mathrm{G}_{1}+\mathrm{G}_{2}}{1-\mathrm{G}_{1}-\mathrm{G}_{2}} \rightarrow
\]
 wrong

$\frac{Y(s)}{R(s)}=\frac{G_{1}+G_{2}}{1+G_{1}+G_{2}} \rightarrow$ wrong
\end{tcolorbox}

\vspace{1.5em}
\hrule
\vspace{1.5em}

\section*{Question 6: State Space Analysis (GATE EC 2025)}
\addcontentsline{toc}{section}{Question 6: State Space Analysis (GATE EC 2025)}
\noindent \textbf{\color{primary}MCQ} \hfill \textbf{\color{secondary}2 Marks}


\noindent Consider a system where ${x}_{1}({t}), {x}_{2}({t})$, and ${x}_{3}({t})$ are three internal state signals and $u(t)$ is the input signal. The differential equations governing the system are given by

\[
\frac{d}{d t}\left[\begin{array}{l}x_{1}(t) \\ x_{2}(t) \\ x_{3}(t)\end{array}\right]=\left[\begin{array}{ccc}2 & 0 & 0 \\ 0 & -2 & 0 \\ 0 & 0 & 0\end{array}\right]\left[\begin{array}{l}x_{1}(t) \\ x_{2}(t) \\ x_{3}(t)\end{array}\right]+\left[\begin{array}{l}1 \\ 1 \\ 1\end{array}\right] u(t)
\]
.

Which of the following statements is/are TRUE?


\begin{enumerate}[label=(\Alph*)]
    \item The signals $x_{1}(t), x_{2}(t)$, and $x_{3}(t)$ are bounded for all bounded inputs
    \item \textbf{(Correct)} There exists a bounded input such that at least one of the signals $x_{1}(t), x_{2}(t)$, and $x_{3}(t)$ is unbounded
    \item There exists a bounded input such that the signals $x_{1}(t), x_{2}(t)$, and $x_{3}(t)$ are unbounded
    \item The signals $x_{1}(t), x_{2}(t)$, and $x_{3}(t)$ are unbounded for all bounded inputs
\end{enumerate}

\noindent \textbf{Correct Answer:} (B)

\begin{tcolorbox}[solbox, title=Detailed Solution -- Question 6]
Here given "A" matrix is in the diagonal canonical form [D.C.F] hence

T.F 
\[
=\frac{\mathrm{k}_{1}}{(\mathrm{~s}-2)}+\frac{\mathrm{k}_{2}}{(\mathrm{~s}+2)}+\frac{\mathrm{k}_{3}}{\mathrm{~s}}
\]
Here poles $\Rightarrow \mathrm{s}=-2,+2,0$

 

\begin{center}\includegraphics[max width=0.85\linewidth]{images/img\_90.png}\end{center}

\begin{center}\includegraphics[max width=0.85\linewidth]{images/img\_90.png}\end{center}

Here one pole is in the right half of s-plane so at least one of the signals ${x}_{1}({t}), {x}_{2}({t})$ and ${x}_{3}({t})$ are unbounded.
\end{tcolorbox}

\vspace{1.5em}
\hrule
\vspace{1.5em}

\section*{Question 7: Stability Analysis \& Routh-Hurwitz (GATE EC 2025)}
\addcontentsline{toc}{section}{Question 7: Stability Analysis \& Routh-Hurwitz (GATE EC 2025)}
\noindent \textbf{\color{primary}MCQ} \hfill \textbf{\color{secondary}2 Marks}


\noindent Consider the polynomial

$p(s)=s^{5}+7 s^{4}+3 s^{3}-33 s^{2}+2 s-40$.

 Let $(L,I,R)$ be defined as follows. 

$L$ is the number of roots of $p(s)$ with negative real parts.

$I$ is the number of roots of $\mathrm{p}(\mathrm{s})$ that are purely imaginary. 

$R$ is the number of roots of $p(s)$ with positive real parts.

Which one of the following options is correct?


\begin{enumerate}[label=(\Alph*)]
    \item \textbf{(Correct)} $\mathrm{L}=2, \mathrm{I}=2$, and $\mathrm{R}=1$
    \item $\mathrm{L}=3, \mathrm{I}=2$, and $\mathrm{R}=0$
    \item $\mathrm{L}=1, \mathrm{I}=2$, and $\mathrm{R}=2$
    \item $\mathrm{L}=0, \mathrm{I}=4$, and $\mathrm{R}=1$
\end{enumerate}

\noindent \textbf{Correct Answer:} (A)

\begin{tcolorbox}[solbox, title=Detailed Solution -- Question 7]
Given $p(s)=s^{5}+7 s^{4}+3 s^{3}-33 s^{2}+2 s-40$

\[
\begin{array}{l|lrl}
s^{5} & 1 & 3 & 2 \\
s^{4} & 7 & -33 & -40 \\
s^{3} & \frac{54}{7} & \frac{54}{7} & 0 \\
s^{2} & -40 & -40 & 0 \\
s^{1} & 0(-80) & 0 & 0 \\
s^{0} & -40 &  &  \\
\end{array}
\]

Auxiliary equation [A.E] $-40 \mathrm{~s}^{2}-40=0$

A.E 
\[
\Rightarrow \mathrm{s}^{2}+1=0 \Rightarrow \mathrm{~s}^{2}=-1 \Rightarrow \mathrm{~s}= \pm 1 \mathrm{j}
\]

So number of imaginary roots $[\mathrm{I}]=2$

In the RH tabular, number of sign changes in the first column $=1$, hence number of poles in the right half of s-plane (or) roots with positive real parts $[R]=1$, number of poles in the left half of s-plane $[\mathrm{L}]=5-2-1=2$
\end{tcolorbox}

\vspace{1.5em}
\hrule
\vspace{1.5em}

\section*{Question 8: Compensators \& Controllers (GATE EC 2025)}
\addcontentsline{toc}{section}{Question 8: Compensators \& Controllers (GATE EC 2025)}
\noindent \textbf{\color{primary}MCQ} \hfill \textbf{\color{secondary}2 Marks}


\noindent Let $\mathrm{G}(\mathrm{s})=\frac{1}{10 \mathrm{~s}^{2}}$ be the transfer function of a secondorder system. A controller M(s) is connected to the system $\mathrm{G}(\mathrm{s})$ in the configuration shown below.

Consider the following statements.

(i) There exists no controller of the form $M(s)=\frac{K_{I}}{s}$, where $K_{I}$ is a positive real number, such that the closed loop system is stable.

(ii) There exists at least one controller of the form 
\[
\mathrm{M}(\mathrm{s})=\mathrm{K}_{\mathrm{P}}+\mathrm{s} \mathrm{K}_{\mathrm{D}}
\]
, where $\mathrm{K}_{\mathrm{P}}$ and $\mathrm{K}_{\mathrm{D}}$ are positive real numbers, such that the closed loop system is stable.

Which one of the following options is correct?

\begin{center}\includegraphics[max width=0.85\linewidth]{images/img\_28.png}\end{center}

\begin{center}\includegraphics[max width=0.85\linewidth]{images/img\_28.png}\end{center}


\begin{enumerate}[label=(\Alph*)]
    \item (i) is TRUE and (ii) is FALSE
    \item (i) is FALSE and (ii) is TRUE
    \item Both (i) and (ii) are FALSE
    \item \textbf{(Correct)} Both (i) and (ii) are TRUE
\end{enumerate}

\noindent \textbf{Correct Answer:} (D)

\begin{tcolorbox}[solbox, title=Detailed Solution -- Question 8]
Statement (i): Given integral controller, integral controller adds pole at origin result system becomes unstable hence for closed loop system stability no integral controller,

so statement (i) is true.

Statement (ii): Given PD controller, which adds zero result system becomes stable, hence statement (ii) is also true.
\end{tcolorbox}

\vspace{1.5em}
\hrule
\vspace{1.5em}

\section*{Question 9: Root Locus Technique (GATE EC 2025)}
\addcontentsline{toc}{section}{Question 9: Root Locus Technique (GATE EC 2025)}
\noindent \textbf{\color{primary}MCQ} \hfill \textbf{\color{secondary}1 Mark}


\noindent Consider the unity-negative-feedback system shown in Figure (i) below, where gain  $K \geq 0$. The root locus of this system is shown in Figure (ii) below. 

For what value(s) of  $K$ will the system in Figure (i) have a pole at  $-1+ji$?  

\begin{center}\includegraphics[max width=0.85\linewidth]{images/img\_101.png}\end{center}

\begin{center}\includegraphics[max width=0.85\linewidth]{images/img\_101.png}\end{center}


\begin{enumerate}[label=(\Alph*)]
    \item $\mathrm{K}=5$
    \item $\mathrm{K}=1 / 5$
    \item \textbf{(Correct)} For no positive value of $\mathrm{K}$
    \item For all positive values of $\mathrm{K}$
\end{enumerate}

\noindent \textbf{Correct Answer:} (C)

\begin{tcolorbox}[solbox, title=Detailed Solution -- Question 9]
We are given the root locus plot in Figure (ii), and we are to determine whether a closed-loop pole lies at the complex location $-1 + j1$ for any $K \geq 0$. 

From control theory, the root locus shows the trajectory of the closed-loop poles as $K$ varies from $0$ to $\infty$. 

We inspect the root locus diagram in Figure (ii): 
The root locus branches are symmetric about the real axis.

The point $-1 + j1$ lies outside the plotted root locus path.

It is not marked or touched by any of the root locus branches. 

Hence, the closed-loop poles never pass through the point $-1 + j1$ for any positive value of $K$.
\end{tcolorbox}

\vspace{1.5em}
\hrule
\vspace{1.5em}

\section*{Question 10: Frequency Response Analysis (GATE EC 2025)}
\addcontentsline{toc}{section}{Question 10: Frequency Response Analysis (GATE EC 2025)}
\noindent \textbf{\color{primary}MCQ} \hfill \textbf{\color{secondary}1 Mark}


\noindent The Nyquist plot of a system is given in the figure. Let $\omega_P, \omega_Q, \omega_R$ and $\omega_S$ be the positive frequencies at the points $P, Q, R,$ and $S,$ respectively. 

Which one of the following statements is TRUE?

\begin{center}\includegraphics[max width=0.85\linewidth]{images/img\_63.png}\end{center}

\begin{center}\includegraphics[max width=0.85\linewidth]{images/img\_63.png}\end{center}


\begin{enumerate}[label=(\Alph*)]
    \item $\omega_S$ is the gain crossover frequency and $\omega_P$ is the phase crossover frequency
    \item $\omega_Q$ is the gain crossover frequency and $\omega_S$ is the phase crossover frequency
    \item $\omega_R$ is the gain crossover frequency and $\omega_S$ is the phase crossover frequency
    \item \textbf{(Correct)} $\omega_S$ is the gain crossover frequency and $\omega_Q$ is the phase crossover frequency
\end{enumerate}

\noindent \textbf{Correct Answer:} (D)

\begin{tcolorbox}[solbox, title=Detailed Solution -- Question 10]
From the Nyquist plot: 

The gain crossover frequency is the frequency at which the Nyquist plot crosses the unit circle, i.e., the magnitude of the open-loop transfer function $|G(j\omega)H(j\omega)| = 1$. 

The phase crossover frequency is the frequency at which the phase is $-180^\circ$, i.e., the Nyquist plot crosses the negative real axis. 

From the diagram: 

Point S lies on the unit circle, so $\omega_S$ is the gain crossover frequency. 

Point P lies on the negative real axis, hence $\omega_P$ is the phase crossover frequency. 

Therefore, the correct statement is:
$\omega_S$ is the gain crossover frequency and $\omega_P$ is the phase crossover frequency.
\end{tcolorbox}

\vspace{1.5em}
\hrule
\vspace{1.5em}

\section*{Question 11: State Space Analysis (GATE EC 2024)}
\addcontentsline{toc}{section}{Question 11: State Space Analysis (GATE EC 2024)}
\noindent \textbf{\color{primary}MSQ} \hfill \textbf{\color{secondary}2 Marks}


\noindent Consider a system $S$ represented in state space as

 
\[
\frac{d x}{d t}=\left[\begin{array}{ll}
0 & -2 \\
1 & -3
\end{array}\right] x+\left[\begin{array}{l}
1 \\
0
\end{array}\right] r, y=\left[\begin{array}{ll}
2 & -5
\end{array}\right] x
\]

Which of the state space representations given below has/have the same transfer function as that of $S$ ?


\begin{enumerate}[label=(\Alph*)]
    \item \textbf{(Correct)} \[
\frac{d x}{d t}=\left[\begin{array}{cc}0 & 1 \\ -2 & -3\end{array}\right] x+\left[\begin{array}{l}1 \\ 0\end{array}\right] r, y=\left[\begin{array}{ll}1 & 2\end{array}\right] x
\]
    \item \[
\frac{d x}{d t}=\left[\begin{array}{cc}0 & 1 \\ -2 & -3\end{array}\right] x+\left[\begin{array}{l}1 \\ 0\end{array}\right] r, y=\left[\begin{array}{ll}0 & 2\end{array}\right] x
\]
    \item \textbf{(Correct)} \[
\frac{d x}{d t}=\left[\begin{array}{cc}-1 & 0 \\ 0 & -2\end{array}\right] x+\left[\begin{array}{c}-1 \\ 3\end{array}\right] r, y=\left[\begin{array}{ll}1 & 1\end{array}\right] x
\]
    \item \[
\frac{d x}{d t}=\left[\begin{array}{cc}-1 & 0 \\ 0 & -2\end{array}\right] x+\left[\begin{array}{l}1 \\ 1\end{array}\right] r, y=\left[\begin{array}{ll}1 & 2\end{array}\right] x
\]
\end{enumerate}

\noindent \textbf{Correct Answer:} (A, C)

\begin{tcolorbox}[solbox, title=Detailed Solution -- Question 11]
From given state model, (observable form)

 $T F=C\left[(s I-A)^{-1} B\right]+D$

 $=\frac{(2 s+1)}{s^{2}+3 s+2}$

From this controllable form can be written as

 
\[
\begin{aligned}
X^{0} & =\left[\begin{array}{cc}
0 & 1 \\
-2 & -3
\end{array}\right] X+\left[\begin{array}{l}
0 \\
1
\end{array}\right] u \\
Y & =\left[\begin{array}{ll}
1 & 2
\end{array}\right] X+[0] u
\end{aligned}
\]

Another possibility, $\quad \mathrm{TF}=\frac{2 s+1}{s^{2}+3 s+2}=\frac{-1}{s+1}+\frac{3}{s+2}$

 
\[
\begin{aligned}
X^{0} & =\left[\begin{array}{cc}
-1 & 0 \\
0 & -2
\end{array}\right] X+\left[\begin{array}{c}
-1 \\
3
\end{array}\right] v \\
Y & =\left[\begin{array}{ll}
1 & 1
\end{array}\right] X+[0] u
\end{aligned}
\]

Diagonal form of state mode.
\end{tcolorbox}

\vspace{1.5em}
\hrule
\vspace{1.5em}

\section*{Question 12: Compensators \& Controllers (GATE EC 2024)}
\addcontentsline{toc}{section}{Question 12: Compensators \& Controllers (GATE EC 2024)}
\noindent \textbf{\color{primary}MCQ} \hfill \textbf{\color{secondary}2 Marks}


\noindent A satellite attitude control system, as shown below, has a plant with transfer function $G(s)=\frac{1}{s^{2}}$ cascaded with a compensator $C(s)=\frac{K(s+\alpha)}{s+4}$, where $K$ and $\alpha$ are positive real constants.

\begin{center}\includegraphics[max width=0.85\linewidth]{images/img\_99.png}\end{center}

\begin{center}\includegraphics[max width=0.85\linewidth]{images/img\_99.png}\end{center}

In order for the closed-loop system to have poles at $-1 \pm j \sqrt{3}$, the value of $\alpha$ must be \underline{\hspace{1.5cm}}


\begin{enumerate}[label=(\Alph*)]
    \item 0
    \item \textbf{(Correct)} 1
    \item 2
    \item 3
\end{enumerate}

\noindent \textbf{Correct Answer:} (B)

\begin{tcolorbox}[solbox, title=Detailed Solution -- Question 12]
$G(S) C(s)=\frac{1}{S^{2}} \frac{k(S+\alpha)}{(S+4)}$

Characteristic equation is

 $S^{3}+4 S^{2}+k S+\alpha=0$

Poles of the system present at $=-1 \pm j \sqrt{3}$

Characteristic equation is

 
\[
\begin{aligned}
&(S+a)\left[(S+1)^{2}+3\right]=0 \\
&(S+a)\left(S^{2}+2 S+4\right)=0 \\
& S^{3}+(2+a) S^{2}+(4+2 a) S+4 a=0
\end{aligned}
\]

Compare equation (I) and (II)

 
\[
\begin{aligned}
2+a & =4 \\
a & =2 \\
4+2 a & =k \\
k & =8 \\
k \alpha & =4 a \\
\alpha & =\frac{8}{8}=1 \\
\alpha & =1
\end{aligned}
\]
\end{tcolorbox}

\vspace{1.5em}
\hrule
\vspace{1.5em}

\section*{Question 13: Time Response Analysis (GATE EC 2024)}
\addcontentsline{toc}{section}{Question 13: Time Response Analysis (GATE EC 2024)}
\noindent \textbf{\color{primary}MCQ} \hfill \textbf{\color{secondary}2 Marks}


\noindent Consider a unity negative feedback control system with forward path gain $G(s)=\frac{K}{(s+1)(s+2)(s+3)}$ as shown.
The impulse response of the closed-loop system decays faster than $e^{-t}$ if \underline{\hspace{1.5cm}} 

\begin{center}\includegraphics[max width=0.85\linewidth]{images/img\_32.png}\end{center}

\begin{center}\includegraphics[max width=0.85\linewidth]{images/img\_32.png}\end{center}


\begin{enumerate}[label=(\Alph*)]
    \item \textbf{(Correct)} $1 \leq K \leq 5$
    \item $7 \leq K \leq 21$
    \item $-4 \leq K \leq-1$
    \item $-24 \leq K \leq-6$
\end{enumerate}

\noindent \textbf{Correct Answer:} (A)

\begin{tcolorbox}[solbox, title=Detailed Solution -- Question 13]
Given,

 $G(s)=\frac{K}{(S+1)(S+2)(S+3)}$

Impulse response of system is

 $Y(S)=\frac{K}{(S+1)(S+2)(S+3)+K}$

Given: At impulse response of closed loop system decay faster than $e^{-t}$ Hence, put $s=s-1$,

$\therefore$ The new characteristic equation is,

 
\[
\begin{aligned}
& (s-1+1)(s-1+2)(s-1+3)+K=0 \\
& s(s+1)(s+2)+K=0 \\
& s^{3}+3 s^{2}+2 s+K=0
\end{aligned}
\]

By $\mathrm{RH}$ criterion,

 
\[
\begin{array}{c|cc}
s^{3} & 1 & 2 \\
s^{2} & 3 & K \\
s^{1} & \frac{6-K}{3} & \\
s^{0} & K &
\end{array}
\]

 $\therefore K > 0 ; \quad \frac{6-K}{3} > 0 \Rightarrow K < 6$

 $\therefore \quad 0 < K < 6$

Hence, option satisfies above range of $K$.
\end{tcolorbox}

\vspace{1.5em}
\hrule
\vspace{1.5em}

\section*{Question 14: Stability Analysis \& Routh-Hurwitz (GATE EC 2024)}
\addcontentsline{toc}{section}{Question 14: Stability Analysis \& Routh-Hurwitz (GATE EC 2024)}
\noindent \textbf{\color{primary}MCQ} \hfill \textbf{\color{secondary}1 Mark}


\noindent In the feedback control system shown in the figure below $G(s)=\frac{6}{s(s+1)(s+2)}$. 

\begin{center}\includegraphics[max width=0.85\linewidth]{images/img\_11.png}\end{center}

\begin{center}\includegraphics[max width=0.85\linewidth]{images/img\_11.png}\end{center}

$R(s), Y(s)$, and $E(s)$ are the Laplace transforms of $r(t), y(t)$, and $e(t)$, respectively. If the input $r(t)$ is a unit step function, then \underline{\hspace{1.5cm}}


\begin{enumerate}[label=(\Alph*)]
    \item $\lim _{t \rightarrow \infty} e(t)=0$
    \item $\lim _{t \rightarrow \infty} e(t)=\frac{1}{3}$
    \item $\lim _{t \rightarrow \infty} e(t)=\frac{1}{4}$
    \item \textbf{(Correct)} $\lim _{t \rightarrow \infty} e(t)$ does not exist, $e(t)$ is oscillatory
\end{enumerate}

\noindent \textbf{Correct Answer:} (D)

\begin{tcolorbox}[solbox, title=Detailed Solution -- Question 14]
Given system is

\begin{center}\includegraphics[max width=0.85\linewidth]{images/img\_45.png}\end{center}

\begin{center}\includegraphics[max width=0.85\linewidth]{images/img\_45.png}\end{center}

 $G(s)=\frac{6}{s(s+1)(s+2)}$

The characteristic equation,

 
\[
\begin{aligned}
1+G(s) & =0 \\
1+\frac{6}{s(s+1)(s+2)} & =0 \\
s(s+1)(s+2)+6 & =0 \\
s^{3}+3 s^{2}+2 s+6 & =0
\end{aligned}
\]

Clearly, here the internal coefficient product is equal to external coefficient product i.e.,

 $3 \times 2=6 \times 1$

Hence, the given system is marginally stable system.

$\therefore e(t)$ not exist because system is oscillatory system.
\end{tcolorbox}

\vspace{1.5em}
\hrule
\vspace{1.5em}

\section*{Question 15: Frequency Response Analysis (GATE EC 2024)}
\addcontentsline{toc}{section}{Question 15: Frequency Response Analysis (GATE EC 2024)}
\noindent \textbf{\color{primary}MCQ} \hfill \textbf{\color{secondary}1 Mark}


\noindent In the context of Bode magnitude plots, $40 \mathrm{~dB} /$ decade is the same as \underline{\hspace{1.5cm}}


\begin{enumerate}[label=(\Alph*)]
    \item \textbf{(Correct)} $12 \mathrm{~dB} /$ octave
    \item $6 \mathrm{~dB} /$ octave
    \item $20 \mathrm{~dB} /$ octave
    \item $10 \mathrm{~dB} /$ octave
\end{enumerate}

\noindent \textbf{Correct Answer:} (A)

\begin{tcolorbox}[solbox, title=Detailed Solution -- Question 15]
\[
\begin{aligned}
20 \times \mathrm{ndB} / \text { decade } & =6 \times \mathrm{n} \mathrm{dB} / \text { octave } \\
40 \mathrm{~dB} / \text { decade } & =12 \mathrm{~dB} / \text { octave }
\end{aligned}
\]
\end{tcolorbox}

\vspace{1.5em}
\hrule
\vspace{1.5em}

\section*{Question 16: Frequency Response Analysis (GATE EC 2023)}
\addcontentsline{toc}{section}{Question 16: Frequency Response Analysis (GATE EC 2023)}
\noindent \textbf{\color{primary}NAT} \hfill \textbf{\color{secondary}2 Marks}


\noindent The asymptotic magnitude Bode plot of a minimum phase system is shown in the figure.The transfer function of the system is $(s)=\frac{k(s+z)^{a}}{s^{b}(s+p)^{c}}$, where $k, z, p, a, b$ and $c$ are positive constants. The value of $(a+b+c)$ is \underline{\hspace{1.5cm}} 

(rounded off to the nearest integer).

\begin{center}\includegraphics[max width=0.85\linewidth]{images/img\_124.jpg}\end{center}

\begin{center}\includegraphics[max width=0.85\linewidth]{images/img\_124.jpg}\end{center}


\begin{tcolorbox}[solbox, title=Detailed Solution -- Question 16]
From the Bode magnitude plot, it is clear that there is one pole at origin,
 $\therefore$ $b=1$

 and at frequency
$\omega_{1}$, system has a zero

$\therefore$ $a=1$ 

 and at frequency
$\omega_{2}$, system has a two poles

\[
\begin{aligned}
\therefore \quad c&=2 \\
\therefore \quad a+b+c&=1+1+2 \\
a+b+c&=4
\end{aligned}
\]
\end{tcolorbox}

\vspace{1.5em}
\hrule
\vspace{1.5em}

\section*{Question 17: Basics, Block Diagrams \& SFGs (GATE EC 2023)}
\addcontentsline{toc}{section}{Question 17: Basics, Block Diagrams \& SFGs (GATE EC 2023)}
\noindent \textbf{\color{primary}MCQ} \hfill \textbf{\color{secondary}2 Marks}


\noindent In the following block diagram, $R(s)$ and $D(s)$ are two inputs. The output $Y(s)$ is expressed as $Y(s)=G_{1}(s) R(s)+G_{2}(s) D(s)$.

$G_{1}(s)$ and $G_{2}(s)$ are given by

\begin{center}\includegraphics[max width=0.85\linewidth]{images/img\_42.jpg}\end{center}

\begin{center}\includegraphics[max width=0.85\linewidth]{images/img\_42.jpg}\end{center}


\begin{enumerate}[label=(\Alph*)]
    \item \textbf{(Correct)} $G_{1}(s)=\frac{G(s)}{1+G(s)+G(s) H(s)}$ and $G_{2}(s)=\frac{G(s)}{1+G(s)+G(s) H(s)}$
    \item $G_{1}(s)=\frac{G(s)}{1+G(s)+H(s)}$ and $G_{2}(s)=\frac{G(s)}{1+G(s)+H(s)}$
    \item $G_{1}(s)=\frac{G(s)}{1+G(s)+H(s)}$ and $G_{2}(s)=\frac{G(s)}{1+G(s)+G(s) H(s)}$
    \item $G_{1}(s)=\frac{G(s)}{1+G(s)+G(s) H(s)}$ and $G_{2}(s)=\frac{G(s)}{1+G(s)+H(s)}$
\end{enumerate}

\noindent \textbf{Correct Answer:} (A)

\begin{tcolorbox}[solbox, title=Detailed Solution -- Question 17]
\begin{center}\includegraphics[max width=0.85\linewidth]{images/img\_68.jpg}\end{center}

\begin{center}\includegraphics[max width=0.85\linewidth]{images/img\_68.jpg}\end{center}

$Y(s)=\underbrace{G_{1}(s) R(s)}_{Y_{1}(s)}+\underbrace{G_{2}(s) D(s)}_{Y_{2}(s)}$Considering first $R(s)$ only, then $Y(s)$ is $Y_{1}(s)$

\begin{center}\includegraphics[max width=0.85\linewidth]{images/img\_48.jpg}\end{center}

\begin{center}\includegraphics[max width=0.85\linewidth]{images/img\_48.jpg}\end{center}

\begin{center}\includegraphics[max width=0.85\linewidth]{images/img\_84.jpg}\end{center}

\begin{center}\includegraphics[max width=0.85\linewidth]{images/img\_84.jpg}\end{center}

\[
\begin{aligned}
\frac{Y_{1}(s)}{R(s)}&=\frac{\frac{G(s)}{1+G(s) H(s)}}{1+\frac{G(s)}{1+G(s) H(s)}} \\
\frac{Y_{1}(s)}{R(s)}&=\frac{G(s)}{1+G(s) H(s)+G(s)} \\
Y_{1}(s)&=\left[\frac{G(s)}{1+G(s)+G(s) H(s)}\right] R(s) \\
G_{1}(s)&=\frac{G(s)}{1+G(s)+G(s) H(s)}
\end{aligned}
\]

Now considering $D(s)$ only, then $Y(s)$ is $Y_{2}(s)$

\begin{center}\includegraphics[max width=0.85\linewidth]{images/img\_18.jpg}\end{center}

\begin{center}\includegraphics[max width=0.85\linewidth]{images/img\_18.jpg}\end{center}

\[
\begin{aligned}
\frac{Y_{2}(s)}{D(s)} & =\frac{G(s)}{1+G(s)[1+H(s)]} \\
Y_{2}(s) & =\left[\frac{G(s)}{1+G(s)+G(s) H(s)}\right] D(s) \\
G_{2}(s) & =\frac{G(s)}{1+G(s)+G(s) H(s)}
\end{aligned}
\]

Hence, $G_{1}(s)$ and $G_{2}(s)$ both are equal.
\end{tcolorbox}

\vspace{1.5em}
\hrule
\vspace{1.5em}

\section*{Question 18: Time Response Analysis (GATE EC 2023)}
\addcontentsline{toc}{section}{Question 18: Time Response Analysis (GATE EC 2023)}
\noindent \textbf{\color{primary}MCQ} \hfill \textbf{\color{secondary}2 Marks}


\noindent A closed loop system is shown in the figure where $k > 0$ and $\alpha > 0$. The steady state error due to a ramp input $\left(R(s)=\alpha / s^{2}\right)$ is given by

\begin{center}\includegraphics[max width=0.85\linewidth]{images/img\_24.jpg}\end{center}

\begin{center}\includegraphics[max width=0.85\linewidth]{images/img\_24.jpg}\end{center}


\begin{enumerate}[label=(\Alph*)]
    \item \textbf{(Correct)} $\frac{2 \alpha}{k}$
    \item $\frac{\alpha}{k}$
    \item $\frac{\alpha}{2 k}$
    \item $\frac{\alpha}{4 k}$
\end{enumerate}

\noindent \textbf{Correct Answer:} (A)

\begin{tcolorbox}[solbox, title=Detailed Solution -- Question 18]
Given, $\quad$ input is $r(t)=\alpha t u(t)$

$R(s)=\frac{\alpha}{s^{2}}$

From the figure,

$G(s) H(s)=\frac{K}{s(s+2)}$

Now steady state error for Ramp input is

\[
e_{s s} =\frac{\alpha}{K_{v}}, \text { where } \alpha \text { is the magnitude of Ramp input }
\]

\[
\begin{aligned}K_{v} & =\lim _{s \rightarrow 0}[s G(s) H(s)] \\
K_{v} & =\lim _{s \rightarrow 0}\left[\frac{s \times K}{s(s+2)}\right]=\frac{K}{2}\\
\therefore \quad e_{s s} & =\frac{\alpha \times 2}{K} \\
e_{s s} & =\frac{2 \alpha}{K}
\end{aligned}
\]
\end{tcolorbox}

\vspace{1.5em}
\hrule
\vspace{1.5em}

\section*{Question 19: Basics, Block Diagrams \& SFGs (GATE EC 2023)}
\addcontentsline{toc}{section}{Question 19: Basics, Block Diagrams \& SFGs (GATE EC 2023)}
\noindent \textbf{\color{primary}MCQ} \hfill \textbf{\color{secondary}1 Mark}


\noindent The open loop transfer function of a unity negative feedback system is $G(s)=\frac{k}{s\left(1+s T_{1}\right)\left(1+s T_{2}\right)}$, where $k, T_{1}$ and $T_{2}$ are positive constants. The phase crossover frequency, in rad/s, is


\begin{enumerate}[label=(\Alph*)]
    \item \textbf{(Correct)} $\frac{1}{\sqrt{T_{1} T_{2}}}$
    \item $\frac{1}{T_{1} T_{2}}$
    \item $\frac{1}{T_{1} \sqrt{T_{2}}}$
    \item $\frac{1}{T_{2} \sqrt{T_{1}}}$
\end{enumerate}

\noindent \textbf{Correct Answer:} (A)

\begin{tcolorbox}[solbox, title=Detailed Solution -- Question 19]
We know phase crossover frequency is that frequency at which phase of the open loop transfer function is $-180^{\circ}$.

\[
\begin{aligned}
\therefore \quad G(s) & =\frac{K}{s\left(1+s T_{1}\right)\left(1+s T_{2}\right)} \\
G(j \omega) & =\frac{K}{(j \omega)\left(1+j \omega T_{1}\right)\left(1+j \omega T_{2}\right)} \\
\text{Phase of }G(j \omega) & =\phi =-90-\tan ^{-1}\left(\omega T_{1}\right) - \tan ^{-1}\left(\omega T_{2}\right) \\
\therefore \quad \text { At } \omega &=\omega_{p c}, \phi =-180  \\
-180 & =-90-\tan ^{-1}\left(\omega_{p c} T_{1}\right)-\tan ^{-1}\left(\omega_{p c} T_{2}\right) \\
9 \quad 90 & =\tan ^{-1}\left(\omega_{p c} T_{1}\right)+\tan ^{-1}\left(\omega_{p c} T_{2}\right) \\
\tan ^{-1}\left(\frac{\omega_{p c} T_{1}+\omega_{p c} T_{2}}{1-\omega_{p c}^{2} T_{1} T_{2}}\right) & =90 \\
1-\omega_{p c}^{2} T_{1} T_{2} & =0 \\
\omega_{p c} & =\frac{1}{\sqrt{T_{1} T_{2}}}
\end{aligned}
\]
\end{tcolorbox}

\vspace{1.5em}
\hrule
\vspace{1.5em}

\section*{Question 20: Time Response Analysis (GATE EC 2022)}
\addcontentsline{toc}{section}{Question 20: Time Response Analysis (GATE EC 2022)}
\noindent \textbf{\color{primary}NAT} \hfill \textbf{\color{secondary}2 Marks}


\noindent The block diagram of a closed-loop control system is shown in the figure. R(s),Y(s) and D(s) are the Laplace transforms of the time-domain signals $r(t),y(t),$ and $d(t)$, respectively. Let the error signal be defined as $e(t)=r(t)-y(t)$. Assuming the
reference input $r(t)=0$ for all $t$, the steady-state error $e(\infty )$, due to a unit step
disturbance $d(t)$, is \underline{\hspace{1.5cm}} (rounded off to two decimal places).

\begin{center}\includegraphics[max width=0.85\linewidth]{images/img\_115.jpg}\end{center}

\begin{center}\includegraphics[max width=0.85\linewidth]{images/img\_115.jpg}\end{center}


\begin{tcolorbox}[solbox, title=Detailed Solution -- Question 20]
\[
Y(s)=\frac{R(s)\cdot \frac{10}{s(s+10)}}{1+\frac{10}{s(s+10)}} +\frac{D(s)\cdot \frac{10}{s(s+10)}}{1+\frac{10}{s(s+10)}}
\]

When,
$r(t)=0\; and \; d(t)=U(t)\underleftrightarrow{L.T.}\frac{1}{s}$

\[
e(\infty )=-\lim_{s \to 0}\frac{s\cdot \frac{1}{s} \times 1}{s(s+10)+10} =-1/10=-0.1\;\;\;(\because e(t)=r(t)-y(t))
\]
\end{tcolorbox}

\vspace{1.5em}
\hrule
\vspace{1.5em}

\section*{Question 21: Time Response Analysis (GATE EC 2022)}
\addcontentsline{toc}{section}{Question 21: Time Response Analysis (GATE EC 2022)}
\noindent \textbf{\color{primary}MSQ} \hfill \textbf{\color{secondary}2 Marks}


\noindent Two linear time-invariant systems with transfer functions
 $G_1(s)=\frac{10}{s^2+s+1}, G_2(s)=\frac{10}{s^2+s\sqrt{10}+10}$

have unit step responses $y_1(t)$ and $y_2(t)$, respectively. Which of the following
statements is/are true?


\begin{enumerate}[label=(\Alph*)]
    \item \textbf{(Correct)} $y_1(t)$ and $y_2(t)$ have the same percentage peak overshoot.
    \item $y_1(t)$ and $y_2(t)$ have the same steady-state value.
    \item $y_1(t)$ and $y_2(t)$ have the same damped frequency of oscillation.
    \item $y_1(t)$ and $y_2(t)$ have the same 2% settling time.
\end{enumerate}

\noindent \textbf{Correct Answer:} (A)

\begin{tcolorbox}[solbox, title=Detailed Solution -- Question 21]
\[
\begin{aligned}
G_1(s)&=\frac{10}{s^2+s+1}\\
\omega _{n1}^2&=1\\
2\xi _1 \times 1&=1\\
\xi _1&=1/2\\
G_2(s)&=\frac{10}{s^2+s\sqrt{10}+10}\\
\omega _{n2}^2&=10\\
2\xi _1 \times \sqrt{10}&=\sqrt{10}\\
\xi _2&=1/2\\
\end{aligned}
\]

$\because M_P$  depends on
$\xi$ only

$y_1(t) \; and \; y_2(t)$ have same percentage peak
overshoot.

$\omega _{d_1}=\omega _{n_1}\sqrt{1-\xi ^2}$

\[
\Rightarrow \omega _{d_1}\neq \omega _{d_2} \;\; \because \omega _{n_1}\neq \omega _{n_2}
\]

damped frequency of oscillation is not same.

\[
\begin{aligned}
T_s &=\frac{4}{\xi \omega _n}(2%\; setting \; time) \\
T_{s_1} &\neq  T_{s_2} \;\because \omega _{n_1}\neq \omega _{n_2}\\
C_1(s) &= \frac{10 \times 1/s}{s^2+s+1}\\
C_{ss} &=\lim_{s \to 0}s\cdot C_1(s) \\
&= 10\\
C_2(s) &= \frac{10 \times 1/s}{s^2+\sqrt{10}s+10}\\
C_{ss} &=\lim_{s \to 0}s\cdot C_2(s) \\
&= \frac{10}{10}1=1
\end{aligned}
\]

Steady - state value is not same.
\end{tcolorbox}

\vspace{1.5em}
\hrule
\vspace{1.5em}

\section*{Question 22: Stability Analysis \& Routh-Hurwitz (GATE EC 2022)}
\addcontentsline{toc}{section}{Question 22: Stability Analysis \& Routh-Hurwitz (GATE EC 2022)}
\noindent \textbf{\color{primary}MCQ} \hfill \textbf{\color{secondary}2 Marks}


\noindent Consider an even polynomial $p(s)$ given by
 $p(s)=s^4+5s^2+4+K$,

where $K$ is an unknown real parameter. The complete range of$K$ for which $p(s)$ has
all its roots on the imaginary axis is \underline{\hspace{1.5cm}}.


\begin{enumerate}[label=(\Alph*)]
    \item \textbf{(Correct)} $-4\leq K\leq \frac{9}{4}$
    \item $-3\leq K\leq \frac{9}{2}$
    \item $-6\leq K\leq \frac{5}{4}$
    \item $-5\leq K\leq 0$
\end{enumerate}

\noindent \textbf{Correct Answer:} (A)

\begin{tcolorbox}[solbox, title=Detailed Solution -- Question 22]
$p(s)=s^4+5s^2+4+k$

\[
\begin{matrix}
s^4 &  1& 5 &(4+k) \\
s^3& 0 & 0 & \\
s^2&  &  & \\
s&  &  & \\
s^0&  &  &
\end{matrix}
\]

$s^4+5s^2+(4+k)=0$

$4s^3+10s=0$

$(4s^2+10)=0$

$s=\pm j\sqrt{5/2}$

\[
\begin{matrix}
s^4 &  1& 5 &(4+k) \\
s^3& 4 & 10 & \\
s^2& 5/2 & (4+k) & \\
s^1&10-\frac{4(4+k)}{5/2}  &  & \\
s^0& (4+k) &  &
\end{matrix}
\]

\[
\begin{aligned}
(4+k) & > 0\\
k & > -4\\
10 \times \frac{5}{2}& > 4(4+K)\\
(4+K)& < \frac{25}{4}\\
K& < 9/4
\end{aligned}
\]

$\Rightarrow$  All roots be on imaginary axis
$-4\leq K\leqslant 9/4$
\end{tcolorbox}

\vspace{1.5em}
\hrule
\vspace{1.5em}

\section*{Question 23: Root Locus Technique (GATE EC 2022)}
\addcontentsline{toc}{section}{Question 23: Root Locus Technique (GATE EC 2022)}
\noindent \textbf{\color{primary}MCQ} \hfill \textbf{\color{secondary}1 Mark}


\noindent The root-locus plot of a closed-loop system with unity negative feedback and
transfer function $KG(s)$ in the forward path is shown in the figure. Note that $K$ is
varied from 0 to $\infty$.
 Select the transfer function $G(s)$ that results in the root-locus plot of the closed-loop
system as shown in the figure.

\begin{center}\includegraphics[max width=0.85\linewidth]{images/img\_121.jpg}\end{center}

\begin{center}\includegraphics[max width=0.85\linewidth]{images/img\_121.jpg}\end{center}


\begin{enumerate}[label=(\Alph*)]
    \item \textbf{(Correct)} $G(s)=\frac{1}{(s+1)^5}$
    \item $G(s)=\frac{1}{(s^5+1)}$
    \item $G(s)=\frac{s-1}{(s+1)^6}$
    \item $G(s)=\frac{s+1}{(s^6+1)}$
\end{enumerate}

\noindent \textbf{Correct Answer:} (A)

\begin{tcolorbox}[solbox, title=Detailed Solution -- Question 23]
Here 5 Root Locus branches are diverging from
same point, this can possible only when if we
have 5 poles in the system at the same point
because Root Locus branch departs from open
loop pole and

Number of Root Locus branches = Number of
open loop poles or Number of zero (Whichever is
greater).

Here, there are 5 multiple Real poles, and
matching with option (A).
\end{tcolorbox}

\vspace{1.5em}
\hrule
\vspace{1.5em}

\section*{Question 24: Frequency Response Analysis (GATE EC 2022)}
\addcontentsline{toc}{section}{Question 24: Frequency Response Analysis (GATE EC 2022)}
\noindent \textbf{\color{primary}MCQ} \hfill \textbf{\color{secondary}1 Mark}


\noindent Consider a closed-loop control system with unity negative feedback and $KG(s)$ in
the forward path, where the gain $K=2$. The complete Nyquist plot of the transfer
function $G(s)$ is shown in the figure. Note that the Nyquist contour has been chosen
to have the clockwise sense. Assume $G(s)$ has no poles on the closed right-half of
the complex plane. The number of poles of the closed-loop transfer function in the
closed right-half of the complex plane is \underline{\hspace{1.5cm}}

\begin{center}\includegraphics[max width=0.85\linewidth]{images/img\_2.jpg}\end{center}

\begin{center}\includegraphics[max width=0.85\linewidth]{images/img\_2.jpg}\end{center}


\begin{enumerate}[label=(\Alph*)]
    \item 0
    \item 1
    \item \textbf{(Correct)} 2
    \item 3
\end{enumerate}

\noindent \textbf{Correct Answer:} (C)

\begin{tcolorbox}[solbox, title=Detailed Solution -- Question 24]
\begin{center}\includegraphics[max width=0.85\linewidth]{images/img\_74.jpg}\end{center}

\begin{center}\includegraphics[max width=0.85\linewidth]{images/img\_74.jpg}\end{center}

For K = 2, point A will be -0.8x2 = -1.6

Hence N = -2, P = 0

(By default Nyquist contoure is considered in
clockwise direction)

P - N = 2

Number of closed loop pole in right side of the
complex plane.
\end{tcolorbox}

\vspace{1.5em}
\hrule
\vspace{1.5em}

\section*{Question 25: Compensators \& Controllers (GATE EC 2021)}
\addcontentsline{toc}{section}{Question 25: Compensators \& Controllers (GATE EC 2021)}
\noindent \textbf{\color{primary}NAT} \hfill \textbf{\color{secondary}2 Marks}


\noindent A unity feedback system that uses proportional-integral ($\text{PI}$
) control is shown in the figure.

\begin{center}\includegraphics[max width=0.85\linewidth]{images/img\_5.jpg}\end{center}

\begin{center}\includegraphics[max width=0.85\linewidth]{images/img\_5.jpg}\end{center}

The stability of the overall system is controlled by tuning the $\text{PI}$ control parameters $K_{P}$
and $K_{I}$. The maximum value of $K_{I}$
that can be chosen so as to keep the overall system stable or, in the worst case, marginally stable (rounded off to three decimal places) is \underline{\hspace{1.5cm}}


\begin{tcolorbox}[solbox, title=Detailed Solution -- Question 25]
\[
\begin{aligned} G H &=\left(\frac{s K_{p}+K_{I}}{s}\right)\left(\frac{2}{s^{3}+4 s^{2}+5 s+2}\right) \\ q(s) &=s^{4}+4 s^{3}+5 s^{2}+s\left(2+2 k_{p}\right)+2 k_{I} \\ \text{Necessary}:\qquad K_{p} & > -1 ; K_{I} > 0\\ &\begin{array}{l|ccc} s^{4} & 1 & 5 & 2 K_{I} \\ s^{3} & 4 & 2+2 K_{p} & \\ s^{2} & \frac{18-2 K_{p}}{4} & 2 K_{I} & \\ s^{1} & \left(9-K_{p}\right)\left(1+K_{p}\right)-8 K_{I} & & \\ s^{0} & 2 K_{I} & & \end{array}\\ \end{aligned}
\]

 Sufficient:

\[
\begin{aligned} \frac{18-2 K_{p}}{4}& > 0\\ \Rightarrow \quad K_{p}& < 9\\ \therefore \quad-1& < K_{p} < 9\\ \left(18-2 K_{p}\right)\left(2+2 K_{p}\right)-32 K_{I}& > 0\\ 32 K_{I}& < 36+32 K_{p}-4 K_{p}^{2}\\ \therefore \qquad \qquad 0& < K_{I} < \frac{36+32 K_{p}-4 K_{p}^{2}}{32} \end{aligned}
\]

\[
\begin{aligned} \text{If }K_{p}&=-1 \Rightarrow k_{I}=0\\ \text{If }K_{p}&=9 \Rightarrow k_{I}=0\\ \end{aligned}
\]
,

\[
\begin{aligned} \therefore \qquad \qquad \frac{d K_{I}}{d K_{p}}&=0\\ \Rightarrow \qquad \qquad 32-8 K_{p}&=0=0 \Rightarrow K_{p}=4 \end{aligned}
\]

$\therefore$ For $K_{p}=4, K_{I}$ is maximum, which is
$K_{I}=\frac{36+32 \times 4-64}{32}=3.125$
 For $K_{p}=4, K_{I} < 3.125$ for stability 
$\therefore K_{I \max }=3.125$

\begin{center}\includegraphics[max width=0.85\linewidth]{images/img\_64.jpg}\end{center}

\begin{center}\includegraphics[max width=0.85\linewidth]{images/img\_64.jpg}\end{center}
\end{tcolorbox}

\vspace{1.5em}
\hrule
\vspace{1.5em}

\section*{Question 26: State Space Analysis (GATE EC 2021)}
\addcontentsline{toc}{section}{Question 26: State Space Analysis (GATE EC 2021)}
\noindent \textbf{\color{primary}MCQ} \hfill \textbf{\color{secondary}2 Marks}


\noindent The electrical system shown in the figure converts input source current $i_{s}(t)$ to output voltage $v_{o}(t)$

\begin{center}\includegraphics[max width=0.85\linewidth]{images/img\_139.jpg}\end{center}

\begin{center}\includegraphics[max width=0.85\linewidth]{images/img\_139.jpg}\end{center}

Current i $i_{L}(t)$ in the inductor and voltage $v_{c}(t)$ across the capacitor are taken as the state variables, both assumed to be initially equal to zero, i.e., $i_{L}(0) =0$ and $v_{c}(0)=0$. The system is


\begin{enumerate}[label=(\Alph*)]
    \item completely state controllable as well as completely observable
    \item completely state controllable but not observable
    \item completely observable but not state controllable
    \item \textbf{(Correct)} neither state controllable nor observable
\end{enumerate}

\noindent \textbf{Correct Answer:} (D)

\begin{tcolorbox}[solbox, title=Detailed Solution -- Question 26]
\[
\begin{aligned} i_{s} &=v_{i}+v_{c} \\ \therefore \qquad v_{i} &=-v_{c}+i_{s} \\ v_{L} &=L i_{L} \\ v_{L} &=\left(i_{s}-i_{L}\right) \\ \therefore \qquad L i_{L} &=i_{s}-i_{L} \\ \therefore \qquad i_{L} &=-i_{L}+i_{s} \\ v_{o} &=v_{c} \\ X &=\left[\begin{array}{cc} -1 & 0 \\ 0 & -1 \end{array}\right] X+\left[\begin{array}{ll} 1 \\ 2 \end{array}\right] U \\ Y &=[0 \quad 1] X+[0] \cup \\ A &=-I \\ Q_{C} &=\left[\begin{array}{ll} B & A B \end{array}\right]=\left[\begin{array}{ll} 1 & -1 \\ 1 & -1 \end{array}\right] \Rightarrow\left|Q_{c}\right|=0 \\ Q_{o} &=\left[\begin{array}{ll} C^{T} & A^{\top} C^{\top} \end{array}\right]=\left[\begin{array}{cc} 0 & 0 \\ 1 & -1 \end{array}\right] \Rightarrow\left|Q_{o}\right|=0 \end{aligned}
\]
\end{tcolorbox}

\vspace{1.5em}
\hrule
\vspace{1.5em}

\section*{Question 27: Frequency Response Analysis (GATE EC 2021)}
\addcontentsline{toc}{section}{Question 27: Frequency Response Analysis (GATE EC 2021)}
\noindent \textbf{\color{primary}MCQ} \hfill \textbf{\color{secondary}1 Mark}


\noindent The complete Nyquist plot of the open-loop transfer function G(s)H(s) of a feedback control system in the figure.

\begin{center}\includegraphics[max width=0.85\linewidth]{images/img\_66.jpg}\end{center}

\begin{center}\includegraphics[max width=0.85\linewidth]{images/img\_66.jpg}\end{center}

If G(s)H(s) has one zero in the right-half of the s-plane, the number of poles that the closed-loop system will have in the right-half of the s-plane is


\begin{enumerate}[label=(\Alph*)]
    \item 0
    \item 1
    \item 4
    \item \textbf{(Correct)} 3
\end{enumerate}

\noindent \textbf{Correct Answer:} (D)

\begin{tcolorbox}[solbox, title=Detailed Solution -- Question 27]
The given Nyquist plot is not matched according to the data given in the question.
\end{tcolorbox}

\vspace{1.5em}
\hrule
\vspace{1.5em}

\section*{Question 28: Basics, Block Diagrams \& SFGs (GATE EC 2021)}
\addcontentsline{toc}{section}{Question 28: Basics, Block Diagrams \& SFGs (GATE EC 2021)}
\noindent \textbf{\color{primary}MCQ} \hfill \textbf{\color{secondary}1 Mark}


\noindent The block diagram of a feedback control system is shown in the figure

\begin{center}\includegraphics[max width=0.85\linewidth]{images/img\_60.jpg}\end{center}

\begin{center}\includegraphics[max width=0.85\linewidth]{images/img\_60.jpg}\end{center}

The transfer function $\dfrac{Y{\left ( s \right )}}{X {\left ( s \right )}}$ of the system is


\begin{enumerate}[label=(\Alph*)]
    \item $\frac{G_{1}+G_{2}+G_{1}G_{2}H}{1+G_{1}H}$
    \item $\frac{G_{1}+G_{2}}{1+G_{1}H+G_{2}H}$
    \item \textbf{(Correct)} $\frac{G_{1}+G_{2}}{1+G_{1}H}$
    \item $\frac{G_{1}+G_{2}+G_{1}G_{2}H}{1+G_{1}H+G_{2}H}$
\end{enumerate}

\noindent \textbf{Correct Answer:} (C)

\begin{tcolorbox}[solbox, title=Detailed Solution -- Question 28]
\begin{center}\includegraphics[max width=0.85\linewidth]{images/img\_87.jpg}\end{center}

\begin{center}\includegraphics[max width=0.85\linewidth]{images/img\_87.jpg}\end{center}

\[
\frac{Y}{R}=\frac{G_{1}(1-0)+G_{2}(1-0)}{1-\left[-G_{1} H\right]}=\frac{G_{1}+G_{2}}{1+G_{1} H}
\]
\end{tcolorbox}

\vspace{1.5em}
\hrule
\vspace{1.5em}

\section*{Question 29: Time Response Analysis (GATE EC 2020)}
\addcontentsline{toc}{section}{Question 29: Time Response Analysis (GATE EC 2020)}
\noindent \textbf{\color{primary}NAT} \hfill \textbf{\color{secondary}2 Marks}


\noindent Consider the following closed loop control system 

\begin{center}\includegraphics[max width=0.85\linewidth]{images/img\_96.jpg}\end{center}

\begin{center}\includegraphics[max width=0.85\linewidth]{images/img\_96.jpg}\end{center}

 where $G(s)=\frac{1}{s(s+1)}$ and $C(s)=K\frac{s+1}{s+3}$.
If the steady state error for a unit ramp input
is 0.1, then the value of K is \underline{\hspace{1.5cm}}.


\begin{tcolorbox}[solbox, title=Detailed Solution -- Question 29]
Open loop transfer function for the system $=C(s)\times G(s)=\frac{K(s+1)}{ (s+3)}\times \frac{1}{s(s+1)}$
Since the system is type-1, so far a given unit ramp input steady state
$e_{ss}=\frac{1}{K_{V}}$
where,  $K_{V}=\lim_{s\rightarrow 0}S\times \frac{K}{S(S+3)}=\frac{K}{3}$
so, $\; e_{ss}=\frac{1}{K/3}=\frac{3}{K}$
 Given that,  $\; e_{ss}=0.1$
So, $0.1=\frac{3}{K}\Rightarrow K=30$
\end{tcolorbox}

\vspace{1.5em}
\hrule
\vspace{1.5em}

\section*{Question 30: Time Response Analysis (GATE EC 2020)}
\addcontentsline{toc}{section}{Question 30: Time Response Analysis (GATE EC 2020)}
\noindent \textbf{\color{primary}NAT} \hfill \textbf{\color{secondary}2 Marks}


\noindent A system with transfer function $G(s)=\frac{1}{(s+1)(s+a)}, a > 0$ is subjected an input $5\cos 3t$.
The steady state output of the system is $\frac{1}{\sqrt{10}} \cos (3t-1.892)$. The value of $a$ is \underline{\hspace{1.5cm}}.


\begin{tcolorbox}[solbox, title=Detailed Solution -- Question 30]
Given that,

\[
\begin{aligned}G(j\omega )&=\frac{1}{(1+j\omega )(\alpha +j\omega )} \\ \left |G(j\omega ) \right |&=\frac{1}{\sqrt{(\omega ^{2}+1)(\omega ^{2}+\alpha ^{2})}} \\ \text{According to }& \text{question,}\\ \left | G(j\omega ) \right |_{\omega =3}&=\frac{1}{5\sqrt{10}}\\ \Rightarrow \frac{1}{\sqrt{(\omega ^{2}+1)(\omega ^{2}+\alpha ^{2})}}&=\frac{1}{5\sqrt{10}} \\ \Rightarrow \frac{1}{\sqrt{10(a^{2}+9)}}&=\frac{1}{5\sqrt{10}} \alpha ^{2}+9=25 \\ \alpha ^{2}&=16 \\ \alpha &=4\end{aligned}
\]
\end{tcolorbox}

\vspace{1.5em}
\hrule
\vspace{1.5em}

\section*{Question 31: Root Locus Technique (GATE EC 2020)}
\addcontentsline{toc}{section}{Question 31: Root Locus Technique (GATE EC 2020)}
\noindent \textbf{\color{primary}MCQ} \hfill \textbf{\color{secondary}2 Marks}


\noindent The characteristic equation of a system is

$s^3 + 3s^2 + (K + 2)s + 3K = 0$

In the root locus plot for the given system, as K varies from 0 to $\infty$, the break-away
or break-in point(s) lie within


\begin{enumerate}[label=(\Alph*)]
    \item \textbf{(Correct)} (-1,0)
    \item (-2,-1)
    \item (-3,-2)
    \item (-$\infty$
,-3)
\end{enumerate}

\noindent \textbf{Correct Answer:} (A)

\begin{tcolorbox}[solbox, title=Detailed Solution -- Question 31]
$Q(s)=1+G(s)H(s)=0$
$s^{3}+3s^{2}+2s+ks+3k=0$
$-k=\frac{s^{3}+3s^{2}+2s}{s+3}$

\[
-\frac{\mathrm{d} k}{\mathrm{d} s}=\frac{(s+3)(3s^{2}+6s+2)-(s^{3}+3s^{2}+2s)}{(s+3)^{2}}=0
\]

$3s^{3}+6s^{2}+2s+9s^{2}+18s+6-s^{3}-3s^{2}-2s=0$
$2s^{3}+12s^{2}+18s+6=0$
$s=-0.46,-3.87,-1.65$

\begin{center}\includegraphics[max width=0.85\linewidth]{images/img\_103.jpg}\end{center}

\begin{center}\includegraphics[max width=0.85\linewidth]{images/img\_103.jpg}\end{center}

$\therefore$
Break-away point lies between (0, -1), i.e. (-1, 0)
\end{tcolorbox}

\vspace{1.5em}
\hrule
\vspace{1.5em}

\section*{Question 32: State Space Analysis (GATE EC 2020)}
\addcontentsline{toc}{section}{Question 32: State Space Analysis (GATE EC 2020)}
\noindent \textbf{\color{primary}MCQ} \hfill \textbf{\color{secondary}2 Marks}


\noindent For the given circuit, which one of the following is the correct state equation? 

\begin{center}\includegraphics[max width=0.85\linewidth]{images/img\_104.jpg}\end{center}

\begin{center}\includegraphics[max width=0.85\linewidth]{images/img\_104.jpg}\end{center}


\begin{enumerate}[label=(\Alph*)]
    \item \textbf{(Correct)} A
    \item B
    \item C
    \item D
\end{enumerate}

\noindent \textbf{Correct Answer:} (A)

\begin{tcolorbox}[solbox, title=Detailed Solution -- Question 32]
From source transformation,

\begin{center}\includegraphics[max width=0.85\linewidth]{images/img\_134.jpg}\end{center}

\begin{center}\includegraphics[max width=0.85\linewidth]{images/img\_134.jpg}\end{center}

KVL in loop 1,
$2i_{1}=2i+0.5\frac{\mathrm{d} i}{\mathrm{d} t}+V$
$\frac{\mathrm{d} i}{\mathrm{d} t}=-2V-4i+4i_{1}$
KCL at node (a),
$i=0.25\frac{\mathrm{d} V}{\mathrm{d} t}+\frac{V-i_{2}}{1}$
$\frac{\mathrm{d} v}{\mathrm{d} t}=-4V+4i+4i_{2}$

\[
\begin{bmatrix} v\\ i \end{bmatrix}=\begin{bmatrix} -4 & 4\\ -2 & -4 \end{bmatrix}\begin{bmatrix} v\\ i \end{bmatrix}+\begin{bmatrix} 0 &4 \\ 4 &0 \end{bmatrix}\begin{bmatrix} i_{1}\\ i_{2} \end{bmatrix}
\]
\end{tcolorbox}

\vspace{1.5em}
\hrule
\vspace{1.5em}

\section*{Question 33: Time Response Analysis (GATE EC 2020)}
\addcontentsline{toc}{section}{Question 33: Time Response Analysis (GATE EC 2020)}
\noindent \textbf{\color{primary}NAT} \hfill \textbf{\color{secondary}1 Mark}


\noindent The loop transfer function of a negative feedback system is

$G(s)H(s)=\frac{K(s+11)}{s(s+2)(s+8)}$

The value of K, for which the system is marginally stable, is \underline{\hspace{1.5cm}}.


\begin{tcolorbox}[solbox, title=Detailed Solution -- Question 33]
Characteristic equation q(s) for the given open loop system will be
$q(s)=s^{3}+10s^{2}+16s+Ks+11K=0$
Using R-H criteria,

\begin{center}\includegraphics[max width=0.85\linewidth]{images/img\_9.jpg}\end{center}

\begin{center}\includegraphics[max width=0.85\linewidth]{images/img\_9.jpg}\end{center}

For System to be Marginally Stable

\[
\begin{aligned}\frac{10(16+K)-11K}{10}&=0 \\ 160+10K-11K&=0 \\ K&=160\end{aligned}
\]
\end{tcolorbox}

\vspace{1.5em}
\hrule
\vspace{1.5em}

\section*{Question 34: Frequency Response Analysis (GATE EC 2020)}
\addcontentsline{toc}{section}{Question 34: Frequency Response Analysis (GATE EC 2020)}
\noindent \textbf{\color{primary}MCQ} \hfill \textbf{\color{secondary}1 Mark}


\noindent The pole-zero map of a rational function G(s) is shown below. When the closed counter $\Gamma$ is mapped into the G(s)-plane, then the mapping encircles. 

\begin{center}\includegraphics[max width=0.85\linewidth]{images/img\_70.jpg}\end{center}

\begin{center}\includegraphics[max width=0.85\linewidth]{images/img\_70.jpg}\end{center}


\begin{enumerate}[label=(\Alph*)]
    \item the origin of the G(s)-plane once in the counter-clockwise direction.
    \item \textbf{(Correct)} the origin of the G(s)-plane once in the clockwise direction.
    \item the point -1+j0 of the G(s)-plane once in the counter-clockwise direction.
    \item the point -1+j0 of the G(s)-plane once in the clockwise direction.
\end{enumerate}

\noindent \textbf{Correct Answer:} (B)

\begin{tcolorbox}[solbox, title=Detailed Solution -- Question 34]
s-plane contour is encircling 2-poles and 3-zeros in clockwise direction hence the corresponding G(s) plane contour encircles origin 2-times in anti-clockwise direction and 3-times in clockwise direction.

Therefore,  Effectively once in clockwise direction.
\end{tcolorbox}

\vspace{1.5em}
\hrule
\vspace{1.5em}

\section*{Question 35: Stability Analysis \& Routh-Hurwitz (GATE EC 2019)}
\addcontentsline{toc}{section}{Question 35: Stability Analysis \& Routh-Hurwitz (GATE EC 2019)}
\noindent \textbf{\color{primary}NAT} \hfill \textbf{\color{secondary}2 Marks}


\noindent Consider a unity feedback system, as in the figure shown, with an integral compensator $\frac{K}{s}$ and open-loop transfer function 

$G(s)=\frac{1}{s^2+3s+2}$

  where $k > 0$. The positive value of K for which there are exactly two poles of the unity feedback system on the $j\omega$ axis is equal to \underline{\hspace{1.5cm}} (rounded off to two decimal places).  

\begin{center}\includegraphics[max width=0.85\linewidth]{images/img\_100.jpg}\end{center}

\begin{center}\includegraphics[max width=0.85\linewidth]{images/img\_100.jpg}\end{center}


\begin{tcolorbox}[solbox, title=Detailed Solution -- Question 35]
$\frac{Y(s)}{X(s)}=\frac{K}{s^{3}+3 s^{2}+2 s+K}$
 Two poles of this system lie the system is moro:
 System is marginally stable. 
$\qquad k_{\text{mar}=3 \times 2=6}$
\end{tcolorbox}

\vspace{1.5em}
\hrule
\vspace{1.5em}

\section*{Question 36: State Space Analysis (GATE EC 2019)}
\addcontentsline{toc}{section}{Question 36: State Space Analysis (GATE EC 2019)}
\noindent \textbf{\color{primary}MCQ} \hfill \textbf{\color{secondary}2 Marks}


\noindent Let the state-space representation of an LTI system be $\dot{x}(t)=Ax(t)+Bu(t)$,  $\dot{y}(t)=Cx(t)+du(t)$  where A, B, C are matrices, d is a scalar, u(t) is the input to the system, and y(t) is its output. Let $B=[0\;0\;1]^T$ and d=0. Which one of the following options for A and C will ensure that the transfer function of this LTI system is 
 $H(s)=\frac{1}{s^3+3s^2+2s+1}$? 

\begin{center}\includegraphics[max width=0.85\linewidth]{images/img\_50.jpg}\end{center}

\begin{center}\includegraphics[max width=0.85\linewidth]{images/img\_50.jpg}\end{center}


\begin{enumerate}[label=(\Alph*)]
    \item \textbf{(Correct)} A
    \item B
    \item C
    \item D
\end{enumerate}

\noindent \textbf{Correct Answer:} (A)

\begin{tcolorbox}[solbox, title=Detailed Solution -- Question 36]
\[
\begin{aligned} X(t) &=A x(t)+B u(t) \\ y(t) &=C x(t) \\ B &=\left[\begin{array}{l} 0 \\ 0 \\ 1 \end{array}\right] \\ \frac{\gamma(s)}{U(s)} &=\\ \frac{\gamma(s)}{X_{1}(s)} \times \frac{X_{1}(s)}{U(s)}=& 1 \times \frac{1}{s^{3}+3 s^{2}+2 s+1} \\ x_{1}(s)\left(s^{3}+3 s^{2}+2 s+1\right]&=U(s) \\ x_{2} &=\dot{x}_{1}(t) \\ x_{2}(s) &=s x_{1}(s) \\ x_{3} &=\dot{x}_{2}(t) \\ \Rightarrow \quad x_{3}(s) &=s x_{2}(s)=s^{2} X_{1}(s) \\ \text { So, } \quad s x_{3}(s) &=-x_{1}(s)-2 x_{2}(s)-3 x_{3}(s)+u(s) \\ \dot{x}_{3}(t) &=-x_{1}(t)-2 r_{2}(t)-3 x_{3}(t)+u(t) \\ y(t) &=x_{1}(t) \\ \dot{x}(t) &=\left[\begin{array}{ccc}0 & 1 & 0 \\ 0 & 0 & 1 \\ -1 & -2 & -3\end{array}\right] x(t)+\left[\begin{array}{l}0 \\ 0 \\ 1\end{array}\right] u(t) \\ y(t) &=[1 \quad 0 \quad 0] x(t) \end{aligned}
\]
\end{tcolorbox}

\vspace{1.5em}
\hrule
\vspace{1.5em}

\section*{Question 37: Basics, Block Diagrams \& SFGs (GATE EC 2019)}
\addcontentsline{toc}{section}{Question 37: Basics, Block Diagrams \& SFGs (GATE EC 2019)}
\noindent \textbf{\color{primary}MCQ} \hfill \textbf{\color{secondary}2 Marks}


\noindent The block diagram of a system is illustrated in the figure shown, where X(s) is the input and Y(s) is the output. The transfer function $H(s)=\frac{Y(s)}{X(s)}$ is 

\begin{center}\includegraphics[max width=0.85\linewidth]{images/img\_140.jpg}\end{center}

\begin{center}\includegraphics[max width=0.85\linewidth]{images/img\_140.jpg}\end{center}


\begin{enumerate}[label=(\Alph*)]
    \item $H(s)=\frac{s^2+1}{s^3+s^2+s+1}$
    \item \textbf{(Correct)} $H(s)=\frac{s^2+1}{s^3+2s^2+s+1}$
    \item $H(s)=\frac{s+1}{s^2+s+1}$
    \item $H(s)=\frac{s^2+1}{2s^2+s+1}$
\end{enumerate}

\noindent \textbf{Correct Answer:} (B)

\begin{tcolorbox}[solbox, title=Detailed Solution -- Question 37]
Using block diagram reduction, we get,

\begin{center}\includegraphics[max width=0.85\linewidth]{images/img\_110.jpg}\end{center}

\begin{center}\includegraphics[max width=0.85\linewidth]{images/img\_110.jpg}\end{center}

$\frac{Y(s)}{X(s)}=H(s)=\frac{s^{2}+1}{s^{3}+2 s^{2}+s+1}$
\end{tcolorbox}

\vspace{1.5em}
\hrule
\vspace{1.5em}

\section*{Question 38: Time Response Analysis (GATE EC 2019)}
\addcontentsline{toc}{section}{Question 38: Time Response Analysis (GATE EC 2019)}
\noindent \textbf{\color{primary}MCQ} \hfill \textbf{\color{secondary}2 Marks}


\noindent Consider a causal second-order system with the transfer function 

$G(s)=\frac{1}{1+2s+s^2}$

 with a unit-step $R(s)=\frac{1}{s}$ as an input. Let C(s) be the corresponding output. The time taken by the system output c(t) to reach 94% of its steady-state value $\lim_{t \to \infty }c(t)$, rounded off to two decimal places, is


\begin{enumerate}[label=(\Alph*)]
    \item 5.25
    \item \textbf{(Correct)} 4.5
    \item 3.89
    \item 2.81
\end{enumerate}

\noindent \textbf{Correct Answer:} (B)

\begin{tcolorbox}[solbox, title=Detailed Solution -- Question 38]
\[
\begin{aligned} G(s)&=\frac{1}{s^{2}+2 s+1}=\frac{1}{(s+1)^{2}} \\ R(s)&=\frac{1}{s} \\ \alpha (s)&=G(s) R(s)=\frac{1}{s(s+1)^{2}} \end{aligned}
\]

 Using partial fraction expansion, we get,

\[
\begin{aligned} C(s) & =\frac{A}{s}+\frac{B}{(s+1)}+\frac{C}{(s+1)^{2}} \\ A\left(s^{2}+2 s+1\right)&+B\left(s^{2}+s\right)+C s=1 \\ A&=1 \\ A+B & =0 \Rightarrow B=-1 \\ 2 A+B+C & =0 \Rightarrow C=-1 \\ \therefore \quad C(s) & =\frac{1}{s}-\frac{1}{(s+1)}-\frac{1}{(s+1)^{2}} \\ \text{and}\quad c(t) & =\left(1-e^{-t}-t e^{-t}\right) u(t) \\ \lim _{t \rightarrow \infty} c(t) & =1 \end{aligned}
\]

 In order to reach 94 % of its steady-state value,
$\left(1-e^{-t}-t e^{-t}\right)=0.94$
 By trial and error, we get,
$t \approx 4.50 \mathrm{sec}$
\end{tcolorbox}

\vspace{1.5em}
\hrule
\vspace{1.5em}

\section*{Question 39: Frequency Response Analysis (GATE EC 2019)}
\addcontentsline{toc}{section}{Question 39: Frequency Response Analysis (GATE EC 2019)}
\noindent \textbf{\color{primary}MCQ} \hfill \textbf{\color{secondary}1 Mark}


\noindent For an LTI system, the Bode plot for its gain is as illustrated in the figure shown. The number of system poles $N_p$ and the number of system zeros $N_z$ in the frequency range $1Hz\leq f\leq 10^7Hz$ is  

\begin{center}\includegraphics[max width=0.85\linewidth]{images/img\_135.jpg}\end{center}

\begin{center}\includegraphics[max width=0.85\linewidth]{images/img\_135.jpg}\end{center}


\begin{enumerate}[label=(\Alph*)]
    \item $N_p=5,N_z=2$
    \item \textbf{(Correct)} $N_p=6,N_z=3$
    \item $N_p=7,N_z=4$
    \item $N_p=4,N_z=2$
\end{enumerate}

\noindent \textbf{Correct Answer:} (B)

\begin{tcolorbox}[solbox, title=Detailed Solution -- Question 39]
\begin{center}\includegraphics[max width=0.85\linewidth]{images/img\_72.jpg}\end{center}

\begin{center}\includegraphics[max width=0.85\linewidth]{images/img\_72.jpg}\end{center}

Number of poles $(N_{P})$= 6

Number of zeros $(N_{Z})$ = 3
\end{tcolorbox}

\vspace{1.5em}
\hrule
\vspace{1.5em}

\section*{Question 40: Frequency Response Analysis (GATE EC 2018)}
\addcontentsline{toc}{section}{Question 40: Frequency Response Analysis (GATE EC 2018)}
\noindent \textbf{\color{primary}NAT} \hfill \textbf{\color{secondary}2 Marks}


\noindent The figure below shows the Bode magnitude and phase plots of a stable transfer function $G(s)=\frac{n_{o}}{s^{2}+d_{2}s^{2}+d_{1}+d_{0}}$  

\begin{center}\includegraphics[max width=0.85\linewidth]{images/img\_69.jpg}\end{center}

\begin{center}\includegraphics[max width=0.85\linewidth]{images/img\_69.jpg}\end{center}
 Consider the negative unity feedback configuration with gain k in the feedforward path. The closed loop is stable for $k < k_{o}$. The maximum value of $k_{o}$ is \underline{\hspace{1.5cm}}.


\begin{tcolorbox}[solbox, title=Detailed Solution -- Question 40]
For G(s) 
$M_{\mathrm{dB}}\left(\omega_{p c}\right)=20 \mathrm{dB}$
 When cascaded with k ,

\[
\begin{aligned} \mathrm{GM}_{\mathrm{dB}} &=-20 \mathrm{dB}-20 \log _{10}(k) > 0 \mathrm{dB} \\ 20+20 \log _{10}(k) & < 0 \\ 20 \log _{10}(k) & < -20 \\ k & < 10^{-1}=0.10\\ \text{So,}\quad k_{0}&=0.10 \end{aligned}
\]
\end{tcolorbox}

\vspace{1.5em}
\hrule
\vspace{1.5em}

\section*{Question 41: Frequency Response Analysis (GATE EC 2018)}
\addcontentsline{toc}{section}{Question 41: Frequency Response Analysis (GATE EC 2018)}
\noindent \textbf{\color{primary}NAT} \hfill \textbf{\color{secondary}2 Marks}


\noindent For a unity feedback control system with the forward path transfer function 

$G(s)=\frac{K}{s(s+2)}$ 

 The peak resonant magnitude $M_{r}$, of the closed-loop frequency response is 2. The corresponding value of the gain K (correct to two decimal places) is \underline{\hspace{1.5cm}}.


\begin{tcolorbox}[solbox, title=Detailed Solution -- Question 41]
Maximum resonant peak, 

\[
\begin{aligned} M_{r} &=\frac{1}{2 \xi \sqrt{1-\xi^{2}}}=2 \\ 2 \xi \sqrt{1-\xi^{2}} &=\frac{1}{2} \\ \xi^{2}\left(1-\xi^{2}\right) &=\frac{1}{16} \\ \xi^{4}-\xi^{2}+\frac{1}{16} &=0 \\ \xi^{2}&=\frac{1}{2} \pm \sqrt{\frac{1-\frac{1}{4}}{4}}=\frac{1}{2} \pm \frac{\sqrt{3}}{4} \\ \text { As } M_{r}=2 > 1, \xi& < \frac{1}{\sqrt{2}} \text { and } \xi^{2} < \frac{1}{2} \\ \text { So. } \quad\xi^{2}&=\frac{1}{2}-\frac{\sqrt{3}}{4}\\ Given, \quad G(s)&=\frac{K}{s(s+2)}=\frac{\omega_{n}^{2}}{s\left(s+2 \xi \omega_{n}\right)} \\ \text{So} \quad\omega_{n} &=\sqrt{K} \\ 2 \xi \sqrt{K} &=2\\ \sqrt{K} &=\frac{1}{\xi} \\ K &=\frac{1}{\xi^{2}}=\frac{1}{\left(\frac{1}{2}-\frac{\sqrt{3}}{4}\right)} \\ &=\frac{4}{2-\sqrt{3}}=14.928 \end{aligned}
\]
\end{tcolorbox}

\vspace{1.5em}
\hrule
\vspace{1.5em}

\section*{Question 42: State Space Analysis (GATE EC 2018)}
\addcontentsline{toc}{section}{Question 42: State Space Analysis (GATE EC 2018)}
\noindent \textbf{\color{primary}MCQ} \hfill \textbf{\color{secondary}2 Marks}


\noindent The state equation and the output equation of a control system are given below:

\[
\dot{x}=\begin{bmatrix} -4 & -1.5\\ 4 & 0 \end{bmatrix}x+\begin{bmatrix} 2\\ 0\end{bmatrix}u,
\]
 
 
\[
y=\begin{bmatrix} 1.5 & 0.625 \end{bmatrix}x.
\]
  

The transfer function representation of the system is


\begin{enumerate}[label=(\Alph*)]
    \item \textbf{(Correct)} $\frac{3s+5}{s^{2}+4s+6}$
    \item $\frac{3s-1.875}{s^{2}+4s+6}$
    \item $\frac{4s+1.5}{s^{2}+4s+6}$
    \item $\frac{4s+1.5}{s^{2}+4s+6}$
\end{enumerate}

\noindent \textbf{Correct Answer:} (A)

\begin{tcolorbox}[solbox, title=Detailed Solution -- Question 42]
Transfer function

\[
\begin{aligned} T(s) &=\frac{Y(s)}{U(s)}=\mathrm{C}(s I-A]^{-1} \mathrm{B} \\ A &=\left[\begin{array}{cc} -4 & -1.5 \\ 4 & 0 \end{array}\right] \\ B &=\left[\begin{array}{c} 2 \\ 0 \end{array}\right] \\ C&=[1.5\quad0.625]\\ {[s I-A]} & =\left[\begin{array}{cc}s+4 & 1.5 \\-4 & s\end{array}\right] \\ {[s I-A]^{-1}} & =\frac{1}{\left(s^{2}+4 s+6\right)}\left[\begin{array}{cc}s & -1.5 \\4 & s+4 \end{array}\right] \\ {[s I-A]^{-1} B} & =\frac{1}{s^{2}+4 s+6}\left[\begin{array}{c}2 s \\8 \end{array}\right] & \\ C[s I-A]^{-1} B & =\frac{1}{s^{2}+4 s+6}[1.5 \quad 0.625]\left[\begin{array}{c}2 s \\ 8\end{array}\right] \\ T(s) & =\frac{3 s+5}{s^{2}+4 s+6} \end{aligned}
\]
\end{tcolorbox}

\vspace{1.5em}
\hrule
\vspace{1.5em}

\section*{Question 43: Frequency Response Analysis (GATE EC 2018)}
\addcontentsline{toc}{section}{Question 43: Frequency Response Analysis (GATE EC 2018)}
\noindent \textbf{\color{primary}MCQ} \hfill \textbf{\color{secondary}1 Mark}


\noindent The Nyquist stability criterion and the Routh criterion both are powerful analysis tools for
determining the stability of feedback controllers. Identify which of the following statements
is FALSE:


\begin{enumerate}[label=(\Alph*)]
    \item Both the criteria provide information relative to the stable gain range of the system.
    \item The general shape of the Nyquist plot is readily obtained from the Bode magnitude plot
for all minimum-phase systems.
    \item The Routh criterion is not applicable in the condition of transport lag, which can be
readily handled by the Nyquist criterion.
    \item \textbf{(Correct)} The closed-loop frequency response for a unity feedback system cannot be obtained
from the Nyquist plot.
\end{enumerate}

\noindent \textbf{Correct Answer:} (D)


\vspace{1.5em}
\hrule
\vspace{1.5em}

\section*{Question 44: State Space Analysis (GATE EC 2017-SET-2)}
\addcontentsline{toc}{section}{Question 44: State Space Analysis (GATE EC 2017-SET-2)}
\noindent \textbf{\color{primary}MCQ} \hfill \textbf{\color{secondary}2 Marks}


\noindent A second - order LTI system is described by the following state equations, 
 $\frac{d}{dt}x_{1}(t)- x_{2}(t)=0$   
        $\frac{d}{dt}x_{2}(t)+2x_{1}(t)+3x_{2}(t)=r(t)$    
     Where $x_{1}(t)$ and  $x_{2}(t)$ are the two state variables and r(t) denotes the input. The output  $c(t) =x_{1}(t)$. The system is.


\begin{enumerate}[label=(\Alph*)]
    \item Undamped (oscillatory)
    \item Under damped
    \item Critically damped
    \item \textbf{(Correct)} Over damped
\end{enumerate}

\noindent \textbf{Correct Answer:} (D)

\begin{tcolorbox}[solbox, title=Detailed Solution -- Question 44]
Given state variable equations are as follows:

\[
\begin{aligned} &\frac{d}{d t} x_{1}(t)-x_{2}(t)=0 &\ldots(i)\\ &\frac{d}{d t} x_{2}(t)+2 x_{2}(t)+3 x_{2}(t)=r(t) &\ldots(ii) \end{aligned}
\]

 Also given that, input $=r(t)$ and output $=x_{1}(t)$
 By applying Laplace transform to equation (i), we
 get,
$s X_{1}(s)=X_{2}(s) \qquad \ldots(iii)$
 By applying Laplace transform to equation (ii).
 we get,

\[
\begin{aligned} s X_{2}(s)+2 X_{1}(s)+3 X_{2}(s)&=R(s) \\ x_{2}(s)&=\frac{R(s)-2 X_{1}(s)}{s+3} &\ldots(iv) \end{aligned}
\]

 By substituting equation (iv) in equation (iii), we get,

\[
\begin{aligned} s X_{\mathrm{i}}(s)&=\frac{R(s)-2 X_{1}(s)}{s+3}\\ \left(s^{2}+3 s+2\right) X_{1}(s)&=R(s) \end{aligned}
\]

So, the transfer function of the given system is,

\[
\begin{aligned}
\frac{X_{1}(s)}{R(s)} &=\frac{1}{s^{2}+3 s+2} \\
v r^{2} &=2 \\
2 \\\zeta_{n} &=3 \\
\xi &=\frac{3}{2 \omega_{n}}=\frac{3}{2 \sqrt{2}}=1.06
\end{aligned}
\]

As  $\xi > 1,$  the given system is overdamped.
\end{tcolorbox}

\vspace{1.5em}
\hrule
\vspace{1.5em}

\section*{Question 45: Stability Analysis \& Routh-Hurwitz (GATE EC 2017-SET-2)}
\addcontentsline{toc}{section}{Question 45: Stability Analysis \& Routh-Hurwitz (GATE EC 2017-SET-2)}
\noindent \textbf{\color{primary}NAT} \hfill \textbf{\color{secondary}2 Marks}


\noindent A unity feedback control system is characterized by the open-loop transfer function   
$G(s)=\frac{2(s+1)}{s^{3}+ks^{2}+2s+1}$ 
 The value of k for which the system oscillates at 2 rad/s is \underline{\hspace{1.5cm}}.


\begin{tcolorbox}[solbox, title=Detailed Solution -- Question 45]
The given open loop transfer function is,
$G(s)=\frac{2(s+1)}{s^{3}+K s^{2}+2 s+1}$
 The closed loop transfer function is,
$T(s)=\frac{G(s)}{1+G(s)}=\frac{2(s+1)}{s^{3}+K s^{2}+4 s+3}$
 When system oscillates, i.e., when system is marginally stable,
$(1)(3)=K(4)$
$\mathrm{So}, \quad K=0.75$
\end{tcolorbox}

\vspace{1.5em}
\hrule
\vspace{1.5em}

\section*{Question 46: Frequency Response Analysis (GATE EC 2017-SET-2)}
\addcontentsline{toc}{section}{Question 46: Frequency Response Analysis (GATE EC 2017-SET-2)}
\noindent \textbf{\color{primary}MCQ} \hfill \textbf{\color{secondary}2 Marks}


\noindent A unity feedback control system is characterized by the open-loop transfer function 
$G(s)=\frac{10K(s+2)}{s^{3}+3s^{2}+10}$ 
  The Nyquist path and the corresponding Nyquist plot of G(s) are shown in the figures below.

\begin{center}\includegraphics[max width=0.85\linewidth]{images/img\_82.jpg}\end{center}

\begin{center}\includegraphics[max width=0.85\linewidth]{images/img\_82.jpg}\end{center}
 
 If $0 < K < 1$, then the number of poles of the closed-loop transfer function that lie in the right half of the s-plane is


\begin{enumerate}[label=(\Alph*)]
    \item 0
    \item 1
    \item \textbf{(Correct)} 2
    \item 3
\end{enumerate}

\noindent \textbf{Correct Answer:} (C)

\begin{tcolorbox}[solbox, title=Detailed Solution -- Question 46]
Given that, the open loop transfer function of a  unity feedback system is,
$G(s) =\frac{10 k(s+2)}{s^{3}+3 s^{2}+10}$
 From the given Nyquist plot, for $0 < k < 1$ the encirclements about the point (-1+j 0) is,

\[
\begin{array}{l} N=0 \\ N=P-Z \end{array}
\]

 P= number of open loop poles in the right half of s-plane
 Z= number of closed loop poles in the right half of s-plane
 To determine the value of P :
 Applying R-H criteria to G(s) 
$G(s)=\frac{10 k(s+2)}{s^{3}+3 s^{2}+10}$

\[
\begin{array}{c|cc} s^{3} & 1 & 0 \\ s^{2} & 3 & 10 \\ s^{1} & -10 / 3 & 0 \\ s^{0} & 10 & 0 \end{array}
\]

Two sign changes are there in the first column of
the RH table. So, two open loop poles are there in
right half of s-plane.

So,
$P=2$

\[
\begin{array}{l}
N=P-Z \\
0=2-Z \\
Z=2
\end{array}
\]

Z=2 indicates, there are two closed loop poles
in right half of s-plane.
\end{tcolorbox}

\vspace{1.5em}
\hrule
\vspace{1.5em}

\section*{Question 47: Basics, Block Diagrams \& SFGs (GATE EC 2017-SET-2)}
\addcontentsline{toc}{section}{Question 47: Basics, Block Diagrams \& SFGs (GATE EC 2017-SET-2)}
\noindent \textbf{\color{primary}NAT} \hfill \textbf{\color{secondary}1 Mark}


\noindent For the system shown in the figure, Y(s)/X(s)= \underline{\hspace{1.5cm}}.

\begin{center}\includegraphics[max width=0.85\linewidth]{images/img\_19.jpg}\end{center}

\begin{center}\includegraphics[max width=0.85\linewidth]{images/img\_19.jpg}\end{center}


\begin{tcolorbox}[solbox, title=Detailed Solution -- Question 47]
\begin{center}\includegraphics[max width=0.85\linewidth]{images/img\_108.jpg}\end{center}

\begin{center}\includegraphics[max width=0.85\linewidth]{images/img\_108.jpg}\end{center}

\[
\begin{aligned} Y(s) &=[X(s)-Y(s)] G(s)+X(s) \\ [1+G(s)] Y(s) &=[G(s)+1] X(s) \\ \frac{Y(s)}{X(s)} &=1 \end{aligned}
\]
\end{tcolorbox}

\vspace{1.5em}
\hrule
\vspace{1.5em}

\section*{Question 48: Compensators \& Controllers (GATE EC 2017-SET-2)}
\addcontentsline{toc}{section}{Question 48: Compensators \& Controllers (GATE EC 2017-SET-2)}
\noindent \textbf{\color{primary}MCQ} \hfill \textbf{\color{secondary}1 Mark}


\noindent Which of the following statements is incorrect?


\begin{enumerate}[label=(\Alph*)]
    \item Lead compensator is used to reduce the settling time.
    \item Lag compensator is used to reduce the steady state error.
    \item Lead compensator may increase the order of a system.
    \item \textbf{(Correct)} Lag compensator always stabilizes an unstable system.
\end{enumerate}

\noindent \textbf{Correct Answer:} (D)

\begin{tcolorbox}[solbox, title=Detailed Solution -- Question 48]
In\_ case of high type systems Lag compensator
fails to give stability.
\end{tcolorbox}

\vspace{1.5em}
\hrule
\vspace{1.5em}

\section*{Question 49: State Space Analysis (GATE EC 2017-SET-2)}
\addcontentsline{toc}{section}{Question 49: State Space Analysis (GATE EC 2017-SET-2)}
\noindent \textbf{\color{primary}NAT} \hfill \textbf{\color{secondary}1 Mark}


\noindent Consider the state space realization 
 
\[
\begin{bmatrix} x_{1}(t)\\ x_{2}(t) \end{bmatrix}=\begin{bmatrix} 0 & 0\\ 0&-9 \end{bmatrix} \begin{bmatrix} x_{1}(t)\\ x_{2}(t) \end{bmatrix}+ \begin{bmatrix} 0\\ 45 \end{bmatrix}u(t)
\]
 , with the initial condition 
\[
\begin{bmatrix} x_{1}(0)\\ x_{2}(0) \end{bmatrix}=\begin{bmatrix} 0\\ 0 \end{bmatrix}
\]
  
 where u(t) denotes the unit step function. The value of  $\lim_{t\rightarrow \infty }|\sqrt{x_{1}^{2}(t)+x_{2}^{2}(t)}|$  is \underline{\hspace{1.5cm}}.


\begin{tcolorbox}[solbox, title=Detailed Solution -- Question 49]
\[
\left[\begin{array}{l} \dot{x}_{1}(t) \\ \dot{x}_{2}(t) \end{array}\right]=\left[\begin{array}{cc} 0 & 0 \\ 0 & -9 \end{array}\right]\left[\begin{array}{l} x_{1}(t) \\ x_{2}(t) \end{array}\right]+\left[\begin{array}{c} 0 \\ 45 \end{array}\right] u(t)
\]

 By applying Laplace transform on both sides. get

\[
\begin{aligned} s x_{1}(s)-x_{1}(0)&=0\\ X_{1}(s)&=\frac{x_{1}(0)}{s}=0 \quad \because x_{1}(0)=0 \\ \text{So,} x_{1}(t)&=0 \\ \text{and }\quad s X_{\lambda}(s)-x_{2}(0)&=-9 X_{1}(s)+\frac{45}{s} \\ X_{2}(s) &=\frac{45}{s(s+9)} \quad \because x_{2}(0)=0 \\ \text { Required value } &=\lim _{t \rightarrow \infty}|\sqrt{x_{1}^{2}(t)+x_{2}^{2}(t)}| \\ &=\left|\lim _{t \rightarrow \infty} x_{2}(t)\right| \quad \because x_{1}(t)=0 \\ \lim _{t \rightarrow \infty} x_{2}(t) &=\lim _{s \rightarrow 0} s X_{2}(s)=\frac{45}{9}=5 \end{aligned}
\]

 So,
 Required value $= |5| =5$
\end{tcolorbox}

\vspace{1.5em}
\hrule
\vspace{1.5em}

\section*{Question 50: Frequency Response Analysis (GATE EC 2017-SET-1)}
\addcontentsline{toc}{section}{Question 50: Frequency Response Analysis (GATE EC 2017-SET-1)}
\noindent \textbf{\color{primary}MCQ} \hfill \textbf{\color{secondary}2 Marks}


\noindent The Nyquist plot of the transfer function 

$G(s)=\frac{K}{(s^{2}+2s+2)(s+2)}$

 does not encircle the point (1+j0) for K=10 but does encircle the point (-1+j0) for K=100. Then the closed loop system (having unity gain feedback) is


\begin{enumerate}[label=(\Alph*)]
    \item stable for K = 10 and stable for K = 100
    \item \textbf{(Correct)} stable for K = 10 and unstable for K = 100
    \item unstable for K = 10 and stable for K = 100
    \item unstable for K = 10 and unstable for K = 100
\end{enumerate}

\noindent \textbf{Correct Answer:} (B)

\begin{tcolorbox}[solbox, title=Detailed Solution -- Question 50]
Given that,
$G(s)=\frac{K}{\left(s^{2}+2 s+2\right)(s+2)}$
 P= open loop poles in RHS of s-plane 
 Z= closed loop poles in RHS of s-plane 
 N= number of encirclements about the

\[
\begin{array}{l} \text { point }(-1+j 0) \\ \text { For } K=10 ; \\ N=P-Z=0 \Rightarrow Z=0: \text { stable } \\ \text { For } K=100 ; \\ N=P-Z=1 \Rightarrow Z \neq 0: \text { unstable } \end{array}
\]
\end{tcolorbox}

\vspace{1.5em}
\hrule
\vspace{1.5em}

\section*{Question 51: Stability Analysis \& Routh-Hurwitz (GATE EC 2017-SET-1)}
\addcontentsline{toc}{section}{Question 51: Stability Analysis \& Routh-Hurwitz (GATE EC 2017-SET-1)}
\noindent \textbf{\color{primary}MCQ} \hfill \textbf{\color{secondary}2 Marks}


\noindent Which one of the following options correctly describes the locations of the roots of the equation $s^{4}+s^{2}+1=0$ on the complex plane?


\begin{enumerate}[label=(\Alph*)]
    \item Four left half plane (LHP) roots
    \item One right half plane (RHP) root, one LHP root and two roots on the imaginary axis
    \item \textbf{(Correct)} Two RHP roots and two LHP roots
    \item All four roots are on the imaginary axis
\end{enumerate}

\noindent \textbf{Correct Answer:} (C)

\begin{tcolorbox}[solbox, title=Detailed Solution -- Question 51]
$q(s)=s^{4}+s^{2}+1=0$

\[
\begin{array}{r|rrr} s^{4} & 1 & 1 & 1\\ s^{3} & 4 & 2 & 0\\ s^{2} & 0.5 & 1 & 0 \\ s^{1} & -6 & 0 & 0 \\ s^{0} & 1 & 0 & 0 \end{array} \frac{d A(s)}{d s}=4 s^{3}+2 s^{1}
\]

 There are two sign Changes in the first column of the R-H table and the order of auxiliary equation is 4. So, four poles are symmetric about origin.
$\therefore$2 RHP roots and 2 LHP roots.

\begin{center}\includegraphics[max width=0.85\linewidth]{images/img\_88.jpg}\end{center}

\begin{center}\includegraphics[max width=0.85\linewidth]{images/img\_88.jpg}\end{center}
\end{tcolorbox}

\vspace{1.5em}
\hrule
\vspace{1.5em}

\section*{Question 52: Root Locus Technique (GATE EC 2017-SET-1)}
\addcontentsline{toc}{section}{Question 52: Root Locus Technique (GATE EC 2017-SET-1)}
\noindent \textbf{\color{primary}MCQ} \hfill \textbf{\color{secondary}2 Marks}


\noindent A linear time invariant (LTI) system with the transfer function 

$G(s)=\frac{K(s^{2}+2s+2)}{(s^{2}-3s+2)}$ 

is connected in unity feedback configuration as shown in the figure.

\begin{center}\includegraphics[max width=0.85\linewidth]{images/img\_25.jpg}\end{center}

\begin{center}\includegraphics[max width=0.85\linewidth]{images/img\_25.jpg}\end{center}

 For the closed loop system shown, the root locus for $0 < K < \infty$ intersects the imaginary axis for K = 1.5. The closed loop system is stable for


\begin{enumerate}[label=(\Alph*)]
    \item \textbf{(Correct)} $K>1.5$
    \item $1 < K < 1.5$
    \item $0 < K < 1$
    \item no positive value of K
\end{enumerate}

\noindent \textbf{Correct Answer:} (A)

\begin{tcolorbox}[solbox, title=Detailed Solution -- Question 52]
Given that,
$G(s)=\frac{K\left(s^{2}+2 s+2\right)}{s^{2}-3 s+2}$

The root locus plot of the given system is
follows:

\begin{center}\includegraphics[max width=0.85\linewidth]{images/img\_125.jpg}\end{center}

\begin{center}\includegraphics[max width=0.85\linewidth]{images/img\_125.jpg}\end{center}

$\therefore$System is stable for $K >  1.5$.
\end{tcolorbox}

\vspace{1.5em}
\hrule
\vspace{1.5em}

\section*{Question 53: Frequency Response Analysis (GATE EC 2017-SET-1)}
\addcontentsline{toc}{section}{Question 53: Frequency Response Analysis (GATE EC 2017-SET-1)}
\noindent \textbf{\color{primary}MCQ} \hfill \textbf{\color{secondary}1 Mark}


\noindent Consider a stable system with transfer function 
$G(s)=\frac{s^{p}+b_{1}S^{p-1}+....+b_p}{s^{q}+a_{1}S^{q-1}+....+a_{q}}$ 

 where $b_{1},...b_{p}$ and $a_{1},...a_{q}$ are real valued constants. The slope of the Bode log magnitude curve of G(s) converges to -60 dB/decade as $\omega \rightarrow \infty$  .  A possible pair of values for p and q is


\begin{enumerate}[label=(\Alph*)]
    \item \textbf{(Correct)} p=0 and q=3
    \item p=1 and q=7
    \item p=2 and q=3
    \item p=3 and q=5
\end{enumerate}

\noindent \textbf{Correct Answer:} (A)

\begin{tcolorbox}[solbox, title=Detailed Solution -- Question 53]
Final slope = -60 dB/decade,
which indicates that, q-p=3 

Among the given options, option (A) satisfies this
condition
\end{tcolorbox}

\vspace{1.5em}
\hrule
\vspace{1.5em}

\section*{Question 54: Time Response Analysis (GATE EC 2017-SET-1)}
\addcontentsline{toc}{section}{Question 54: Time Response Analysis (GATE EC 2017-SET-1)}
\noindent \textbf{\color{primary}NAT} \hfill \textbf{\color{secondary}1 Mark}


\noindent The open loop transfer function G(s)=$\frac{(s+1)}{s^{P}(s+2)(s+3)}$ 
Where p is an integer, is connected in unity feedback configuration as shown in figure.

\begin{center}\includegraphics[max width=0.85\linewidth]{images/img\_130.jpg}\end{center}

\begin{center}\includegraphics[max width=0.85\linewidth]{images/img\_130.jpg}\end{center}

  Given that the steady state error is zero for unit step input and is 6 for unit ramp input, the value of the parameter p is \underline{\hspace{1.5cm}}.


\begin{tcolorbox}[solbox, title=Detailed Solution -- Question 54]
Steady state error of type 1 system, for step input
is zero, for ramp input is $\frac{1}{K_{v}}$ and for parabolic input is infinity. 
 So, the given system must be of type- 1 and p=1
\end{tcolorbox}

\vspace{1.5em}
\hrule
\vspace{1.5em}

\section*{Question 55: Compensators \& Controllers (GATE EC 2017-SET-1)}
\addcontentsline{toc}{section}{Question 55: Compensators \& Controllers (GATE EC 2017-SET-1)}
\noindent \textbf{\color{primary}MCQ} \hfill \textbf{\color{secondary}1 Mark}


\noindent Which of the following can be pole-zero configuration of a phase-lag controller (lag compensator)? 

\begin{center}\includegraphics[max width=0.85\linewidth]{images/img\_12.jpg}\end{center}

\begin{center}\includegraphics[max width=0.85\linewidth]{images/img\_12.jpg}\end{center}


\begin{enumerate}[label=(\Alph*)]
    \item \textbf{(Correct)} A
    \item B
    \item C
    \item D
\end{enumerate}

\noindent \textbf{Correct Answer:} (A)

\begin{tcolorbox}[solbox, title=Detailed Solution -- Question 55]
Phase lag controller transfer function is
$G_{C}(s)=\frac{s+Z}{s+P} ;|Z|>|P|$

\begin{center}\includegraphics[max width=0.85\linewidth]{images/img\_133.jpg}\end{center}

\begin{center}\includegraphics[max width=0.85\linewidth]{images/img\_133.jpg}\end{center}
\end{tcolorbox}

\vspace{1.5em}
\hrule
\vspace{1.5em}

\section*{Question 56: Root Locus Technique (GATE EC 2016-SET-3)}
\addcontentsline{toc}{section}{Question 56: Root Locus Technique (GATE EC 2016-SET-3)}
\noindent \textbf{\color{primary}NAT} \hfill \textbf{\color{secondary}2 Marks}


\noindent The forward-path transfer function and the feedback-path transfer function of a single loop negative feedback control system are given as 

$G(s)=\frac{K(s+2)}{s^{2}+2s+2}$ and H(s)=1, 

respectively. If the variable parameter K is real positive, then the location of the breakaway point on
the root locus diagram of the system is \underline{\hspace{1.5cm}}


\begin{tcolorbox}[solbox, title=Detailed Solution -- Question 56]
\[
\begin{aligned} G(s) &=\frac{K(s+2)}{s^{2}+2 s+2} \\ H(s)&=1 \\ 1+G(s) H(s) &=0 \\ 1+\frac{K}{\left(s^{2}\right)+2 s+2}(s+2) &=0\\ K&=-\frac{\left(s^{2}+2 s+2\right)}{(s+2)} \end{aligned}
\]

\begin{center}\includegraphics[max width=0.85\linewidth]{images/img\_95.jpg}\end{center}

\begin{center}\includegraphics[max width=0.85\linewidth]{images/img\_95.jpg}\end{center}

\[
\frac{d K}{d s}=-\left(\frac{(2 s+2)(s+2)-\left(s^{2}+2 s+2\right)}{(s+2)^{2}}\right)
\]

 For break away points, $\frac{d K}{d s}=0$

\[
\begin{array}{l} (2 s+2)(s+2)-\left(s^{2}+2 s+2\right)=0 \\ 2 s^{2}+6 s+4-s^{2}-2 s-2=0 \\ s^{2}+4 s+2=0 \end{array}
\]

\[
\begin{aligned} s&=\frac{-4 \pm \sqrt{16-8}}{2} \\ &=\frac{-4 \pm 2 \sqrt{2}}{2}=-2 \pm \sqrt{2} \\ s&=-0.58 ; s=-3.41 \\ \text{But,}\quad s&=-3.41 \text{ lies on root locus }\\ \text{Hence}\quad s&=-3.41 \end{aligned}
\]
\end{tcolorbox}

\vspace{1.5em}
\hrule
\vspace{1.5em}

\section*{Question 57: State Space Analysis (GATE EC 2016-SET-3)}
\addcontentsline{toc}{section}{Question 57: State Space Analysis (GATE EC 2016-SET-3)}
\noindent \textbf{\color{primary}MCQ} \hfill \textbf{\color{secondary}2 Marks}


\noindent A second-order linear time-invariant system is described by the following state equations 

 $\frac{d}{dt}x_{1}(t)+2x_{1}(t)=3u(t)$ 
$\frac{d}{dt}x_{2}(t)+x_{2}(t)=u(t)$

 where $x_{1}(t) \; and \;  x_{2}(t)$ are the two state variables and u(t) denotes the input. If the output c(t) = $x_{1}(t)$, then the system is


\begin{enumerate}[label=(\Alph*)]
    \item \textbf{(Correct)} controllable but not observable
    \item observable but not controllable
    \item both controllable and observable
    \item neither controllable nor observable
\end{enumerate}

\noindent \textbf{Correct Answer:} (A)


\vspace{1.5em}
\hrule
\vspace{1.5em}

\section*{Question 58: Stability Analysis \& Routh-Hurwitz (GATE EC 2016-SET-3)}
\addcontentsline{toc}{section}{Question 58: Stability Analysis \& Routh-Hurwitz (GATE EC 2016-SET-3)}
\noindent \textbf{\color{primary}MCQ} \hfill \textbf{\color{secondary}2 Marks}


\noindent The first two rows in the Routh table for the characteristic equation of a certain closed-loop control system are given as 

\begin{center}\includegraphics[max width=0.85\linewidth]{images/img\_76.jpg}\end{center}

\begin{center}\includegraphics[max width=0.85\linewidth]{images/img\_76.jpg}\end{center}
 
The range of K for which the system is stable is


\begin{enumerate}[label=(\Alph*)]
    \item $-2.0 < K < 0.5$
    \item $0 < K < 0.5$
    \item $0< K< \infty$
    \item \textbf{(Correct)} $0.5< K< \infty$
\end{enumerate}

\noindent \textbf{Correct Answer:} (D)

\begin{tcolorbox}[solbox, title=Detailed Solution -- Question 58]
\[
\begin{array}{c|cc} s^{3}& 1 & (2 k+3) \\ s^{2}&2 k & 4 \\ s&\frac{4 k^{2}+6 k-4}{2 k} & 0\\ s^{0}&4 \end{array}
\]

 For stability, $K > 0$ 

\[
\begin{aligned} 4 K^{2}+6 K-4& > 0 \\ 2 K^{2}+3 K-2& > 0 \\ (2 K-1)(K+2)& > 0 \\ \Rightarrow \quad K& > \frac{1}{2} or, k < -2 \\ \text{Hence, }\quad 0.5& < K < \infty \end{aligned}
\]
\end{tcolorbox}

\vspace{1.5em}
\hrule
\vspace{1.5em}

\section*{Question 59: Time Response Analysis (GATE EC 2016-SET-3)}
\addcontentsline{toc}{section}{Question 59: Time Response Analysis (GATE EC 2016-SET-3)}
\noindent \textbf{\color{primary}MCQ} \hfill \textbf{\color{secondary}1 Mark}


\noindent For the unity feedback control system shown in the figure, the open-loop transfer function G(s) is given as  

$G(s)=\frac{2}{s(s+1)}$

. The steady state error $e_{ss}$ due to a unit step input is

\begin{center}\includegraphics[max width=0.85\linewidth]{images/img\_73.jpg}\end{center}

\begin{center}\includegraphics[max width=0.85\linewidth]{images/img\_73.jpg}\end{center}


\begin{enumerate}[label=(\Alph*)]
    \item \textbf{(Correct)} 0
    \item 0.5
    \item 1
    \item $\infty$
\end{enumerate}

\noindent \textbf{Correct Answer:} (A)

\begin{tcolorbox}[solbox, title=Detailed Solution -- Question 59]
\[
\begin{aligned} E(s) &=\frac{R(s)}{1+G(s) H(s)} \\ R(s) &=\frac{1}{s} ; G(s)=\frac{2}{s(s+1)}; \\ H(s) &=1 \\ E(s) &=\frac{\frac{1}{s}}{1+\frac{s}{s(s+1)}}=\frac{s+1}{s^{2}+s+2} \\ e_{s s} &=\lim _{s \rightarrow 0} s E(s) \\ &=\lim _{s \rightarrow 0} \frac{s(s+1)}{s^{2}+s+2} \\ e_{s s} &=0 \end{aligned}
\]
\end{tcolorbox}

\vspace{1.5em}
\hrule
\vspace{1.5em}

\section*{Question 60: Basics, Block Diagrams \& SFGs (GATE EC 2016-SET-3)}
\addcontentsline{toc}{section}{Question 60: Basics, Block Diagrams \& SFGs (GATE EC 2016-SET-3)}
\noindent \textbf{\color{primary}MCQ} \hfill \textbf{\color{secondary}1 Mark}


\noindent The block diagram of a feedback control system is shown in the figure. The overall closed-loop gain G of the system is

\begin{center}\includegraphics[max width=0.85\linewidth]{images/img\_6.jpg}\end{center}

\begin{center}\includegraphics[max width=0.85\linewidth]{images/img\_6.jpg}\end{center}


\begin{enumerate}[label=(\Alph*)]
    \item $G=\frac{G_{1}G_{2}}{1+G_{1}H_{1}}$
    \item \textbf{(Correct)} $G=\frac{G_{1}G_{2}}{1+G_{1}G_{2}+G_{1}H_{1}}$
    \item $G=\frac{G_{1}G_{2}}{1+G_{1}G_{2}H_{1}}$
    \item $G=\frac{G_{1}G_{2}}{1+G_{1}G_{2}+G_{1}G_{2}H_{1}}$
\end{enumerate}

\noindent \textbf{Correct Answer:} (B)

\begin{tcolorbox}[solbox, title=Detailed Solution -- Question 60]
\begin{center}\includegraphics[max width=0.85\linewidth]{images/img\_122.jpg}\end{center}

\begin{center}\includegraphics[max width=0.85\linewidth]{images/img\_122.jpg}\end{center}

\[
\begin{aligned} \frac{Y(s)}{X(s)} &=\frac{\frac{G_{2} G_{1}}{1+G_{1} H_{1}}}{1+\frac{G_{2} G_{1}}{1+G_{1} H_{1}}} \\ &=\frac{G_{2} G_{1}}{1+G_{1} G_{2}+G_{1} H_{1}} \end{aligned}
\]
\end{tcolorbox}

\vspace{1.5em}
\hrule
\vspace{1.5em}

\section*{Question 61: Frequency Response Analysis (GATE EC 2016-SET-2)}
\addcontentsline{toc}{section}{Question 61: Frequency Response Analysis (GATE EC 2016-SET-2)}
\noindent \textbf{\color{primary}NAT} \hfill \textbf{\color{secondary}2 Marks}


\noindent The asymptotic Bode phase plot of $G(s)=\frac{1}{(s+0.1)(s+10)(s+p_{1})}$, with k and $p_{1}$ both positive, is shown below. The value of $p_{1}$ is \underline{\hspace{1.5cm}} 

\begin{center}\includegraphics[max width=0.85\linewidth]{images/img\_116.jpg}\end{center}

\begin{center}\includegraphics[max width=0.85\linewidth]{images/img\_116.jpg}\end{center}


\begin{tcolorbox}[solbox, title=Detailed Solution -- Question 61]
\[
\begin{aligned} G(s)=\frac{K}{(s+0.1)(s+10)(s+p_{1})} \\ \angle G(s)=-\tan ^{-1} \frac{\omega}{0.1}-\tan ^{-1} \frac{\omega}{10}-\tan ^{-1} \frac{\omega}{p_{1}} \\ .\angle G(s)|_{\omega=1}=-\tan ^{-1} \frac{1}{0.1}-\tan ^{-1} \frac{1}{10}-\tan ^{-1} \frac{1}{p_{1}} \\ =-135^{\circ} \\ -\tan ^{-1} 10-\tan ^{-1} 0.1-\tan ^{-1} \frac{1}{p_{1}}=-135^{\circ} \\ -84.28^{\circ}-5.71^{\circ}-\tan ^{-1} \frac{1}{p_{1}}=-135^{\circ} \\ -\tan ^{-1} \frac{1}{p_{1}}=-135^{\circ}+90^{\circ}\\ \tan ^{-1} \frac{1}{p_{1}} =45^{\circ} \\ \frac{1}{p_{1}} =1 \Rightarrow p_{1}=1 \end{aligned}
\]
\end{tcolorbox}

\vspace{1.5em}
\hrule
\vspace{1.5em}

\section*{Question 62: Frequency Response Analysis (GATE EC 2016-SET-2)}
\addcontentsline{toc}{section}{Question 62: Frequency Response Analysis (GATE EC 2016-SET-2)}
\noindent \textbf{\color{primary}NAT} \hfill \textbf{\color{secondary}2 Marks}


\noindent In the feedback system shown below $G(s)=\frac{1}{(s+1)(s+2)(s+3)}$

\begin{center}\includegraphics[max width=0.85\linewidth]{images/img\_93.jpg}\end{center}

\begin{center}\includegraphics[max width=0.85\linewidth]{images/img\_93.jpg}\end{center}
 
The positive value of k for which the gain margin of the loop is exactly 0 dB and the phase margin of the loop is exactly zero degree is \underline{\hspace{1.5cm}}


\begin{tcolorbox}[solbox, title=Detailed Solution -- Question 62]
$1+G(s) H(s)=1+\frac{K}{(s+1)(s+2)(s+3)}=0$
$(s+1)\left(s^{2}+5 s+6\right)+K=0$
$s^{3}+5 s^{2}+6 s+s^{2}+5 s+6+k=0$
$\Rightarrow s^{3}+6 s^{2}+11 s+6+K=$
 Gain margin =0dB and phase margin $=0^{\circ}$
 It implies marginal stable system 
 By Routh Array

\[
\begin{array}{r|rr} s^{3} & 1 & 11 \\ s^{2} & 6 &(6+k) \\ s & \frac{66-6-K}{6} & 0 \\ s^{\circ} & 6+K & \end{array}
\]

 For marginal stable system, 
$60-K=0$
$K=60$
\end{tcolorbox}

\vspace{1.5em}
\hrule
\vspace{1.5em}

\section*{Question 63: Time Response Analysis (GATE EC 2016-SET-2)}
\addcontentsline{toc}{section}{Question 63: Time Response Analysis (GATE EC 2016-SET-2)}
\noindent \textbf{\color{primary}NAT} \hfill \textbf{\color{secondary}2 Marks}


\noindent In the feedback system shown below $G(s)=\frac{1}{(s^{2}+2s)}$. The step response of the closed-loop system should have minimum settling time and have no overshoot.

\begin{center}\includegraphics[max width=0.85\linewidth]{images/img\_36.jpg}\end{center}

\begin{center}\includegraphics[max width=0.85\linewidth]{images/img\_36.jpg}\end{center}

 The required value of gain k to achieve this is \underline{\hspace{1.5cm}}


\begin{tcolorbox}[solbox, title=Detailed Solution -- Question 63]
\[
\begin{aligned} G(s)&=\frac{1}{s^{2}+2 s} \\ \frac{Y(s)}{R(s)}&=\frac{\frac{K}{s^{2}+2 s}}{1+\frac{K}{s^{2}+2 s}}=\frac{K}{s^{2}+2 s+K} \end{aligned}
\]

 Minimum settling time and no overshoot implies critical damping i.e.

\[
\begin{aligned} \xi &=1 \\ \omega_{n} &=\sqrt{K} \\ \therefore \quad 2 \times \xi \cdot \omega_{n} &=2 \\ \omega_{n} &=1 \\ \Rightarrow \quad \sqrt{K} &=1 \text { or } K=1 \end{aligned}
\]
\end{tcolorbox}

\vspace{1.5em}
\hrule
\vspace{1.5em}

\section*{Question 64: Frequency Response Analysis (GATE EC 2016-SET-2)}
\addcontentsline{toc}{section}{Question 64: Frequency Response Analysis (GATE EC 2016-SET-2)}
\noindent \textbf{\color{primary}MCQ} \hfill \textbf{\color{secondary}1 Mark}


\noindent The number and direction of encirclements around the point -1+j0 in the complex plane by the Nyquist plot of $G(s)=\frac{1-s}{4+2s}$ is


\begin{enumerate}[label=(\Alph*)]
    \item \textbf{(Correct)} zero
    \item one, anti-clockwise
    \item one, clockwise.
    \item two, clockwise.
\end{enumerate}

\noindent \textbf{Correct Answer:} (A)

\begin{tcolorbox}[solbox, title=Detailed Solution -- Question 64]
$G(s)=\frac{1-s}{4+2 s}$
 and $\angle G(s)=-\tan ^{-1} \omega-\tan ^{-1} \frac{\omega}{2}$
 at $s=0, \quad|G(s)|=\frac{1}{4}=0.25$
 and $.\angle G(s)|_{\omega=0}=0^{\circ}$
 at $s=\infty, \quad|G(s)|=-\frac{1}{2}=-0.5$
 and $.\angle G(s)|_{m \times \infty}=-180^{\circ}$
 Nyquist plot is

\begin{center}\includegraphics[max width=0.85\linewidth]{images/img\_40.jpg}\end{center}

\begin{center}\includegraphics[max width=0.85\linewidth]{images/img\_40.jpg}\end{center}

 Hence, number of encirclements of (-1 + j0) = 0
\end{tcolorbox}

\vspace{1.5em}
\hrule
\vspace{1.5em}

\section*{Question 65: Time Response Analysis (GATE EC 2016-SET-2)}
\addcontentsline{toc}{section}{Question 65: Time Response Analysis (GATE EC 2016-SET-2)}
\noindent \textbf{\color{primary}NAT} \hfill \textbf{\color{secondary}1 Mark}


\noindent The response of the system $G(s)=\frac{s-2}{(s+1)(s+3)}$ to the unit step input u(t) is y(t). The value of $\frac{dy}{dt}$ at t=$0^{+}$ is \underline{\hspace{1.5cm}}


\begin{tcolorbox}[solbox, title=Detailed Solution -- Question 65]
\[
\begin{aligned} L(\frac{d y}{d t}) &=(s Y(s)-y(0)) \\ Y(s) &=G(s) \times \frac{1}{s}=\frac{s-2}{s(s+1)(s+3)} \\ y(0)&=\lim _{s arrow \infty} s Y(s) \\ \text{(Applying } & \text{initial value theorem) } \\ &=\lim _{s arrow \infty} \frac{s-2}{(s+1)(s+3)} \\ &=\frac{(1-\frac{2}{s})}{s(1+\frac{1}{s})(1+\frac{3}{s})} \\ y(0)&=0 \\ L(\frac{d y}{d t}) &=s Y(s)=\frac{s \times(s-2)}{s(s+1)(s+3)} \\ &=\frac{s-2}{(s+1)(s+3)} \\ .\frac{d y}{d t}|_{t=0} &=\lim _{s arrow \infty} s L(\frac{d y}{d t}) \\ &=\lim _{s arrow \infty} \frac{s \times(s-2)}{(s+1)(s+3)}\\ &=\frac{(1-\frac{2}{s})}{(1+\frac{1}{s})(1+\frac{3}{s})}=1 \end{aligned}
\]
\end{tcolorbox}

\vspace{1.5em}
\hrule
\vspace{1.5em}

\section*{Question 66: Stability Analysis \& Routh-Hurwitz (GATE EC 2016-SET-1)}
\addcontentsline{toc}{section}{Question 66: Stability Analysis \& Routh-Hurwitz (GATE EC 2016-SET-1)}
\noindent \textbf{\color{primary}NAT} \hfill \textbf{\color{secondary}2 Marks}


\noindent The transfer function of a linear time invariant system is given by  
$H(s)=2s^{4}-5s^{3}+5s-2$ 
 The number of zeros in the right half of the s-plane is \underline{\hspace{1.5cm}}


\begin{tcolorbox}[solbox, title=Detailed Solution -- Question 66]
$2 s^{4}-5 s^{3}+5 s-2=0$
 By Routh Array, 

\[
\begin{array}{c|ccc} s^{4} & 2 & 0 & -2 \\ s^{3} & -5 & 5 & \\ s^{2} & 2 & -2 & \\ s^{1} & 0(2) & & \\ s^{0} & -2 & & \end{array}
\]

 Number at sign changes = number at roots (zeros)  in right half of s-plane =3
\end{tcolorbox}

\vspace{1.5em}
\hrule
\vspace{1.5em}

\section*{Question 67: Time Response Analysis (GATE EC 2016-SET-1)}
\addcontentsline{toc}{section}{Question 67: Time Response Analysis (GATE EC 2016-SET-1)}
\noindent \textbf{\color{primary}NAT} \hfill \textbf{\color{secondary}2 Marks}


\noindent The open-loop transfer function of a unity-feedback control system is given by 
 
$G(s)=\frac{K}{s(s+2)}$
 
 For the peak overshoot of the closed-loop system to a unit step input to be 10%, the value of K is\underline{\hspace{1.5cm}}


\begin{tcolorbox}[solbox, title=Detailed Solution -- Question 67]
$G(s)=\frac{K}{s(s+2)} ; H(s)=1$
 Characteristic equation $=1+G(s) H(s)=0$

\[
\begin{aligned} 1+\frac{K}{s(s+2)} &=0 \\ s^{2}+2 s+K &=0 \Rightarrow \omega_{n}=\sqrt{K} \\ 2 \xi \omega_{n} &=2 ; \xi=\frac{1}{\sqrt{K}} \\ M_{p}&=e^{-\pi\xi/\sqrt{1-\xi^{2}}}-0.1\\ -\frac{\pi \xi}{\sqrt{1-\xi^{2}}} &=\ln (0.1) \\ \Rightarrow \quad\frac{\pi \xi}{\sqrt{1-\xi^{2}}} &=2.3 \\ \pi^{2} \xi^{2} &=(2.3)^{2}\left(1-\xi^{2}\right) \\ 15.16 \xi^{2} &=(2.3)^{2} \\ \Rightarrow \qquad \xi&=0.59 \\ \Rightarrow \quad \text{Also}, \quad K&=\frac{1}{\xi^{2}}=2.8 \end{aligned}
\]
\end{tcolorbox}

\vspace{1.5em}
\hrule
\vspace{1.5em}

\section*{Question 68: Root Locus Technique (GATE EC 2016-SET-1)}
\addcontentsline{toc}{section}{Question 68: Root Locus Technique (GATE EC 2016-SET-1)}
\noindent \textbf{\color{primary}NAT} \hfill \textbf{\color{secondary}2 Marks}


\noindent The open-loop transfer function of a unity-feedback control system is 
$G(s)=\frac{K}{s^{2}+5s+5}$ 
The value of K at the breakaway point of the feedback control system's root-locus plot is \underline{\hspace{1.5cm}}


\begin{tcolorbox}[solbox, title=Detailed Solution -- Question 68]
Characteristic equation  $1+G(s)H(s)=0$

$1+\frac{K}{s^2+5s+5}=0$

$K=-s^2-5s-5$
For break away point $\frac{dK}{ds}=0$
$\frac{dK}{ds}=-2s+5=0\;\Rightarrow \; s=-2.5$
According to magnitude condition,

\[
\begin{aligned} |G(s)H(s)|_{s=-2.5} &=1 \\ |G(s)H(s)|_{s=-2.5} &=\frac{|K|}{|(2.5)^2 +5 \times -2.5+5|}=1 \\ |K|&=|(6.25+5-12.5)|=1.25 \\ K&=\pm 1.25 \end{aligned}
\]
\end{tcolorbox}

\vspace{1.5em}
\hrule
\vspace{1.5em}

\section*{Question 69: Frequency Response Analysis (GATE EC 2016-SET-1)}
\addcontentsline{toc}{section}{Question 69: Frequency Response Analysis (GATE EC 2016-SET-1)}
\noindent \textbf{\color{primary}MCQ} \hfill \textbf{\color{secondary}1 Mark}


\noindent A closed-loop control system is stable if the Nyquist plot of the corresponding open-loop transfer function


\begin{enumerate}[label=(\Alph*)]
    \item \textbf{(Correct)} encircles the s-plane point (-1 + j0) in the counterclockwise direction as many times as the number of right-half s-plane poles.
    \item encircles the s-plane point (0 - j1) in the clockwise direction as many times as the number of right-half s-plane poles.
    \item encircles the s-plane point (-1 + j0) in the counterclockwise direction as many times as the number of left-half s-plane poles
    \item encircles the s-plane point (-1 + j0) in the counterclockwise direction as many times as the number of right-half s-plane zeros.
\end{enumerate}

\noindent \textbf{Correct Answer:} (A)

\begin{tcolorbox}[solbox, title=Detailed Solution -- Question 69]
N= P-Z

N = Number of encirclements of (-1 + /J). It is
positive if nyquist plot encircles the point -1 + j0
in counterclockwise direction.

Z = Number of closed loop poles lying in the right
half of s-plane

P = Number of open loop poles lying in right half
of s-plane

For stability Z = 0 $\Rightarrow$ N = P
\end{tcolorbox}

\vspace{1.5em}
\hrule
\vspace{1.5em}

\section*{Question 70: Stability Analysis \& Routh-Hurwitz (GATE EC 2016-SET-1)}
\addcontentsline{toc}{section}{Question 70: Stability Analysis \& Routh-Hurwitz (GATE EC 2016-SET-1)}
\noindent \textbf{\color{primary}MCQ} \hfill \textbf{\color{secondary}1 Mark}


\noindent Match the inferences X, Y, and Z, about a system, to the corresponding properties of the elements of first column in Routh?s Table of the system characteristic equation. 

\begin{center}\includegraphics[max width=0.85\linewidth]{images/img\_138.jpg}\end{center}

\begin{center}\includegraphics[max width=0.85\linewidth]{images/img\_138.jpg}\end{center}


\begin{enumerate}[label=(\Alph*)]
    \item X$\rightarrow$P, Y$\rightarrow$Q, Z$\rightarrow$R
    \item X$\rightarrow$Q, Y$\rightarrow$P, Z$\rightarrow$R
    \item X$\rightarrow$R, Y$\rightarrow$Q, Z$\rightarrow$P
    \item \textbf{(Correct)} X$\rightarrow$P, Y$\rightarrow$R, Z$\rightarrow$Q
\end{enumerate}

\noindent \textbf{Correct Answer:} (D)

\begin{tcolorbox}[solbox, title=Detailed Solution -- Question 70]
When all elements are positive, the system is
stable. When any element is zero, the test breaks
down. When there is change in sign of
coefficients, the system is unstable.
\end{tcolorbox}

\vspace{1.5em}
\hrule
\vspace{1.5em}

\section*{Question 71: Stability Analysis \& Routh-Hurwitz (GATE EC 2015-SET-3)}
\addcontentsline{toc}{section}{Question 71: Stability Analysis \& Routh-Hurwitz (GATE EC 2015-SET-3)}
\noindent \textbf{\color{primary}NAT} \hfill \textbf{\color{secondary}2 Marks}


\noindent The characteristic equation of an LTI system is given by $F(s)=s^{5}+2s^{4}+3s^{3}+6s^{2}-4s-8=0$ The number of roots that lie strictly in the left half s-plane is \underline{\hspace{1.5cm}}.


\begin{tcolorbox}[solbox, title=Detailed Solution -- Question 71]
Using Routh's tabular or:

\[
\begin{array}{c|lll} s^{5} & 1 & 3 & -4 \\ s^{4} & 2 & 6 & -8 \\ s^{3} & 0 8 & 0(12) & 0 \\ s^{2} & 3 & -8 & 0 \\ s^{1} & \frac{100}{3} & 0 & 0 \\ s^{0} & -8 & & \end{array}
\]

 Auxiliary equation:

\[
\begin{aligned} A &=2 s^{4}+6 s^{2}-8=0 \\ \frac{d A}{d s} &=8 s^{3}+12 s=0 \end{aligned}
\]

 As one time row of zero occur in Routh's table therefore a pair of imaginary pole exists, also  number of sign change in Routh's table is one. Hence, two poles lie strictly in the left half of s-plane.
\end{tcolorbox}

\vspace{1.5em}
\hrule
\vspace{1.5em}

\section*{Question 72: Root Locus Technique (GATE EC 2015-SET-3)}
\addcontentsline{toc}{section}{Question 72: Root Locus Technique (GATE EC 2015-SET-3)}
\noindent \textbf{\color{primary}NAT} \hfill \textbf{\color{secondary}2 Marks}


\noindent For the system shown in the figure, s = -2.75 lies on the root locus if K is \underline{\hspace{1.5cm}}.

\begin{center}\includegraphics[max width=0.85\linewidth]{images/img\_61.jpg}\end{center}

\begin{center}\includegraphics[max width=0.85\linewidth]{images/img\_61.jpg}\end{center}


\begin{tcolorbox}[solbox, title=Detailed Solution -- Question 72]
The open loop transfer function of the system:
$G(s) H(s)=\frac{K(s+3) \times 10}{(s+2)}$
 In order to lie a point s=-2.75, the angle condition must be satisfy and therefore the magnitude condition is

\[
\begin{aligned} &\mid G(s) H(s) \|_{s=-275}=1 \\ &=.\frac{10 K \sqrt{s^{2}+3^{2}}}{\sqrt{s^{2}+2^{2}}}|_{s=-2.75}=1 \\ &=\frac{10 K \sqrt{(3-2.75)^{2}}}{\sqrt{(2-2.75)^{2}}}=1 \\ &=\frac{0.25 \times 10 K}{0.75} or K=0.3 \end{aligned}
\]
\end{tcolorbox}

\vspace{1.5em}
\hrule
\vspace{1.5em}

\section*{Question 73: Basics, Block Diagrams \& SFGs (GATE EC 2015-SET-3)}
\addcontentsline{toc}{section}{Question 73: Basics, Block Diagrams \& SFGs (GATE EC 2015-SET-3)}
\noindent \textbf{\color{primary}NAT} \hfill \textbf{\color{secondary}2 Marks}


\noindent The position control of a DC servo-motor is given in the figure. The values of the parameters are $K_{T}=1 N-m/A$, $R_{a}=1 \Omega$, $L_{a}=0.1H,J=5kg-m^{2}$, $B=1 N-m(rad/sec)$ and $K_b=1V/(rad/sec)$ . The steady-state position response (in radians) due to unit impulse disturbance torque $T_{d}$ is\underline{\hspace{1.5cm}}.

\begin{center}\includegraphics[max width=0.85\linewidth]{images/img\_123.jpg}\end{center}

\begin{center}\includegraphics[max width=0.85\linewidth]{images/img\_123.jpg}\end{center}


\begin{tcolorbox}[solbox, title=Detailed Solution -- Question 73]
The transfer function due to the disturbance torque $T_{d}(s) \text { is }$

\[
\begin{aligned} \frac{\theta(s)}{T_{d}(s)}&=\frac{-\frac{1}{(J s+B)} \times \frac{1}{s}}{1+\left(\frac{1}{J s+B}\right)\left(\frac{K_{b} K_{T}}{R_{a}+L_{a} s}\right)} \\ &=\frac{-\left(R_{a}+L_{a} s\right) \cdot \frac{1}{s}}{\left(R_{a}+L_{a} s\right)(J s+B)+K_{b} K_{T}} \end{aligned}
\]

 The steady value of response for unit impulse input

\[
\begin{aligned} &=\frac{-s\left(R_{a}+L_{a} s\right) \cdot \frac{1}{s}}{\left(R_{a}+L_{a} s\right)(J s+B)+K_{b} K_{T}}.T_{d}\\ &=-\frac{R_{a}}{R_{a} B+K_{b} K_{T}} \cdot 1\\ \text{Given:}\\ K_{T}&=1 \mathrm{N}-\mathrm{m} / \mathrm{A} R_{a}=1 \Omega \\ B&=1 \mathrm{N}-\mathrm{m} / \mathrm{rad} / \mathrm{sec}\\ \text{and }\quad K_{b}&=1 \mathrm{V} / \mathrm{rad} / \mathrm{sec} \end{aligned}
\]

 Substituting the given values into above equation,  we get
$\theta(0)=-\frac{1}{2}=-0.5 \mathrm{rad}$
\end{tcolorbox}

\vspace{1.5em}
\hrule
\vspace{1.5em}

\section*{Question 74: State Space Analysis (GATE EC 2015-SET-3)}
\addcontentsline{toc}{section}{Question 74: State Space Analysis (GATE EC 2015-SET-3)}
\noindent \textbf{\color{primary}MCQ} \hfill \textbf{\color{secondary}2 Marks}


\noindent A network is described by the state model as 
$\dot{x_1}=2x_{1}-x_{2}+3u$
$\dot{x_2}=-4x_{2}-u$
$y=3x_{1}-2x_{2}$ 
 The transfer function $H(s)(=\frac{Y(s)}{U(s)})$ is


\begin{enumerate}[label=(\Alph*)]
    \item \textbf{(Correct)} $\frac{11s+35}{(s-2)(s+4)}$
    \item $\frac{11s-35}{(s-2)(s+4)}$
    \item $\frac{11s+38}{(s-2)(s+4)}$
    \item $\frac{11s-38}{(s-2)(s+4)}$
\end{enumerate}

\noindent \textbf{Correct Answer:} (A)

\begin{tcolorbox}[solbox, title=Detailed Solution -- Question 74]
State matrix form:

\[
\begin{array}{l} \dot{x}=\left[\begin{array}{cc} 2 & -1 \\ 0 & -4 \end{array}\right] x+\left[\begin{array}{c} 3 \\ -1 \end{array}\right] u \\ y=\left[\begin{array}{ll} 3 & -2 \end{array}\right] x \end{array}
\]

 The transfer function from tate matrix can be calculated as:
$T(S)=C(s l-A]^{-1} B+D$
 Here,

\[
[s l-A]^{-1}=\left[\begin{array}{cc} s-2 & 1 \\ 0 & s+4 \end{array}\right]^{-1}
\]

\[
\begin{aligned} &=\frac{1}{(s-2)(s+4)}\left[\begin{array}{cc} (s+4) & -1 \\ 0 & (s-2) \end{array}\right] \\ \Rightarrow &C[s /-A]^{-1} B \\ =& \frac{1}{(s-2)(s+4)}[3-2]\left[\begin{array}{cc} (s+4) & -1 \\ 0 & (s-2) \end{array}\right]\left[\begin{array}{c} 3 \\ -1 \end{array}\right] \\ =& \frac{1}{(s-2)(s+4)}[3-2]\left[\begin{array}{c} 3 s+12+1 \\ -s+2 \end{array}\right] \\ =& \frac{1}{(s-2)(s+4)}[9 s+39+2 s-4] \\ =& \frac{11 s+35}{(s-2)(s+4)} \end{aligned}
\]
\end{tcolorbox}

\vspace{1.5em}
\hrule
\vspace{1.5em}

\section*{Question 75: Compensators \& Controllers (GATE EC 2015-SET-3)}
\addcontentsline{toc}{section}{Question 75: Compensators \& Controllers (GATE EC 2015-SET-3)}
\noindent \textbf{\color{primary}MCQ} \hfill \textbf{\color{secondary}1 Mark}


\noindent The transfer function of a first-order controller is given as 

$G_{c}(s)=\frac{K(s+a)}{s+b}$ 

 where K, a and b are positive real numbers. The condition for this controller to act as a phase lead compensator is


\begin{enumerate}[label=(\Alph*)]
    \item \textbf{(Correct)} $a < b$
    \item $a > b$
    \item $K < ab$
    \item $K > ab$
\end{enumerate}

\noindent \textbf{Correct Answer:} (A)

\begin{tcolorbox}[solbox, title=Detailed Solution -- Question 75]
For phase lead compensator

\[
\begin{aligned} G_{c}(s)&=\frac{(1+s \tau)}{(1+\alpha s \tau)} \quad ; \quad \alpha < 1\\ \text{Here,}\quad \tau&=\frac{1}{a}\\ \text{and}\quad \alpha \tau&=\frac{1}{b}\\ \text{or,}\quad \alpha&=\frac{a}{b} < 1\\ \text{or,}\quad a& < b \end{aligned}
\]
\end{tcolorbox}

\vspace{1.5em}
\hrule
\vspace{1.5em}

\section*{Question 76: Frequency Response Analysis (GATE EC 2015-SET-3)}
\addcontentsline{toc}{section}{Question 76: Frequency Response Analysis (GATE EC 2015-SET-3)}
\noindent \textbf{\color{primary}NAT} \hfill \textbf{\color{secondary}1 Mark}


\noindent The phase margin (in degrees) of the system $G(s)=\frac{10}{s(s+10)}$ is \underline{\hspace{1.5cm}}.


\begin{tcolorbox}[solbox, title=Detailed Solution -- Question 76]
\[
\begin{aligned} G(s) &=\frac{10}{s(s+10)} \\ P M &=180^{\circ}+.\angle G(\omega) H(\omega)|_{\omega=\omega_{g c}} \end{aligned}
\]

 The gain cross over frequency can be calculated by

\[
\begin{aligned} |G(\omega) H(\omega)|&=1 \\ \frac{10}{\omega \sqrt{\omega^{2}+100}}&=1 \\ \omega_{g c}^{4}+100 \omega_{g c}^{2}&-100=0\\ \text { Or } \quad \omega_{g c}^{2}&=0.99, \quad-100.99 \\ \text { Or } \quad \omega_{g c}&=0.99 \text { rad/sec. } \\ \therefore P M &=180^{\circ}+\left(-90^{\circ}-\tan ^{-1} \frac{\omega_{g c}}{10}\right) \\ &=180^{\circ}-90^{\circ}-\tan ^{-1} \frac{0.99}{10} \\ & \simeq 84.364^{\circ} \end{aligned}
\]
\end{tcolorbox}

\vspace{1.5em}
\hrule
\vspace{1.5em}

\section*{Question 77: Frequency Response Analysis (GATE EC 2015-SET-3)}
\addcontentsline{toc}{section}{Question 77: Frequency Response Analysis (GATE EC 2015-SET-3)}
\noindent \textbf{\color{primary}NAT} \hfill \textbf{\color{secondary}1 Mark}


\noindent Consider the Bode plot shown in the figure. Assume that all the poles and zeros are real-valued. 

\begin{center}\includegraphics[max width=0.85\linewidth]{images/img\_27.jpg}\end{center}

\begin{center}\includegraphics[max width=0.85\linewidth]{images/img\_27.jpg}\end{center}

 The value of $f_{H}-f_{L}$ (in Hz) is \underline{\hspace{1.5cm}}.


\begin{tcolorbox}[solbox, title=Detailed Solution -- Question 77]
As we know,

\[
\begin{aligned} \text { Slope } &=\frac{y_{2}-y_{1}}{x_{2}-x_{1}} \\ \therefore \quad 40 \mathrm{dB} / \mathrm{dec} &=\frac{40-0}{\log 300-\log _{L}}\\ or, \log \left(\frac{300}{f_{L}}\right)&=\frac{40}{40}=1 f_{L}=30 \mathrm{Hz} \end{aligned}
\]

 Similarly,

\[
\begin{aligned} -40 \mathrm{dB} / \mathrm{dec}&=\frac{40-0}{\log 900-\log f_{H}}\\ \log \left(\frac{900}{f_{H}}\right) &=-1 \\ \text { or, } \quad f_{H}&=9000 \mathrm{Hz} \\ \text { Therefore, } f_{H}-f_{L}&=9000-30 =8970 \mathrm{Hz} \end{aligned}
\]
\end{tcolorbox}

\vspace{1.5em}
\hrule
\vspace{1.5em}

\section*{Question 78: Frequency Response Analysis (GATE EC 2015-SET-2)}
\addcontentsline{toc}{section}{Question 78: Frequency Response Analysis (GATE EC 2015-SET-2)}
\noindent \textbf{\color{primary}NAT} \hfill \textbf{\color{secondary}2 Marks}


\noindent The transfer function of a mass-spring-damper system is given by 

$G(s)=\frac{1}{Ms^{2}+Bs+K}$

  The frequency response data for the system are given in the following table. 
\begin{center}\includegraphics[max width=0.85\linewidth]{images/img\_77.jpg}\end{center}

\begin{center}\includegraphics[max width=0.85\linewidth]{images/img\_77.jpg}\end{center}

The unit step response of the system approaches a steady state value of \underline{\hspace{1.5cm}} .


\begin{tcolorbox}[solbox, title=Detailed Solution -- Question 78]
$\frac{\alpha(s)}{R(s)}=G(s)=\frac{1}{M s^{2}+B s+K}$
 Steady state value of response for unit step input

\[
\begin{aligned} C(s)&=\lim _{s \rightarrow 0} s \times \frac{1}{M s^{2}+B s+K} \times \frac{1}{s}=\frac{1}{K} \\ \text{Also}&\\ G(\approx 0)&=\frac{1}{K}=-18.5 \mathrm{dB} \\ &\qquad\text{(as per the given table)}\\ \therefore \frac{1}{K}&=10^{\left(-\frac{18.5}{20}\right)}=0.12 \end{aligned}
\]
\end{tcolorbox}

\vspace{1.5em}
\hrule
\vspace{1.5em}

\section*{Question 79: Time Response Analysis (GATE EC 2015-SET-2)}
\addcontentsline{toc}{section}{Question 79: Time Response Analysis (GATE EC 2015-SET-2)}
\noindent \textbf{\color{primary}MCQ} \hfill \textbf{\color{secondary}2 Marks}


\noindent The output of a standard second-order system for a unit step input is given as $y(t)=1-\frac{2}{\sqrt{3}}e^{-t}cos(\sqrt{3}t-\frac{\pi }{6})$. The transfer function of the system is


\begin{enumerate}[label=(\Alph*)]
    \item $\frac{2}{(s+2)(s+\sqrt{3})}$
    \item $\frac{1}{s^2+2s+1}$
    \item $\frac{3}{s^2+2s+3}$
    \item \textbf{(Correct)} $\frac{4}{s^2+2s+4}$
\end{enumerate}

\noindent \textbf{Correct Answer:} (D)

\begin{tcolorbox}[solbox, title=Detailed Solution -- Question 79]
The output of a second order system for unit step 
 input

\[
y(t)=1-\frac{e^{-\xi \omega_{n} t}}{\sqrt{1-\xi^{2}}} \sin \left(\omega_{d} t-\theta\right) \qquad\ldots(i)
\]

 From given data

\[
\begin{aligned} y(t)&=1-\frac{2}{\sqrt{3}} e^{-t} \cos \left(\sqrt{3} t-\frac{\pi}{6}\right)\\ \therefore \quad \alpha&=\xi \omega_{n}=1 \\ \text{and }\quad \omega_{d}&=\sqrt{3} \\ \omega_{n}&=\sqrt{\alpha^{2}+\omega_{d}^{2}}=2 \\ \xi&=\frac{1}{2} \end{aligned}
\]

 Therefore, the transfer function

\[
T(s)=\frac{\omega_{n}^{2}}{s^{2}+2 \xi \omega_{n} s+\omega_{n}^{2}}=\frac{4}{s^{2}+2 s+4}
\]
\end{tcolorbox}

\vspace{1.5em}
\hrule
\vspace{1.5em}

\section*{Question 80: State Space Analysis (GATE EC 2015-SET-2)}
\addcontentsline{toc}{section}{Question 80: State Space Analysis (GATE EC 2015-SET-2)}
\noindent \textbf{\color{primary}MCQ} \hfill \textbf{\color{secondary}2 Marks}


\noindent The state variable representation of a system is given as 
\[
x=\begin{bmatrix} 0 & 1\\ 0&-1 \end{bmatrix}x; \; x(0)=\begin{bmatrix} 1\\ 0 \end{bmatrix} \; y= \begin{bmatrix} 0 &1 \end{bmatrix}x
\]

 The response y(t) is


\begin{enumerate}[label=(\Alph*)]
    \item sin(t)
    \item $1-e^{t}$
    \item 1-cos(t)
    \item \textbf{(Correct)} 0
\end{enumerate}

\noindent \textbf{Correct Answer:} (D)

\begin{tcolorbox}[solbox, title=Detailed Solution -- Question 80]
The response of the system is given by 
$x(t)=L^{-1}\left[(s l-A)^{-1} x(0)+(s /-A)^{-1} B \cdot U(s)\right]$
 and the complete response
$y(t)=[C] x(t)$
 From the given state model,

\[
\begin{array}{l} A=\left[\begin{array}{ll}0 & 1 \\ 0 & -1\end{array}\right] ; B=0: C=[0 \quad 1] \; and \;  x(0)=\left[\begin{array}{c}1 \\ 0\end{array}\right] \\ \therefore \quad(s l-A)^{-1}=\left[\begin{array}{cc}s & -1 \\ 0 & s+1\end{array}\right]^{-1} \\ =\frac{1}{s(s+1)}\left[\begin{array}{cc}s+1 & 1 \\ 0 & s\end{array}\right] \\ (s l-A)^{-1} x(0)=\frac{1}{s(s+1)}\left[\begin{array}{cc} s+1 & 1 \\ 0 & s \end{array}\right]\left[\begin{array}{c} 1 \\ 0 \end{array}\right] \\ =\frac{1}{s(s+1)}\left[\begin{array}{c} s+1 \\ 0 \end{array}\right]=\left[\begin{array}{c} 1 / s \\ 0 \end{array}\right] \\ \Rightarrow x(t)=L^{-1}\left[(s I-A)^{-1} x(0)\right] \\ =\left[\begin{array}{c} 1 \\ 0 \end{array}\right] u(t) \end{array}
\]

and the complete response 

\[
y(t)=\left[\begin{array}{lll}0 & 1\end{array}\right] x(t)=\left[\begin{array}{ll}0 & 1\end{array}\right]\left[\begin{array}{l}1 \\ 0\end{array}\right]=0
\]
\end{tcolorbox}

\vspace{1.5em}
\hrule
\vspace{1.5em}

\section*{Question 81: Time Response Analysis (GATE EC 2015-SET-2)}
\addcontentsline{toc}{section}{Question 81: Time Response Analysis (GATE EC 2015-SET-2)}
\noindent \textbf{\color{primary}NAT} \hfill \textbf{\color{secondary}1 Mark}


\noindent A unity negative feedback system has an open-loop transfer function $G(s)=\frac{K}{s(s+10)}$. The gain for the system to have a damping ratio of 0.25 is \underline{\hspace{1.5cm}}.


\begin{tcolorbox}[solbox, title=Detailed Solution -- Question 81]
$T(s)=\frac{G(s)}{1+G(s) H(s)}=\frac{K}{s^{2}+10 s+K}$
 Comparing with standard second order transfer  function

\[
\begin{aligned} T(s)&=\frac{\omega_{n}^{2}}{s^{2}+2 \xi \omega_{n} s+\omega_{n}^{2}}\\ \text{We have, }\omega_{n}^{2}&=K \\ \text{and }\quad 2 \xi \omega_{n}&=10 \\ \because \quad \xi&=0.25 \quad (given)\\ \therefore \quad \omega_{n}&=\frac{10}{2 \times 0.25}=20 \\ \text{and }\quad k&=\omega_{n}^{2}=(20)^{2}=400 \end{aligned}
\]
\end{tcolorbox}

\vspace{1.5em}
\hrule
\vspace{1.5em}

\section*{Question 82: Basics, Block Diagrams \& SFGs (GATE EC 2015-SET-2)}
\addcontentsline{toc}{section}{Question 82: Basics, Block Diagrams \& SFGs (GATE EC 2015-SET-2)}
\noindent \textbf{\color{primary}MCQ} \hfill \textbf{\color{secondary}1 Mark}


\noindent For the signal flow graph shown in the figure, the value of $\frac{C(s)}{R(s)}$  is 

\begin{center}\includegraphics[max width=0.85\linewidth]{images/img\_20.jpg}\end{center}

\begin{center}\includegraphics[max width=0.85\linewidth]{images/img\_20.jpg}\end{center}


\begin{enumerate}[label=(\Alph*)]
    \item \[
\frac{G_{1}G_{2}G_{3}G_{4}}{1-G_{1}G_{2}H_{1}-G_{3}G_{4}H_{2}-G_{2}G_{3}H_{3}+G_{1}G_{2}G_{3}G_{4}H_{1}H_{2}}
\]
    \item \textbf{(Correct)} \[
\frac{G_{1}G_{2}G_{3}G_{4}}{1+G_{1}G_{2}H_{1}+G_{3}G_{4}H_{2}+G_{2}G_{3}H_{3}+G_{1}G_{2}G_{3}G_{4}H_{1}H_{2}}
\]
    \item \[
\frac{1}{1+G_{1}G_{2}H_{1}+G_{3}G_{4}H_{2}+G_{2}G_{3}H_{3}+G_{1}G_{2}G_{3}G_{4}H_{1}H_{2}}
\]
    \item \[
\frac{1}{1-G_{1}G_{2}H_{1}-G_{3}G_{4}H_{2}-G_{2}G_{3}H_{3}+G_{1}G_{2}G_{3}G_{4}H_{1}H_{2}}
\]
\end{enumerate}

\noindent \textbf{Correct Answer:} (B)

\begin{tcolorbox}[solbox, title=Detailed Solution -- Question 82]
Transfer function

\[
\begin{aligned} \frac{C(s)}{R(s)}&=\frac{\Sigma P_{k} \Delta_{k}}{\Delta} \\ =& \frac{G_{1} G_{2} G_{3} G_{4}}{1-\left(-G_{1} G_{2} H_{1}-G_{2} G_{4} H_{2}-G_{2} G_{3} H_{3}\right)+G_{1} G_{2} H_{1} \cdot G_{3} G_{4} H_{2}} \\ \frac{C(s)}{R(s)}&= \frac{G_{1}G_{2}G_{3}G_{4}}{1+G_{1}G_{2}H_{1}+G_{3}G_{4}H_{2}+G_{2}G_{3}H_{3}+G_{1}G_{2}G_{3}G_{4}H_{1}H_{2}} \end{aligned}
\]
\end{tcolorbox}

\vspace{1.5em}
\hrule
\vspace{1.5em}

\section*{Question 83: Basics, Block Diagrams \& SFGs (GATE EC 2015-SET-2)}
\addcontentsline{toc}{section}{Question 83: Basics, Block Diagrams \& SFGs (GATE EC 2015-SET-2)}
\noindent \textbf{\color{primary}MCQ} \hfill \textbf{\color{secondary}1 Mark}


\noindent By performing cascading and/or summing/differencing operations using transfer function blocks $G_{1}(s) \; and \;  G_{2}(s)$, one CANNOT realize a transfer function of the form


\begin{enumerate}[label=(\Alph*)]
    \item ${G_{1}(s)}{G_{1}(s)}$
    \item \textbf{(Correct)} $\frac{G_{1}(s)}{G_{1}(s)}$
    \item $G_{1}(s)(\frac{1}{G_{1}(s)}+G_{2}(s))$
    \item $G_{1}(s)(\frac{1}{G_{1}(s)}-G_{2}(s))$
\end{enumerate}

\noindent \textbf{Correct Answer:} (B)

\begin{tcolorbox}[solbox, title=Detailed Solution -- Question 83]
Given blocks has transfer function $G_{1}(s)$ and $G_{2}(s)$
 In cascade connection $G_{1}(s) \cdot G_{2}(s)$
 In parallel connection $G_{1}(s)+G_{2}(s)$
 From the given options:
 (For cascading/ summing/differncing)
 (A) $G_{1}(s) G_{2}(s) \Rightarrow$Realization is possible.
 (B) $\frac{G_{1}(s)}{G_{2}(s)} \Rightarrow$ Realization is not possible because 
 the gain of blocks given as $G_{1}(s)$ and $G_{2}(s)$
 (C) $G_{1}(s)\left(\frac{1}{G_{1}(s)}+G_{2}(s)\right)=1+G_{1}(s) G_{2}(s)$
 Realization is possible.
 (D) $G_{1}(s)\left(\frac{1}{G_{1}(s)}-G_{2}(s)\right)=1-G_{1}(s) G_{2}(s)$
 Realization is possible.
\end{tcolorbox}

\vspace{1.5em}
\hrule
\vspace{1.5em}

\section*{Question 84: Stability Analysis \& Routh-Hurwitz (GATE EC 2015-SET-1)}
\addcontentsline{toc}{section}{Question 84: Stability Analysis \& Routh-Hurwitz (GATE EC 2015-SET-1)}
\noindent \textbf{\color{primary}MCQ} \hfill \textbf{\color{secondary}2 Marks}


\noindent A plant transfer function is given as $G(s)=(K_{P}+\frac{K_{I}}{s})\frac{1}{s(s+2)}$. When the plant operates in a unity feedback configuration, the condition for the stability of the closed loop system is


\begin{enumerate}[label=(\Alph*)]
    \item \textbf{(Correct)} $K_{P}> \frac{K_{I}}{2} >  0$
    \item $2K_{I}> K_{P} > 0$
    \item $2K_{I} < K_{P}$
    \item $2K_{I} > K_{P}$
\end{enumerate}

\noindent \textbf{Correct Answer:} (A)

\begin{tcolorbox}[solbox, title=Detailed Solution -- Question 84]
$G(s)=\left(K_{p}+\frac{K_{I}}{s}\right)\left(\frac{1}{s(s+2)}\right)$
 The closed loop transfer function for unity feedback

\[
\begin{aligned} \frac{G(s)}{1+G(s)} &=\frac{\left(K_{p} s+K_{I}\right)}{s^{2}(s+2)+\left(K_{p} s+K_{I}\right)} \\ &=\frac{K_{p} s+K_{I}}{s^{3}+2 s^{2}+K_{p} s+K_{I}} \end{aligned}
\]

 Using Routh's tabular form:

\[
\begin{array}{c|cc} s^{3} & 1 & K_{p} \\ s^{2} & 2 & K_{I} \\ s^{1} & \frac{2 K_{p}-K_{I}}{2} & 0 \\ s^{0} & K_{p} & \end{array}
\]

 For system to be stable;

\[
\begin{aligned} K_{p}& > 0\\ \text{and }\frac{2 K_{p}-K_{I}}{2}& > 0 \\ \text{or} \quad 2 K_{p}-K_{I}& > 0 \\ \text{or} \quad K_{p}& > \frac{K_{I}}{2} > 0 \end{aligned}
\]
\end{tcolorbox}

\vspace{1.5em}
\hrule
\vspace{1.5em}

\section*{Question 85: Compensators \& Controllers (GATE EC 2015-SET-1)}
\addcontentsline{toc}{section}{Question 85: Compensators \& Controllers (GATE EC 2015-SET-1)}
\noindent \textbf{\color{primary}NAT} \hfill \textbf{\color{secondary}2 Marks}


\noindent A lead compensator network includes a parallel combination of R and C in the feed-forward path. If the transfer function of the compensator is $G(s)=\frac{s+2}{s+4}$ , the value of RC is \underline{\hspace{1.5cm}}.


\begin{tcolorbox}[solbox, title=Detailed Solution -- Question 85]
$G_{C}(s)=\frac{s+2}{s+4} \quad\dots(i)$
 For lead compensator,

\begin{center}\includegraphics[max width=0.85\linewidth]{images/img\_52.jpg}\end{center}

\begin{center}\includegraphics[max width=0.85\linewidth]{images/img\_52.jpg}\end{center}

 Transfer function $=\frac{1+s \tau}{1+\alpha s \tau} \quad\dots(ii)$
 Where, 
$\tau=$ Lead time constant $=R_{1} C$
 and $\alpha=\frac{R_{2}}{R_{1}+R_{2}}$
 Comparing equation (i) and (ii), we get
$\tau=\frac{1}{2} \quad$ and $\quad \alpha \tau=\frac{1}{4}$
 or, $\alpha=\frac{1}{2}$
$\therefore$ RC time constant =0.5
\end{tcolorbox}

\vspace{1.5em}
\hrule
\vspace{1.5em}

\section*{Question 86: Root Locus Technique (GATE EC 2015-SET-1)}
\addcontentsline{toc}{section}{Question 86: Root Locus Technique (GATE EC 2015-SET-1)}
\noindent \textbf{\color{primary}NAT} \hfill \textbf{\color{secondary}2 Marks}


\noindent The open-loop transfer function of a plant in a unity feedback configuration is given as $G(s)=\frac{K(s+4)}{(s+8)(s^{2}-9)}$ . The value of the gain K(> 0) for which -1+ j2 lies on the root locus is\underline{\hspace{1.5cm}}.


\begin{tcolorbox}[solbox, title=Detailed Solution -- Question 86]
\[
\begin{aligned} G(s) &=\frac{K(s+4)}{(s+8)\left(s^{2}-9\right)} \\ &=\frac{K(s+4)}{(s+8)(s+3)(s-3)} \end{aligned}
\]

 For the point (-1+2 j) to lie on the root locus the angle condition must be satisfy 
 i.e. $\angle G(s) H(s)=\pm(2 Q+1) 180^{\circ}$
 Taking LHS:

\[
\begin{array}{l} .\angle G(s) H(s)|_{s=-1+2 j} \\ =\frac{\angle K+\angle(s+4)}{\angle(s+8)+\angle(s+3)+\angle(s-3)} \\ =\frac{\angle K+\angle(-1+2 j+4)}{\angle(-1+2 j+8)+\angle(-1+2 j+3)+\angle(-1+2 j-3)} \\ =\frac{0+\tan ^{-1}(\frac{2}{3})}{\tan ^{-1}(\frac{2}{7})+\tan ^{-1}(1)+180^{\circ}-\tan ^{-1}(\frac{1}{2})} \\ =33.69^{\circ}-15.945^{\circ}-45^{\circ}-180^{\circ}+26.565^{\circ} \\ \simeq 180^{\circ} \end{array}
\]

As angle condition is satisfied the value of system gain K can be obtained by using magnitude
 condition

\[
\begin{array}{l} i.e. |G(s) H(s)|_{s=-1+2 j}=1 \\ =\frac{K \sqrt{(-1+4)^{2}+2^{2}}}{\sqrt{(-1+8)^{2}+2^{2}} \cdot \sqrt{(-1+3)^{2}+2^{2}} \cdot \sqrt{(-1-3)^{2}+2^{2}}} \\ =1 \\ =\frac{K \sqrt{9+4}}{\sqrt{49+4} \sqrt{4+4} \sqrt{16+4}}=1 \\ =\frac{K \sqrt{13}}{\sqrt{53} \sqrt{8} \sqrt{20}}=1 \\ K=\frac{\sqrt{53 \times 8 \times 20}}{\sqrt{13}}=25.54 \end{array}
\]
\end{tcolorbox}

\vspace{1.5em}
\hrule
\vspace{1.5em}

\section*{Question 87: Frequency Response Analysis (GATE EC 2015-SET-1)}
\addcontentsline{toc}{section}{Question 87: Frequency Response Analysis (GATE EC 2015-SET-1)}
\noindent \textbf{\color{primary}MCQ} \hfill \textbf{\color{secondary}1 Mark}


\noindent The polar plot of the transfer function $G(s)=\frac{10(s+1)}{s+10}$ for $0\leq \omega <  \infty$ will be in the


\begin{enumerate}[label=(\Alph*)]
    \item \textbf{(Correct)} first quadrant
    \item second quadrant
    \item third quadrant
    \item fourth quadrant
\end{enumerate}

\noindent \textbf{Correct Answer:} (A)

\begin{tcolorbox}[solbox, title=Detailed Solution -- Question 87]
$G(s)=\frac{10(s+1)}{(s+10)}$
 For $\omega=0 \Rightarrow M_{1}=1$ and $\phi_{1}=0^{\circ}$
$\omega=0 \Rightarrow M_{2}=10$
 and $\phi_{2}=-\tan ^{-1} \infty+\tan ^{-1} \infty=0^{\circ}$
$\therefore$ Polar plot

\begin{center}\includegraphics[max width=0.85\linewidth]{images/img\_3.jpg}\end{center}

\begin{center}\includegraphics[max width=0.85\linewidth]{images/img\_3.jpg}\end{center}

As the zero is nearer to the imaginary axis hence
the direction of polar plot is clockwise.
\end{tcolorbox}

\vspace{1.5em}
\hrule
\vspace{1.5em}

\section*{Question 88: Root Locus Technique (GATE EC 2015-SET-1)}
\addcontentsline{toc}{section}{Question 88: Root Locus Technique (GATE EC 2015-SET-1)}
\noindent \textbf{\color{primary}NAT} \hfill \textbf{\color{secondary}1 Mark}


\noindent A unity negative feedback system has the open-loop transfer function $G(s)=\frac{K}{s(s+1)(s+3)}$. The value of the gain K ($>$0) at which the root locus crosses the imaginary axis is \underline{\hspace{1.5cm}}.


\begin{tcolorbox}[solbox, title=Detailed Solution -- Question 88]
$G(s)=\frac{K}{s(s+1)(s+3)}$
 The characteristic equation

\[
\begin{array}{l} =1+G(s) H(s)=0 \\ =s(s+1)(s+3)+K=0 \\ =s^{3}+4 s^{2}+3 s+K=0 \end{array}
\]

 Using Routh's tabular form

\[
\begin{array}{c|cc} s^{3} & 1 & 3 \\ s^{2} & 4 & K \\ s^{1} & \frac{12-K}{4} & 0 \\ s^{0} & K & \end{array}
\]

 In order to cross the imaginary axis, system 
 should be marginally stable

\[
\begin{aligned} \frac{12-K}{4}&=0\\ \text{or }\quad K&=12 \end{aligned}
\]
\end{tcolorbox}

\vspace{1.5em}
\hrule
\vspace{1.5em}

\section*{Question 89: Basics, Block Diagrams \& SFGs (GATE EC 2015-SET-1)}
\addcontentsline{toc}{section}{Question 89: Basics, Block Diagrams \& SFGs (GATE EC 2015-SET-1)}
\noindent \textbf{\color{primary}MCQ} \hfill \textbf{\color{secondary}1 Mark}


\noindent Negative feedback in a closed-loop control system DOES NOT


\begin{enumerate}[label=(\Alph*)]
    \item reduce the overall gain
    \item \textbf{(Correct)} reduce bandwidth
    \item improve disturbance rejection
    \item reduce sensitivity to parameter variation
\end{enumerate}

\noindent \textbf{Correct Answer:} (B)


\vspace{1.5em}
\hrule
\vspace{1.5em}

\section*{Question 90: Time Response Analysis (GATE EC 2014-SET-4)}
\addcontentsline{toc}{section}{Question 90: Time Response Analysis (GATE EC 2014-SET-4)}
\noindent \textbf{\color{primary}NAT} \hfill \textbf{\color{secondary}2 Marks}


\noindent The characteristic equation of a unity negative feedback system is $1 + KG(s) = 0$. The open loop transfer function G(s) has one pole at 0 and two poles at -1. The root locus of the system for varying K is shown in the figure. 

\begin{center}\includegraphics[max width=0.85\linewidth]{images/img\_26.jpg}\end{center}

\begin{center}\includegraphics[max width=0.85\linewidth]{images/img\_26.jpg}\end{center}

 The constant damping ratio line  for $\xi= 0.5$, intersects the root locus at point A. The distance from the origin to point A is given as 0.5. The value of K at point A is \underline{\hspace{1.5cm}}.


\begin{tcolorbox}[solbox, title=Detailed Solution -- Question 90]
According to the question 
$G(s) H(s)=\frac{K}{s(s+1)^{2}} \qquad \ldots(i)$

\begin{center}\includegraphics[max width=0.85\linewidth]{images/img\_22.jpg}\end{center}

\begin{center}\includegraphics[max width=0.85\linewidth]{images/img\_22.jpg}\end{center}

 The damping ratio is given as $\xi=0.5$
$\therefore \quad \theta=\cos ^{-1} \xi=60^{\circ}$
 Let the co-ordinates of point 'A' is (-a+jb) 
 where, $\cos 60^{\circ}=\frac{O B}{O A} \qquad \ldots(ii)$
 Also, $\quad O A=0.5$ (Given)
$\therefore O B=0.25$
 Similarly,

\[
\begin{array}{l} \sin 60^{\circ}=\frac{B A}{O A} \qquad \ldots(iii)\\ \therefore \quad B A=j b=0.433 j \end{array}
\]

 Hence, co-ordinates of $A = -0.25 + 0.43j$
$\therefore$ The value of system gain K can be obtained as
$|G(s) H(s)|=1\qquad \ldots(iv)$
 From (i) and (iv) we get $s=-0.25+j 4.33$
$K=0.37$
\end{tcolorbox}

\vspace{1.5em}
\hrule
\vspace{1.5em}

\section*{Question 91: Stability Analysis \& Routh-Hurwitz (GATE EC 2014-SET-4)}
\addcontentsline{toc}{section}{Question 91: Stability Analysis \& Routh-Hurwitz (GATE EC 2014-SET-4)}
\noindent \textbf{\color{primary}NAT} \hfill \textbf{\color{secondary}2 Marks}


\noindent Consider a transfer function $G_{p}(s)=\frac{ps^{2}+3ps-2}{s^{2}+(3+p)s+(2-p)}$ with $p$ a positive real parameter. The maximum value of $p$ until which $G_{p}$ remains stable is \underline{\hspace{1.5cm}}.


\begin{tcolorbox}[solbox, title=Detailed Solution -- Question 91]
Given transfer function
$G_{p}(s)=\frac{p s^{2}+3 p s-2}{s^{2}+(3+p) s+(2-p)} \qquad\ldots(i)$
 The characteristic equation is
$=s^{2}+(3+p) s+(2-p)\qquad\ldots(ii)$
 Using Routh Hurtwiz criterion,

\[
\begin{array}{l|ll} s^{2} & 1 & (2-P) \\ s^{1} & (3+p) & 0 \\ s^{0} & (2-p) & 0 \end{array}
\]

 For system to be remain stable

\[
\begin{aligned} 2-p & \geq 0 \\ p_{\max } &=1.99 \end{aligned}
\]
\end{tcolorbox}

\vspace{1.5em}
\hrule
\vspace{1.5em}

\section*{Question 92: State Space Analysis (GATE EC 2014-SET-4)}
\addcontentsline{toc}{section}{Question 92: State Space Analysis (GATE EC 2014-SET-4)}
\noindent \textbf{\color{primary}MCQ} \hfill \textbf{\color{secondary}2 Marks}


\noindent The state transition matrix $\phi (t)$ of a system 
\[
\begin{bmatrix} \dot{x}_{1}\\ \dot{x}_{2} \end{bmatrix}=\begin{bmatrix} 0 & 1\\ 0&0 \end{bmatrix}\begin{bmatrix} x_{1}\\ x_{2} \end{bmatrix}
\]
 is


\begin{enumerate}[label=(\Alph*)]
    \item \[
\begin{bmatrix} t & 1\\ 1&0 \end{bmatrix}
\]
    \item \[
\begin{bmatrix} 1 & 0\\ t&1 \end{bmatrix}
\]
    \item \[
\begin{bmatrix} 0 & 1\\ 1&t \end{bmatrix}
\]
    \item \textbf{(Correct)} \[
\begin{bmatrix} 1 & t\\ 0&1 \end{bmatrix}
\]
\end{enumerate}

\noindent \textbf{Correct Answer:} (D)

\begin{tcolorbox}[solbox, title=Detailed Solution -- Question 92]
\[
\begin{aligned} \text { Given } &\left[\begin{array}{l} \dot{x}_{1} \\ \dot{x}_{2} \end{array}\right]=\left[\begin{array}{ll} 0 & 1 \\ 0 & 0 \end{array}\right]\left[\begin{array}{l} x_{1} \\ x_{2} \end{array}\right] \\ A &=\left[\begin{array}{ll} 0 & 1 \\ 0 & 0 \end{array}\right] \\ [s I-A] &=\left[\begin{array}{ll} s & 0 \\ 0 & s \end{array}\right]-\left[\begin{array}{ll} 0 & 1 \\ 0 & 0 \end{array}\right]=\left[\begin{array}{ll} s & -1 \\ 0 & s \end{array}\right] \\ [s I-A]=s^{2} & \\ \phi(t) &=\mathcal{L}^{-1}[s I-A]^{-1} \\ &=\frac{1}{s^{2}}\left[\begin{array}{ll} s & 1 \\ 0 & s \end{array}\right] \\ &=L^{-1}\left[\begin{array}{cc} \frac{1}{s} & \frac{1}{s^{2}} \\ 0 & \frac{1}{s} \end{array}\right]=\left[\begin{array}{ll} 1 & t \\ 0 & 1 \end{array}\right] \end{aligned}
\]
\end{tcolorbox}

\vspace{1.5em}
\hrule
\vspace{1.5em}

\section*{Question 93: Time Response Analysis (GATE EC 2014-SET-4)}
\addcontentsline{toc}{section}{Question 93: Time Response Analysis (GATE EC 2014-SET-4)}
\noindent \textbf{\color{primary}MCQ} \hfill \textbf{\color{secondary}1 Mark}


\noindent For the second order closed-loop system shown in the figure, the natural frequency (in rad/s) is 

\begin{center}\includegraphics[max width=0.85\linewidth]{images/img\_117.jpg}\end{center}

\begin{center}\includegraphics[max width=0.85\linewidth]{images/img\_117.jpg}\end{center}


\begin{enumerate}[label=(\Alph*)]
    \item 16
    \item 4
    \item \textbf{(Correct)} 2
    \item 1
\end{enumerate}

\noindent \textbf{Correct Answer:} (C)

\begin{tcolorbox}[solbox, title=Detailed Solution -- Question 93]
The closed loop transfer function of the given 
 system is
$\frac{C(s)}{R(s)}=\frac{4}{s(s+4)+4}=\frac{4}{s^{2}+4 s+4}$
 Comparing with standard second order transfer 

\[
\begin{aligned} \text{function }\quad &\frac{K \omega_{n}^{2}}{S^{2}+2 \xi \omega_{n}+\omega_{n}^{2}} \\ \text{we get, }\quad \omega_{n}^{2}&=4 \\ \text{or }\quad\omega_{n}&=\sqrt{4}=2 \mathrm{rad} / \mathrm{sec} \end{aligned}
\]
\end{tcolorbox}

\vspace{1.5em}
\hrule
\vspace{1.5em}

\section*{Question 94: Frequency Response Analysis (GATE EC 2014-SET-4)}
\addcontentsline{toc}{section}{Question 94: Frequency Response Analysis (GATE EC 2014-SET-4)}
\noindent \textbf{\color{primary}MCQ} \hfill \textbf{\color{secondary}1 Mark}


\noindent In a Bode magnitude plot, which one of the following slopes would be exhibited at high frequencies by a $4^{th}$ order all-pole system ?


\begin{enumerate}[label=(\Alph*)]
    \item \textbf{(Correct)} $-80 dB/decade$
    \item $-40 dB/decade$
    \item $+40 dB/decade$
    \item $+80 dB/decade$
\end{enumerate}

\noindent \textbf{Correct Answer:} (A)

\begin{tcolorbox}[solbox, title=Detailed Solution -- Question 94]
A $4^{\text {th }}$  order all-pole system means that the system
must be having no zero or s-terms in numerator or
and $s^{4}$  terms in denominator or 4 -poles.

i.e.  $H(s) \propto \frac{1}{s^{4}}$

One pole exhibits -20 dB/decade slope, so 4 pole
exhibits a slope of -80 dB/decade.
\end{tcolorbox}

\vspace{1.5em}
\hrule
\vspace{1.5em}

\section*{Question 95: Root Locus Technique (GATE EC 2014-SET-3)}
\addcontentsline{toc}{section}{Question 95: Root Locus Technique (GATE EC 2014-SET-3)}
\noindent \textbf{\color{primary}MCQ} \hfill \textbf{\color{secondary}2 Marks}


\noindent In the root locus plot shown in the figure, the pole/zero marks and the arrows have been removed. Which one of the following transfer functions has this root locus ? 

\begin{center}\includegraphics[max width=0.85\linewidth]{images/img\_57.jpg}\end{center}

\begin{center}\includegraphics[max width=0.85\linewidth]{images/img\_57.jpg}\end{center}


\begin{enumerate}[label=(\Alph*)]
    \item $\frac{s+1}{(s+2)(s+4)(s+7)}$
    \item \textbf{(Correct)} $\frac{s+4}{(s+1)(s+2)(s+7)}$
    \item $\frac{s+7}{(s+1)(s+2)(s+4)}$
    \item $\frac{(s+1)(s+2)}{(s+4)(s+7)}$
\end{enumerate}

\noindent \textbf{Correct Answer:} (B)

\begin{tcolorbox}[solbox, title=Detailed Solution -- Question 95]
\begin{center}\includegraphics[max width=0.85\linewidth]{images/img\_14.jpg}\end{center}

\begin{center}\includegraphics[max width=0.85\linewidth]{images/img\_14.jpg}\end{center}

Since the root locus always emerges from the break-away points i.e. when two poles are at a part of same root locus lie. So $\sigma$ =-1 and -2 will surely be poles.
 Also we know that locus emerges from poles and terminates at zero. So, s=-4 is a zero and s=-7 is a pole 
$\therefore T(s)=\frac{(s+4)}{(s+1)(s+2)(s+7)}$
\end{tcolorbox}

\vspace{1.5em}
\hrule
\vspace{1.5em}

\section*{Question 96: State Space Analysis (GATE EC 2014-SET-3)}
\addcontentsline{toc}{section}{Question 96: State Space Analysis (GATE EC 2014-SET-3)}
\noindent \textbf{\color{primary}MCQ} \hfill \textbf{\color{secondary}2 Marks}


\noindent The state equation of a second-order linear system is given by 

$\dot{x}(t) = Ax(t), \;\; x(0) = x_{0}$

 For  
\[
x_{0}=\begin{bmatrix} 1\\ -1 \end{bmatrix}, \; x(t)=\begin{bmatrix} e^{-t}\\ -e^{-t} \end{bmatrix}
\]
 and  for 
\[
x_{0}=\begin{bmatrix} 0\\ 1 \end{bmatrix},
\]
 
\[
x(t)=\begin{bmatrix} e^{-t}-e^{-2t}\\ -e^{-t}+2e^{-2t} \end{bmatrix}
\]
 

 when  
\[
x_{0}=\begin{bmatrix} 3\\ 5 \end{bmatrix}
\]
, x(t) is


\begin{enumerate}[label=(\Alph*)]
    \item \[
\begin{bmatrix} -8e^{-t} +11e^{-2t}\\ 8e^{-t} -22e^{-2t} \end{bmatrix}
\]
    \item \textbf{(Correct)} \[
\begin{bmatrix} 11e^{-t} -8e^{-2t}\\ -11e^{-t} +16e^{-2t} \end{bmatrix}
\]
    \item \[
\begin{bmatrix} 3e^{-t} -5e^{-2t}\\ -3e^{-t} +10e^{-2t} \end{bmatrix}
\]
    \item \[
\begin{bmatrix} 5e^{-t} -3e^{-2t}\\ -5e^{-t} +6e^{-2t} \end{bmatrix}
\]
\end{enumerate}

\noindent \textbf{Correct Answer:} (B)

\begin{tcolorbox}[solbox, title=Detailed Solution -- Question 96]
Given $\dot{x}(t)=A x(t), \quad x(0)=x_{0}$
 Taking the Laplace transform

\[
\begin{aligned} s X(s)-x(0) &=A X(s) \\ [s I-A] X(s) &=x(0) \\ X(s) &=[s I-A]^{-1} x(0) \\ x(t) &=\mathcal{L}^{-1}[s I-A]^{-1} x(0) &\ldots(i) \end{aligned}
\]

 Condition given are 

\[
\begin{aligned} For x_{0}&=\left[\begin{array}{c}1 \\ -1\end{array}\right], x(t)=\left[\begin{array}{c}e^{-t} \\ -e^{-t}\end{array}\right] \\ For x_{0}&=\left[\begin{array}{l}0 \\ 1\end{array}\right], x(t)=\left[\begin{array}{l}e^{-t}-e^{-2 t} \\ -e^{-t}+2 e^{-2 t}\end{array}\right] \end{aligned}
\]

Using the linearity property in equation (i) 
$K_{1} x_{1}(t)=\mathcal{L}^{-1}[s I-A]^{-1} x_{1}(0) K_{1}$
$K_{2} x_{2}(t)=\mathcal{L}^{-1}[s I-A]^{-1} x_{2}(0) K_{2}$
 Using the linearity property as 

\[
\begin{aligned} K_{1} x_{1}(t)+K_{2} x_{2}(t)&=\mathcal{L}^{-1}[s I-A]^{-1} \\  \left[K_{1} x_{1}(0)+K_{2} x_{2}(0)\right] & \quad \ldots(ii)\\ \text { Also } \quad X_{3}(s)&=[s I-A]^{-1} x_{3}(0) \\ \text{So, } \quad K_{1}\left[\begin{array}{r}1 \\ -1\end{array}\right]+&K_{2}\left[\begin{array}{l}0 \\ 1\end{array}\right]=\left[\begin{array}{l}3 \\ 5\end{array}\right] \\ \left[\begin{array}{l}K_{1}+O K_{2} \\ -K_{1}+K_{2}\end{array}\right]&=\left[\begin{array}{l}3 \\ 5\end{array}\right] \\ \Rightarrow \quad K_{1}&=3 \\ K_{2}&=8 \end{aligned}
\]

 So, from equation (ii), we get x(t)

\[
\begin{aligned} x(t) &=K_{1} x_{1}(t)+K_{2} x_{2}(t) \\ &=3\left[\begin{array}{c} e^{-t} \\ -e^{-t} \end{array}\right]+8\left[\begin{array}{c} e^{-t}-e^{-2 t} \\ -e^{-t}+2 e^{-2 t} \end{array}\right] \\ &=\left[\begin{array}{c} 11 e^{-t}-8 e^{-2 t} \\ -1 t e^{-t}+16 e^{-2 t} \end{array}\right] \end{aligned}
\]
\end{tcolorbox}

\vspace{1.5em}
\hrule
\vspace{1.5em}

\section*{Question 97: Time Response Analysis (GATE EC 2014-SET-3)}
\addcontentsline{toc}{section}{Question 97: Time Response Analysis (GATE EC 2014-SET-3)}
\noindent \textbf{\color{primary}NAT} \hfill \textbf{\color{secondary}2 Marks}


\noindent The steady state error of the system shown in the figure for a unit step input is
\underline{\hspace{1.5cm}}. 

\begin{center}\includegraphics[max width=0.85\linewidth]{images/img\_31.jpg}\end{center}

\begin{center}\includegraphics[max width=0.85\linewidth]{images/img\_31.jpg}\end{center}


\begin{tcolorbox}[solbox, title=Detailed Solution -- Question 97]
After converting the system into equivalent unity feedback system. 

\begin{center}\includegraphics[max width=0.85\linewidth]{images/img\_85.jpg}\end{center}

\begin{center}\includegraphics[max width=0.85\linewidth]{images/img\_85.jpg}\end{center}

$\therefore \quad OLTF =\frac{4(s+4)}{s(s+2)}$
 Type =1 
$e_{s s}$ for step input =0.5
\end{tcolorbox}

\vspace{1.5em}
\hrule
\vspace{1.5em}

\section*{Question 98: Basics, Block Diagrams \& SFGs (GATE EC 2014-SET-3)}
\addcontentsline{toc}{section}{Question 98: Basics, Block Diagrams \& SFGs (GATE EC 2014-SET-3)}
\noindent \textbf{\color{primary}MCQ} \hfill \textbf{\color{secondary}1 Mark}


\noindent Consider the following block diagram in the figure. 

\begin{center}\includegraphics[max width=0.85\linewidth]{images/img\_102.jpg}\end{center}

\begin{center}\includegraphics[max width=0.85\linewidth]{images/img\_102.jpg}\end{center}
 The transfer function $\frac{C(s)}{R(s)}$ is


\begin{enumerate}[label=(\Alph*)]
    \item $\frac{G_{1}G_{2}}{1+G_{1}G_{2}}$
    \item $G_{1}G_{2}+G_{1}+1$
    \item \textbf{(Correct)} $G_{1}G_{2}+G_{2}+1$
    \item $\frac{G_{1}}{1+G_{1}G_{2}}$
\end{enumerate}

\noindent \textbf{Correct Answer:} (C)

\begin{tcolorbox}[solbox, title=Detailed Solution -- Question 98]
\begin{center}\includegraphics[max width=0.85\linewidth]{images/img\_10.jpg}\end{center}

\begin{center}\includegraphics[max width=0.85\linewidth]{images/img\_10.jpg}\end{center}

Converting the block d iagram into signal flow

\begin{center}\includegraphics[max width=0.85\linewidth]{images/img\_118.jpg}\end{center}

\begin{center}\includegraphics[max width=0.85\linewidth]{images/img\_118.jpg}\end{center}

 graph on 
 Forward paths, 
$P_{1}=G_{1} G_{2} ; P_{2}=G_{2} \cdot 1 ; P_{3}=1 \cdot 1=1$
 So, the transfer function is
$\frac{C(s)}{R(s)}=G_{1} G_{2}+G_{2}+1$
\end{tcolorbox}

\vspace{1.5em}
\hrule
\vspace{1.5em}

\section*{Question 99: State Space Analysis (GATE EC 2014-SET-2)}
\addcontentsline{toc}{section}{Question 99: State Space Analysis (GATE EC 2014-SET-2)}
\noindent \textbf{\color{primary}MCQ} \hfill \textbf{\color{secondary}2 Marks}


\noindent Consider the state space system expressed by the signal flow diagram shown in
the figure. 

\begin{center}\includegraphics[max width=0.85\linewidth]{images/img\_65.jpg}\end{center}

\begin{center}\includegraphics[max width=0.85\linewidth]{images/img\_65.jpg}\end{center}

 The corresponding system is


\begin{enumerate}[label=(\Alph*)]
    \item \textbf{(Correct)} always controllable
    \item always observable
    \item always stable
    \item always unstable
\end{enumerate}

\noindent \textbf{Correct Answer:} (A)

\begin{tcolorbox}[solbox, title=Detailed Solution -- Question 99]
The state equation from the given state diagram is

\[
\begin{aligned} \dot{x}_{1} &=x_{2} \\ \dot{x}_{2} &=x_{3} \\ \dot{x}_{3} &=a_{3} x_{3}+a_{2} \cdot x_{2}+a_{1} x_{1}+u \\ \text { Also, } y &=c_{1} x_{1}+c_{2} x_{2}+c_{3} x_{3} \end{aligned}
\]

 Thus, state matrix 

\[
\left[\begin{array}{c}\dot{x}_{1} \\ \dot{x}_{2} \\ \dot{x}_{3}\end{array}\right]=\left[\begin{array}{ccc}0 & 1 & 0 \\ 0 & 0 & 1 \\ a_{1} & a_{2} & a_{3}\end{array}\right]\left[\begin{array}{c}x_{1} \\ x_{2} \\ x_{3}\end{array}\right]+\left[\begin{array}{c}0 \\ 0 \\ 1\end{array}\right] u
\]

 and 
\[
y=\left[\begin{array}{lll}c_{1} & c_{2} & c_{3}\end{array}\right]\left[\begin{array}{l}x_{1} \\ x_{2} \\ x_{3}\end{array}\right]
\]

 Check for controllability:
 The system is said to be controllable if, the rank  of controllability matrix $Q_{c}$ is equal to the rank of the state matrix A . However, if the controllability matrix $Q_{c}$ is a square matrix then the condition for controllability is
$\left|o_{c}\right| \neq 0$
 where 

\[
\begin{aligned} Q_{c}&=\left[B: A B: A^{2} B\right]\\ \therefore Q_{c}&=\left[\begin{array}{ccc}0 & 0 & 1 \\ 0 & 1 & a_{3} \\ 1 & a_{3} & a_{2}+a_{3}^{2}\end{array}\right] \\ \because\left|Q_{C}\right| \neq & 0 \text{ and Rank of } Q_{c}=\text{ Rank of }A=1 \end{aligned}
\]

$\therefore$ The system is controllable.
 Check for observability:
 The system is said to be observable if, the rank of observability matrix $Q_{o}$ is equal to the rank of the state matrix A . However, if the observability matrix $Q_{o}$ is a square matrix then the condition for observability is 
$\left|0_{0}\right| \neq 0$
 where 
 
\[
\begin{aligned} Q_{0}&=\left[C^{T}: A^{T} C^{T}:\left(A^{T}\right)^{2} C^{T}\right] \\ O_{0}&=\left[\begin{array}{ccc}c_{1} & c_{2} & c_{3} \\ s_{3} a_{1} & c_{1}+c_{2} a_{2} & c_{2}+c_{2} a_{3} \\ c_{1}\left(c_{2}+c_{3} a_{3}\right) & c_{3} a_{1}+\left(c_{2}+c_{3} a_{3}\right) a_{2} & c_{1}+c_{2} a_{2}+\left(c_{2}+c_{3} a_{3}\right) a_{3}\end{array}\right] \end{aligned}
\]

$\because\left|Q_{0}\right|$ depends on value of unknown. 
 Hence, not observable always.
\end{tcolorbox}

\vspace{1.5em}
\hrule
\vspace{1.5em}

\section*{Question 100: Frequency Response Analysis (GATE EC 2014-SET-2)}
\addcontentsline{toc}{section}{Question 100: Frequency Response Analysis (GATE EC 2014-SET-2)}
\noindent \textbf{\color{primary}NAT} \hfill \textbf{\color{secondary}2 Marks}


\noindent The Bode asymptotic magnitude plot of a minimum phase system is shown in
the figure 

\begin{center}\includegraphics[max width=0.85\linewidth]{images/img\_83.jpg}\end{center}

\begin{center}\includegraphics[max width=0.85\linewidth]{images/img\_83.jpg}\end{center}

 If the system is connected in a unity negative feedback configuration, the steady state error of the closed loop system, to a unit ramp input, is\underline{\hspace{1.5cm}}.


\begin{tcolorbox}[solbox, title=Detailed Solution -- Question 100]
\begin{center}\includegraphics[max width=0.85\linewidth]{images/img\_86.jpg}\end{center}

\begin{center}\includegraphics[max width=0.85\linewidth]{images/img\_86.jpg}\end{center}

The open loop transfer function of the system 

\[
G(s) H(s)=\frac{K\left(\frac{s}{2}+1\right)}{s\left(\frac{s}{10}+1\right)}\qquad\ldots(i)
\]

 where,
$.26.02|_{\omega=0.1}=20 \log K-20 \log \omega\qquad\ldots(ii)$
 Or, $\quad K=1.99 \simeq 2$ From equation (i) and (ii), we get
$G(s)=\frac{20(s+2)}{2 s(s+10)}\qquad\ldots(iii)$
 Steady state error for ramp input is
$e_{s s}=\frac{A}{K_{V}} \text { where } K_{V}=\lim _{s \rightarrow 0} s G(s)$
 From equation (iii),
$K_{V}=\lim _{s \rightarrow 0} s \times \frac{10(s+2)}{s(s+10)}=2$
$\therefore \quad e_{s s}=\frac{1}{2}=0.5$
\end{tcolorbox}

\vspace{1.5em}
\hrule
\vspace{1.5em}

\section*{Question 101: State Space Analysis (GATE EC 2014-SET-2)}
\addcontentsline{toc}{section}{Question 101: State Space Analysis (GATE EC 2014-SET-2)}
\noindent \textbf{\color{primary}MCQ} \hfill \textbf{\color{secondary}2 Marks}


\noindent An unforced linear time invariant (LTI) system is represented by
 
\[
\begin{vmatrix} \dot{x}_{1}\\ \dot{x}_{2} \end{vmatrix}=\begin{vmatrix} -1 & 0\\ 0&-2 \end{vmatrix}\begin{vmatrix} x_{1}\\ x_{2} \end{vmatrix}
\]

 If the initial conditions are $x_{1}(0)=1 \; and \; x_{2}(0)=-1$, the solution of the state equation is


\begin{enumerate}[label=(\Alph*)]
    \item $x_{1}(t)=-1, x_{2}(t)=2$
    \item $x_{1}(t)=-e^{-t}, x_{2}(t)=2e^{-t}$
    \item \textbf{(Correct)} $x_{1}(t)=e^{-t}, x_{2}(t)=-e^{-2t}$
    \item $x_{1}(t)=-e^{-t}, x_{2}(t)=-2e^{-t}$
\end{enumerate}

\noindent \textbf{Correct Answer:} (C)

\begin{tcolorbox}[solbox, title=Detailed Solution -- Question 101]
\[
\begin{aligned}\left[\begin{array}{l}\dot{x}_{1} \\ \dot{x}_{2}\end{array}\right] &=\left[\begin{array}{cc}-1 & 0 \\ 0 & -2\end{array}\right]\left[\begin{array}{l}x_{1} \\ x_{2}\end{array}\right] \\ \dot{x}_{1} &=-x_{1}+0 x_{2} \\ \dot{x}_{2} &=0 x_{1}-2 x_{2} \end{aligned}
\]

 Applying Laplace transform on both equations 

\[
\begin{aligned} \Rightarrow s X_{1}(s)-x_{1}(0)&=-X_{1}(s) \\ X_{1}(s)[s+1]&=X_{1}(0) \quad &\text{given }x_{1}(0)=1 \\ X_{1}(s)&=\frac{1}{s+1} \\ \Rightarrow \quad e^{-t}&=x_{1}(t)\\ \Rightarrow \quad s X_{2}(s)-X_{2}(0)&=-2 X_{2}(s) \\ (s+2) X_{2}(s)&=X_{2}(0) &X_{2}(0)=-1 \\ X_{2}(s)&=-\frac{1}{s+2} \\ x_{2}(t)&=-e^{-2 t} \end{aligned}
\]
\end{tcolorbox}

\vspace{1.5em}
\hrule
\vspace{1.5em}

\section*{Question 102: Basics, Block Diagrams \& SFGs (GATE EC 2014-SET-2)}
\addcontentsline{toc}{section}{Question 102: Basics, Block Diagrams \& SFGs (GATE EC 2014-SET-2)}
\noindent \textbf{\color{primary}MCQ} \hfill \textbf{\color{secondary}1 Mark}


\noindent For the following system, 

\begin{center}\includegraphics[max width=0.85\linewidth]{images/img\_8.jpg}\end{center}

\begin{center}\includegraphics[max width=0.85\linewidth]{images/img\_8.jpg}\end{center}

 when $X_{1}(s) = 0$, the transfer function $\frac{Y(s)}{X_{2}(s)}$  is


\begin{enumerate}[label=(\Alph*)]
    \item $\frac{s+1}{s^{2}}$
    \item $\frac{1}{s+1}$
    \item $\frac{s+2}{s(s+1)}$
    \item \textbf{(Correct)} $\frac{s+1}{s(s+2)}$
\end{enumerate}

\noindent \textbf{Correct Answer:} (D)

\begin{tcolorbox}[solbox, title=Detailed Solution -- Question 102]
Redrawing the block diagram with $X_{1}(s)=0$

\begin{center}\includegraphics[max width=0.85\linewidth]{images/img\_58.jpg}\end{center}

\begin{center}\includegraphics[max width=0.85\linewidth]{images/img\_58.jpg}\end{center}

 The transfer function
$T(s)=\frac{Y(s)}{X_{2}(s)}=\frac{G(s)}{1+G(s) H(s)} \quad\ldots(i)$
 Here, $G(s)=\frac{1}{s}$ and $H(s)=\frac{s}{s+1}$

\[
\frac{Y(s)}{X_{2}(s)}=\frac{1 / s}{1+\frac{1}{s} \times \frac{s}{s+1}}=\frac{(s+1)}{s(s+2)}
\]
\end{tcolorbox}

\vspace{1.5em}
\hrule
\vspace{1.5em}

\section*{Question 103: Time Response Analysis (GATE EC 2014-SET-2)}
\addcontentsline{toc}{section}{Question 103: Time Response Analysis (GATE EC 2014-SET-2)}
\noindent \textbf{\color{primary}NAT} \hfill \textbf{\color{secondary}1 Mark}


\noindent The natural frequency of an undamped second-order system is 40 rad/s. If the
system is damped with a damping ratio 0.3, the damped natural frequency in
rad/s is \underline{\hspace{1.5cm}}.


\begin{tcolorbox}[solbox, title=Detailed Solution -- Question 103]
\[
\begin{aligned} \text { Given, } & \omega_{n}=40 \mathrm{rad} / \mathrm{sec} \\ & \xi=0.3 \\ \therefore \quad & \omega_{d}=\omega_{n} \sqrt{1-\xi^{2}}=38.16 \mathrm{rad} / \mathrm{sec} \end{aligned}
\]
\end{tcolorbox}

\vspace{1.5em}
\hrule
\vspace{1.5em}

\section*{Question 104: Compensators \& Controllers (GATE EC 2014-SET-1)}
\addcontentsline{toc}{section}{Question 104: Compensators \& Controllers (GATE EC 2014-SET-1)}
\noindent \textbf{\color{primary}MCQ} \hfill \textbf{\color{secondary}2 Marks}


\noindent For the following feedback system $G(s)=\frac{1}{(s+1)(s+2)}$. The 2% settling time of the step response is required to be less than 2 seconds.
Which one of the following compensators C(s) achieves this ? 

\begin{center}\includegraphics[max width=0.85\linewidth]{images/img\_120.jpg}\end{center}

\begin{center}\includegraphics[max width=0.85\linewidth]{images/img\_120.jpg}\end{center}


\begin{enumerate}[label=(\Alph*)]
    \item $3 (\frac{1}{s+5})$
    \item $5 (\frac{0.03}{s}+1)$
    \item \textbf{(Correct)} $2 (s+4)$
    \item $4 (\frac{s+8}{s+3})$
\end{enumerate}

\noindent \textbf{Correct Answer:} (C)

\begin{tcolorbox}[solbox, title=Detailed Solution -- Question 104]
Given open loop transfer function is

\[
\begin{aligned} G(s)&=\frac{1}{(s+1)(s+2)} \\ \therefore T(s)&=\frac{G(s)}{1+G(s)}=\frac{1}{(s+1)(s+2)+1} \\ T(s)&=\frac{1}{s^{2}+3 s+3} &\ldots(i) \end{aligned}
\]

 Comparing equation (i) with standard transfer function
$\xi \omega_{n}=\frac{3}{2}=1.5$
$\therefore$ 2% settling time
$\tau_{s}=\frac{4}{\xi \omega_{n}}=\frac{4}{1.5}=2.67>2 \mathrm{sec}$
 Thus, in order to make settling time $\left(\tau_{s}\right)$ less than 2 sec, PD controller should be used. Hence, option (C) is the correct choice. 
 Where, $C(s)=2(s+4)$
$\therefore$ New transfer function $=T^{\prime}(G)$
$=\frac{C(s) G(s)}{1+C(s) G(s)}=\frac{2(s+4)}{s^{2}+3 s+2+2 s+8}$
 or $T^{\prime}(s)=\frac{2(s+4)}{s^{2}+5 s+10}$

\[
\therefore \quad \tau_{s}=\frac{4}{\xi \omega_{n}}=\frac{4}{5 / 2}=\frac{4}{2.5}=1.6
\]
\end{tcolorbox}

\vspace{1.5em}
\hrule
\vspace{1.5em}

\section*{Question 105: Frequency Response Analysis (GATE EC 2014-SET-1)}
\addcontentsline{toc}{section}{Question 105: Frequency Response Analysis (GATE EC 2014-SET-1)}
\noindent \textbf{\color{primary}NAT} \hfill \textbf{\color{secondary}2 Marks}


\noindent The phase margin in degrees of $G(s)=\frac{1}{(s+0.1)(s+1)(s+10)}$
calculated using the asymptotic Bode plot is \underline{\hspace{1.5cm}}.


\begin{tcolorbox}[solbox, title=Detailed Solution -- Question 105]
\[
\begin{array}{l} \frac{10}{0.7 \times 10\left(\frac{s}{0.1}+1\right)\left(\frac{s}{1}+1\right)\left(\frac{s}{10}+1\right)} \\ =\frac{10}{\left(\frac{s}{0.1}+1\right)\left(\frac{s}{1}+1\right)\left(\frac{s}{10}+1\right)} \\ \text { Amount of shitt }=2010 \mathrm{g} \mathrm{K}=20 \mathrm{log} 10=20 \mathrm{dB} \end{array}
\]

\begin{center}\includegraphics[max width=0.85\linewidth]{images/img\_37.jpg}\end{center}

\begin{center}\includegraphics[max width=0.85\linewidth]{images/img\_37.jpg}\end{center}

\[
\begin{aligned} -20dB&=\frac{20-0}{\log0.1-\log\omega_{gc}}\\ \text{or},\omega_{gc}&=1\mathrm{rad}/\mathrm{sec}&\ldots(i)\\ \mathrm{PM}&=180^{\circ}+\angle GH(j\omega|_{at\omega=\omega_{wc}}&\ldots(ii) \end{aligned}
\]

 From the given transfer function,

\[
\begin{aligned} \angle G(j \omega)&=-\tan ^{-1}\left(\frac{\omega_{g c}}{0.1}\right)-\tan ^{-1}\left(\frac{\omega_{g c}}{1}\right) \\ &-\tan ^{-1}\left(\frac{\omega_{g c}}{10}\right) \\ \alpha \angle G(j(1)) &=-84.289^{\circ}-45^{\circ}-5.711^{\circ} \\ &=-135^{\circ} &\ldots(iii) \end{aligned}
\]

$\therefore$ From equation(ii) and (iii) we get,
$\text { Phase margin } \mathrm{PM}=180^{\circ}-135^{\circ}=45^{\circ}$
\end{tcolorbox}

\vspace{1.5em}
\hrule
\vspace{1.5em}

\section*{Question 106: State Space Analysis (GATE EC 2014-SET-1)}
\addcontentsline{toc}{section}{Question 106: State Space Analysis (GATE EC 2014-SET-1)}
\noindent \textbf{\color{primary}MCQ} \hfill \textbf{\color{secondary}2 Marks}


\noindent Consider the state space model of a system, as given below 
\[
\begin{bmatrix} \dot{x}_{1}\\ \dot{x}_{2}\\ \dot{x}_{3} \end{bmatrix}=\begin{vmatrix} -1 & 1&0 \\ 0 & -1 & 0\\ 0&0 & -2 \end{vmatrix}\begin{bmatrix} x_{1}\\ x_{2}\\ x_{3} \end{bmatrix}+\begin{bmatrix} 0\\ 4\\ 0 \end{bmatrix}u;y\begin{bmatrix} 1 & 1 &1 \end{bmatrix}\begin{bmatrix} x_{1}\\ x_{2}\\ x_{3} \end{bmatrix}
\]
 The system is


\begin{enumerate}[label=(\Alph*)]
    \item controllable and observable
    \item \textbf{(Correct)} uncontrollable and observable
    \item uncontrollable and unobservable
    \item controllable and unobservable
\end{enumerate}

\noindent \textbf{Correct Answer:} (B)

\begin{tcolorbox}[solbox, title=Detailed Solution -- Question 106]
\[
\begin{aligned} \left[\begin{array}{c}\dot{x}_{1} \\ \dot{x}_{2} \\ \dot{x}_{3}\end{array}\right]&=\left[\begin{array}{rrr}-1 & 1 & 0 \\ 0 & -1 & 0 \\ 0 & 0 & -2\end{array}\right]\left[\begin{array}{l}x_{1} \\ x_{2} \\ x_{3}\end{array}\right]+\left[\begin{array}{l}0 \\ 4 \\ 0\end{array}\right] u \\ y&=[1 \;1 \;1 ]\left[\begin{array}{r}x_{3} \\ x_{3} \\ x_{3}\end{array}\right] \\ A&=\left[\begin{array}{rrr}-1 & 1 & 0 \\ 0 & -1 & 0 \\ 0 & 0 & -2\end{array}\right], B=\left[\begin{array}{l}x_{1} \\ x_{2} \\ 0\end{array}\right], C=\left[1\;1\;1\right] \\ \end{aligned}
\]

 Check for controllability:

\[
\begin{aligned} Q_{C} &=\left[B: A B: A^{2} B\right] \\ &=\left[\begin{array}{rrr} 0 & 4 & -8 \\ 4 & -4 & 4 \\ 0 & 0 & 0 \end{array}\right] \end{aligned}
\]

For controllable.
$\left|Q_{c}\right| \neq 0$
 Here , $\quad\left|Q_{c}\right|=4(0)=0 \quad \therefore$ Uncontrollable
 Check for observability:

\[
\begin{aligned} Q_{0} &=\left[C^{T}: A^{T} C^{T}: A^{2 T} C^{T}\right] \\ &=\left[\begin{array}{rrr} 1 & -1 & 1 \\ 1 & 0 & -1 \\ 1 & -2 & 4 \end{array}\right] \end{aligned}
\]

 For observable,
$\left|Q_{0}\right| \neq 0$
 Here, $\left|Q_{o}\right|=1 \quad \therefore$ Observable
\end{tcolorbox}

\vspace{1.5em}
\hrule
\vspace{1.5em}

\section*{Question 107: Frequency Response Analysis (GATE EC 2014-SET-1)}
\addcontentsline{toc}{section}{Question 107: Frequency Response Analysis (GATE EC 2014-SET-1)}
\noindent \textbf{\color{primary}MCQ} \hfill \textbf{\color{secondary}1 Mark}


\noindent Consider the feedback system shown in the figure. The Nyquist plot of G(s)  is
also shown. Which one of the following conclusions is correct ? 

\begin{center}\includegraphics[max width=0.85\linewidth]{images/img\_94.jpg}\end{center}

\begin{center}\includegraphics[max width=0.85\linewidth]{images/img\_94.jpg}\end{center}


\begin{enumerate}[label=(\Alph*)]
    \item G(s) h is an all pass filter
    \item G(s) h is a strictly proper transfer function
    \item G(s) h is a stable and minimum phase transfer function
    \item \textbf{(Correct)} The closed-loop system is unstable for sufficiently large and positive k
\end{enumerate}

\noindent \textbf{Correct Answer:} (D)

\begin{tcolorbox}[solbox, title=Detailed Solution -- Question 107]
\begin{center}\includegraphics[max width=0.85\linewidth]{images/img\_53.jpg}\end{center}

\begin{center}\includegraphics[max width=0.85\linewidth]{images/img\_53.jpg}\end{center}

From the Nyquist plot, it can be seen that it does
not encircle the critical point (-1,0) 
$N=P \cdot Z, \quad N=0$

The closed loop-system is unstable for
sufficiently large and positive K
\end{tcolorbox}

\vspace{1.5em}
\hrule
\vspace{1.5em}

\section*{Question 108: Stability Analysis \& Routh-Hurwitz (GATE EC 2014-SET-1)}
\addcontentsline{toc}{section}{Question 108: Stability Analysis \& Routh-Hurwitz (GATE EC 2014-SET-1)}
\noindent \textbf{\color{primary}NAT} \hfill \textbf{\color{secondary}1 Mark}


\noindent The forward path transfer function of a unity negative feedback system is given
by 
 $G(s)=\frac{K}{(s+2)(s-1)}$ 
 The value of K which will place both the poles of the closed-loop system at the same location, is \underline{\hspace{1.5cm}}.


\begin{tcolorbox}[solbox, title=Detailed Solution -- Question 108]
Given that, $\quad G(s)=\frac{K}{(s+2)(s-1)}$
 Using root locus method, the break point can be obtain as

\begin{center}\includegraphics[max width=0.85\linewidth]{images/img\_4.jpg}\end{center}

\begin{center}\includegraphics[max width=0.85\linewidth]{images/img\_4.jpg}\end{center}

\[
\begin{aligned} \Rightarrow \quad 1+G(s)&=0 \\ 1+\frac{K}{(s+2)(s-1)} &=0 \\ \text{or }\quad K&=-(s+2)(s-1)\\ \frac{d K}{d s} &=-2 s-1=0 \\ \text{or }\quad s &=-0.5 \end{aligned}
\]

 To have, both the poles at the same directions

\[
\begin{aligned} |G(s)|_{s=0.5} &=1 \\ K &=2.25 \end{aligned}
\]
\end{tcolorbox}

\vspace{1.5em}
\hrule
\vspace{1.5em}

\section*{Question 109: State Space Analysis (GATE EC 2013)}
\addcontentsline{toc}{section}{Question 109: State Space Analysis (GATE EC 2013)}
\noindent \textbf{\color{primary}MCQ} \hfill \textbf{\color{secondary}1 Mark}


\noindent The state diagram of a system is shown below. A system is described by the state-variable equations $\dot{X}= AX+Bu$ ; $y=CX+Du$

\begin{center}\includegraphics[max width=0.85\linewidth]{images/img\_7.jpg}\end{center}

\begin{center}\includegraphics[max width=0.85\linewidth]{images/img\_7.jpg}\end{center}
 
The state transition matrix $e^{At}$ of the system shown in the figure above is


\begin{enumerate}[label=(\Alph*)]
    \item \textbf{(Correct)} \[
\begin{bmatrix} e^{-t} &0 \\ te^{-t} & e^{-t} \end{bmatrix}
\]
    \item \[
\begin{bmatrix} e^{-t} &0 \\ -te^{-t} & e^{-t} \end{bmatrix}
\]
    \item \[
\begin{bmatrix} e^{-t} &0 \\ e^{-t} & e^{-t} \end{bmatrix}
\]
    \item \[
\begin{bmatrix} e^{-t} &-te^{-t} \\ 0 & e^{-t} \end{bmatrix}
\]
\end{enumerate}

\noindent \textbf{Correct Answer:} (A)

\begin{tcolorbox}[solbox, title=Detailed Solution -- Question 109]
\[
\begin{aligned} A &=\left[\begin{array}{rr} -1 & 0 \\ 1 & -1 \end{array}\right] \\ s I-A&=\left[\begin{array}{cc} s+1 & 0 \\ -1 & s+1 \end{array}\right] \\ [s I-A]^{-1}&=\frac{1}{(s+1) \times(s+1)}\left[\begin{array}{cc} s+1 & 0 \\ 1 & s+1 \end{array}\right]\\ [s I-A]^{-1}&=\left[\begin{array}{cc}\frac{1}{s+1} & 0 \\ \frac{1}{(s+1)^{2}} & \frac{1}{s+1}\end{array}\right] \\ \phi(t)&=e^{A t}=L^{-1}{(s I-A)^{-1}} \\ e^{A t}&=\left[\begin{array}{cc}e^{-t} & 0 \\ t e^{-t} & e^{-t}\end{array}\right] \end{aligned}
\]
\end{tcolorbox}

\vspace{1.5em}
\hrule
\vspace{1.5em}

\section*{Question 110: State Space Analysis (GATE EC 2013)}
\addcontentsline{toc}{section}{Question 110: State Space Analysis (GATE EC 2013)}
\noindent \textbf{\color{primary}MCQ} \hfill \textbf{\color{secondary}1 Mark}


\noindent The state diagram of a system is shown below. A system is described by the state-variable equations
$\dot{X}= AX+Bu$;  $y=CX+Du$

\begin{center}\includegraphics[max width=0.85\linewidth]{images/img\_7.jpg}\end{center}

\begin{center}\includegraphics[max width=0.85\linewidth]{images/img\_7.jpg}\end{center}
 
The state-variable equations of the system shown in the figure above are


\begin{enumerate}[label=(\Alph*)]
    \item \textbf{(Correct)} \[
\dot{X}= \begin{bmatrix} -1 &0 \\ 1 & -1 \end{bmatrix}X+\begin{bmatrix} -1\\ 1 \end{bmatrix}u
\]

\[
y=\begin{bmatrix} 1 & -1 \end{bmatrix}X+u
\]
    \item \[
\dot{X}= \begin{bmatrix} -1 &0 \\ -1 & -1 \end{bmatrix}X+\begin{bmatrix} -1\\ 1 \end{bmatrix}u
\]
 
\[
y=\begin{bmatrix} -1 & -1 \end{bmatrix}X+u
\]
    \item \[
\dot{X}= \begin{bmatrix} -1 &0 \\ -1 & -1 \end{bmatrix}X+\begin{bmatrix} -1\\ 1 \end{bmatrix}u
\]
 
\[
y=\begin{bmatrix} -1 & -1 \end{bmatrix}X-u
\]
    \item \[
\dot{X}= \begin{bmatrix} -1 &-1 \\ 0 & -1 \end{bmatrix}X+\begin{bmatrix} -1\\ 1 \end{bmatrix}u
\]
 
\[
y=\begin{bmatrix} 1 & -1 \end{bmatrix}X-u
\]
\end{enumerate}

\noindent \textbf{Correct Answer:} (A)

\begin{tcolorbox}[solbox, title=Detailed Solution -- Question 110]
\begin{center}\includegraphics[max width=0.85\linewidth]{images/img\_89.jpg}\end{center}

\begin{center}\includegraphics[max width=0.85\linewidth]{images/img\_89.jpg}\end{center}

\[
\begin{aligned} \mathrm{So}, \dot{x}_{1} &=-x_{1}-u \\ \dot{x}_{2} &=-\left(x_{2}+\dot{x}_{1}\right)=-\left(x_{2}-x_{1}-u\right) \\ \dot{x}_{2} &=x_{1}-x_{2}+u \\ y &=\dot{x}_{2} \\ y &=x_{1}-x_{2}+u \\ \dot{x}_{2} &=\left[\begin{array}{rr}-1 & 0 \\1 & -1\end{array}\right]\left[\begin{array}{c} x_{1} \\x_{2}\end{array}\right]+\left[\begin{array}{c}-1 \\1 \end{array}\right] u \\ \dot{x} &=\left[\begin{array}{rr} -1 & 0 \\ 1 & -1 \end{array}\right] X+\left[\begin{array}{c} -1 \\ 1 \end{array}\right] u \end{aligned}
\]
\end{tcolorbox}

\vspace{1.5em}
\hrule
\vspace{1.5em}

\section*{Question 111: Basics, Block Diagrams \& SFGs (GATE EC 2013)}
\addcontentsline{toc}{section}{Question 111: Basics, Block Diagrams \& SFGs (GATE EC 2013)}
\noindent \textbf{\color{primary}MCQ} \hfill \textbf{\color{secondary}1 Mark}


\noindent The signal flow graph for a system is given below. The transfer function $\frac{Y\left ( s \right )}{U\left ( s \right )}$ for this system is

\begin{center}\includegraphics[max width=0.85\linewidth]{images/img\_13.jpg}\end{center}

\begin{center}\includegraphics[max width=0.85\linewidth]{images/img\_13.jpg}\end{center}


\begin{enumerate}[label=(\Alph*)]
    \item \textbf{(Correct)} $\frac{s+1}{5s^{2}+6s+2}$
    \item $\frac{s+1}{s^{2}+6s+2}$
    \item $\frac{s+1}{s^{2}+4s+2}$
    \item $\frac{1}{5s^{2}+6s+2}$
\end{enumerate}

\noindent \textbf{Correct Answer:} (A)

\begin{tcolorbox}[solbox, title=Detailed Solution -- Question 111]
Using Mason's gain formula,

\[
\begin{aligned} \Delta &=1-\left[-2 s^{-1}-2 s^{-2}-4-4 s^{-1}\right] \\ \Delta &=1+\frac{2}{s}+\frac{2}{s^{2}}+4+\frac{4}{s} \\ \Delta &=\frac{5 s^{2}+6 s+2}{s^{2}} \\ P_{1} &=s^{2}=\frac{1}{s^{2}} \\ P_{2} &=s^{-1}=\frac{1}{s} \\ \Delta_{1} &=1 \\ \Delta_{2} &=1 \\ \frac{Y(s)}{U(s)} &=\frac{\sum P_{k} \Delta_{k}}{\Delta}=\frac{\frac{1}{s^{2}} \times 1+\frac{1}{s} \times 1}{\frac{5 s^{2}+6 s+2}{s^{2}}} \\ \frac{Y(s)}{U(s)} &=\frac{s+1}{5 s^{2}+6 s+2} \end{aligned}
\]
\end{tcolorbox}

\vspace{1.5em}
\hrule
\vspace{1.5em}

\section*{Question 112: Time Response Analysis (GATE EC 2013)}
\addcontentsline{toc}{section}{Question 112: Time Response Analysis (GATE EC 2013)}
\noindent \textbf{\color{primary}MCQ} \hfill \textbf{\color{secondary}1 Mark}


\noindent A system described by a linear, constant coefficient, ordinary, first order differential equation has an exact solution given by y(t) for t > 0, when the forcing function is x(t) and the initial condition is y(0). If one wishes to modify the system so that the solution becomes -2y(t) for t > 0, we need to


\begin{enumerate}[label=(\Alph*)]
    \item change the initial condition to -y(0) and the forcing function to 2x(t)
    \item change the initial condition to 2y(0) and the forcing function to -x(t)
    \item change the initial condition to $j\sqrt{2}y\left ( 0 \right )$ and the forcing function to $j\sqrt{2}x\left ( t \right )$
    \item \textbf{(Correct)} change the initial condition to -2y(0) and the forcing function to -2x(t)
\end{enumerate}

\noindent \textbf{Correct Answer:} (D)

\begin{tcolorbox}[solbox, title=Detailed Solution -- Question 112]
$\frac{d y(t)}{d t}+k y(t)=x(t)$
 Taking Laplace transform of both sides, we have 
$s Y(s)-y(0)+k Y(s)=X(s)$
$Y(s)[s+k]=X(s)+y(0)$
$\Rightarrow Y(s)=\frac{X(s)}{s+k}+\frac{y(0)}{s+k}$
 Taking inverse Laplace transform, we have 
$y(t)=e^{-k t} x(t)+y(0) e^{-k t}$
 So if we want -2y(t) as a solution both x(t) and  y(0) has to be multiplied by -2; hence change x(t) by -2x(t) and y(0) by -2y(0).
\end{tcolorbox}

\vspace{1.5em}
\hrule
\vspace{1.5em}

\section*{Question 113: Time Response Analysis (GATE EC 2013)}
\addcontentsline{toc}{section}{Question 113: Time Response Analysis (GATE EC 2013)}
\noindent \textbf{\color{primary}MCQ} \hfill \textbf{\color{secondary}1 Mark}


\noindent The open-loop transfer function of a dc motor is given as $\frac{w\left ( s \right )}{V_{a}\left ( s \right )}= \frac{10}{1+10s}$. When connected in feedback as shown below, the approximate value of $K_{a}$ that will reduce the time constant of the closed loop system by one hundred times as compared to that of the open-loop system is

\begin{center}\includegraphics[max width=0.85\linewidth]{images/img\_15.jpg}\end{center}

\begin{center}\includegraphics[max width=0.85\linewidth]{images/img\_15.jpg}\end{center}


\begin{enumerate}[label=(\Alph*)]
    \item $1$
    \item $5$
    \item \textbf{(Correct)} $10$
    \item $100$
\end{enumerate}

\noindent \textbf{Correct Answer:} (C)

\begin{tcolorbox}[solbox, title=Detailed Solution -- Question 113]
\[
\begin{aligned} \frac{\omega(s)}{V_{a}(s)} &=\frac{10}{1+10 s} \\ 1+T s &=1+10 s \\ T s &=10 s \\ \therefore \quad T &=10\\ T^{\prime}=& \frac{T}{100}=\frac{10}{100}=\frac{1}{10} \\ \frac{\omega(s)}{R(s)}=& \frac{K_{a} \times \frac{10}{1+10 s}}{1+K_{a} \frac{10}{1+10 s}} \\ \frac{\omega(s)}{R(s)}&=\frac{10 K_{a}}{1+10 s+10 K_{a}}=\frac{10 K_{a}}{\left(1+10 K_{a}\right)+10s} \\ \frac{\omega(s)}{R(s)}&=\frac{10k_{a}}{\left(1+10 K_{a}\right)\left[1+\frac{10 s}{1+10 K_{a}}\right]} \\ T^{\prime}&=\frac{10}{1+10K_{a}}=\frac{1}{10}\\ 1+10 K_{a}&=100 \\ 10 K_{a}&=99 \\ K_{a}&=9.9 \\ K_{a} &\cong 10 \end{aligned}
\]
\end{tcolorbox}

\vspace{1.5em}
\hrule
\vspace{1.5em}

\section*{Question 114: Frequency Response Analysis (GATE EC 2013)}
\addcontentsline{toc}{section}{Question 114: Frequency Response Analysis (GATE EC 2013)}
\noindent \textbf{\color{primary}MCQ} \hfill \textbf{\color{secondary}1 Mark}


\noindent The Bode plot of a transfer function G(s) is shown in the figure below.

\begin{center}\includegraphics[max width=0.85\linewidth]{images/img\_126.jpg}\end{center}

\begin{center}\includegraphics[max width=0.85\linewidth]{images/img\_126.jpg}\end{center}
 The gain $\left ( 20\log\left | G\left ( s \right ) \right | \right )$ is 32 dB and -8 dB at 1 rad/s and 10 rad/s respectively. The phase is negative for all $\omega$. Then G(s) is


\begin{enumerate}[label=(\Alph*)]
    \item $\frac{39.8}{s}$
    \item \textbf{(Correct)} $\frac{39.8}{s^{2}}$
    \item $\frac{32}{s}$
    \item $\frac{32}{s^{2}}$
\end{enumerate}

\noindent \textbf{Correct Answer:} (B)

\begin{tcolorbox}[solbox, title=Detailed Solution -- Question 114]
$10 \mathrm{rad} / \mathrm{s}$ to $1 \mathrm{rad} / \mathrm{s}$ is 1 decade
$32-(-8)=40 d B$
 So, the slope is 40dB/decade it means there are two poles at origin, it means either option (B) or  option (D) is correct put $\omega=1$rad/sec in both the options. 
$20 \log \left[\frac{39.8}{(1)^{2}}\right]=32 \mathrm{dB}$
$20 \log \left[\frac{32}{(1)^{2}}\right]=30.1 \mathrm{dB}$
 So, option (B) is correct option $=\frac{39.8}{s^{2}}$
\end{tcolorbox}

\vspace{1.5em}
\hrule
\vspace{1.5em}

\section*{Question 115: Compensators \& Controllers (GATE EC 2012)}
\addcontentsline{toc}{section}{Question 115: Compensators \& Controllers (GATE EC 2012)}
\noindent \textbf{\color{primary}MCQ} \hfill \textbf{\color{secondary}1 Mark}


\noindent The transfer function of a compensator is given as 
$G_{c}(s)=\frac{s+a}{s+b}$

The phase of the above lead compensator is maximum at


\begin{enumerate}[label=(\Alph*)]
    \item \textbf{(Correct)} $\sqrt{2}$ rad/s
    \item $\sqrt{3}$rad/s
    \item $\sqrt{6}$ rad/s
    \item 1/$\sqrt{3}$ rad/s
\end{enumerate}

\noindent \textbf{Correct Answer:} (A)

\begin{tcolorbox}[solbox, title=Detailed Solution -- Question 115]
For phase to be maximum,

\[
\begin{aligned} \frac{\partial P}{\partial \omega}&=0 \\ \frac{1}{1+\frac{\omega^{2}}{a^{2}}}- \frac{1/b}{1+\frac{\omega^{2}}{b^{2}}}&=0 \\ \frac{1}{1+\omega^{2}}- \frac{1/b}{1+\frac{\omega^{2}}{4}}&=0\\ \frac{1}{1+\omega^{2}}-\frac{2}{\omega^{2}+4} &=0 \\ \omega^{2}+4-2-2 \omega^{2} &=0 \\ \omega^{2} &=2 \\ \omega &=\sqrt{2} \mathrm{rad} / \mathrm{sec} \end{aligned}
\]
\end{tcolorbox}

\vspace{1.5em}
\hrule
\vspace{1.5em}

\section*{Question 116: Compensators \& Controllers (GATE EC 2012)}
\addcontentsline{toc}{section}{Question 116: Compensators \& Controllers (GATE EC 2012)}
\noindent \textbf{\color{primary}MCQ} \hfill \textbf{\color{secondary}1 Mark}


\noindent The transfer function of a compensator is given as 
$G_{c}(s)=\frac{s+a}{s+b}$

$G_{c}(S)$ is a lead compensator if


\begin{enumerate}[label=(\Alph*)]
    \item \textbf{(Correct)} a =1, b = 2
    \item a = 3, b = 2
    \item a = -3, b = -1
    \item a = 3, b = 1
\end{enumerate}

\noindent \textbf{Correct Answer:} (A)

\begin{tcolorbox}[solbox, title=Detailed Solution -- Question 116]
$G_{c}(s)=\left(\frac{s+a}{s+b}\right)$
 Phase $P=\tan ^{-1}(\omega / a)-\tan ^{-1} \omega / b$
 for lead comparators phase must be + ve.
 for this $\frac{\omega}{a} > \frac{\omega}{b} \Rightarrow a < b$
 So $a=1, b=2$
\end{tcolorbox}

\vspace{1.5em}
\hrule
\vspace{1.5em}

\section*{Question 117: State Space Analysis (GATE EC 2012)}
\addcontentsline{toc}{section}{Question 117: State Space Analysis (GATE EC 2012)}
\noindent \textbf{\color{primary}MCQ} \hfill \textbf{\color{secondary}1 Mark}


\noindent The state variable description of an LTI system is given by 
\[
\begin{pmatrix} \dot{x}_{1}\\ \dot{x}_{2}\\ \dot{x}_{3} \end{pmatrix}=\begin{pmatrix} 0 & a_{1}& 0\\ 0& 0 & a_{2}\\ a_{3} & 0 & 0 \end{pmatrix}\begin{pmatrix} x_{1}\\ x_{2}\\ x_{3} \end{pmatrix}+\begin{pmatrix} 0\\ 0\\ 1 \end{pmatrix}u
\]
 

\[
y=\begin{pmatrix} 1 &0 &0 \end{pmatrix}\begin{pmatrix} x_{1}\\ x_{2}\\ x_{3} \end{pmatrix}
\]
 
 where y is the output and u is the input. The system is controllable for


\begin{enumerate}[label=(\Alph*)]
    \item $a_{1}\neq 0, a_{2}=0,a_{3} \neq 0$
    \item $a_{1}= 0, a_{2}\neq 0,a_{3} \neq 0$
    \item $a_{1}= 0, a_{2}\neq 0,a_{3} = 0$
    \item \textbf{(Correct)} $a_{1}\neq 0, a_{2}\neq 0,a_{3} = 0$
\end{enumerate}

\noindent \textbf{Correct Answer:} (D)

\begin{tcolorbox}[solbox, title=Detailed Solution -- Question 117]
\[
\begin{aligned} \left[\begin{array}{l}\dot{x}_{1} \\\dot{x}_{2} \\\dot{x}_{3}\end{array}\right]&=\left[\begin{array}{lll} 0 & a_{1} & 0 \\0 & 0 & a_{2} \\a_{3} & 0 & 0\end{array}\right]\left[\begin{array}{l} x_{1} \\x_{2} \\x_{3}\end{array}\right]+\left[\begin{array}{l}0 \\0 \\1 \end{array}\right] u\\ y&=\left[\begin{array}{lll}1 & 0 & 0\end{array}\right]\left[\begin{array}{l} x_{1} \\x_{2} \\x_{3}\end{array}\right]\\ Q_{c}&=\left[B A B A^{2} B\right]\\ A B&=\left[\begin{array}{lll} 0 & a_{1} & 0 \\ 0 & 0 & a_{2} \\ a_{3} & 0 & 0 \end{array}\right]\left[\begin{array}{l} 0 \\ 0 \\ 1 \end{array}\right]=\left[\begin{array}{l} 0 \\ a_{2} \\ 0 \end{array}\right]\\ A^{2} B&=A A B=\left[\begin{array}{lll} 0 & a_{1} & 0 \\ 0 & 0 & a_{2} \\ a_{3} & 0 & 0 \end{array}\right]\left[\begin{array}{l} 0 \\ a_{2} \\ 0 \end{array}\right]\\ A^{2} B&=\left[\begin{array}{c} a_{1} a_{2} \\ 0 \\ 0 \end{array}\right]\\ Q_{c}&=\left[\begin{array}{ccc}0 & 0 & a_{1} a_{2} \\ 0 & a_{2} & 0 \\ 1 & 0 & 0\end{array}\right] \end{aligned}
\]

For system to be controllable,

\[
\begin{array}{l} |Q_{i}| \neq 0 \\ \Rightarrow\left(0-a, a_{i}\right) \neq 0 \\ \Rightarrow a_{1} \neq 0 \\ a_{2} \neq 0 \end{array}
\]
\end{tcolorbox}

\vspace{1.5em}
\hrule
\vspace{1.5em}

\section*{Question 118: Stability Analysis \& Routh-Hurwitz (GATE EC 2012)}
\addcontentsline{toc}{section}{Question 118: Stability Analysis \& Routh-Hurwitz (GATE EC 2012)}
\noindent \textbf{\color{primary}MCQ} \hfill \textbf{\color{secondary}1 Mark}


\noindent The feedback system shown below oscillates at 2 rad/s when 

\begin{center}\includegraphics[max width=0.85\linewidth]{images/img\_92.jpg}\end{center}

\begin{center}\includegraphics[max width=0.85\linewidth]{images/img\_92.jpg}\end{center}


\begin{enumerate}[label=(\Alph*)]
    \item \textbf{(Correct)} K=2 and a=0.75
    \item K=3 and a=0.75
    \item K=4 and a=0.5
    \item K=2 and a=0.5
\end{enumerate}

\noindent \textbf{Correct Answer:} (A)

\begin{tcolorbox}[solbox, title=Detailed Solution -- Question 118]
Characteristic equation is
$1+G(s)H(s)=0$
$1+\frac{K(s+1)}{s^{3}+a s^{2}+2 s+1}=0$
$s^{3}+a s^{2}+(2+K) s+(1+K)=0$
 array for this is Routh

\[
\begin{array}{l|lr} s^{3} & 1 & (2+K) \\ s^{2} & a & (1+K) \\ s^{1} & \frac{a(2+K)-(1+K)}{a} \\ s^{0} & (1+K) & \end{array}
\]

 For oscillation $\frac{a(2+K)-(1+K)}{a}=0$
$\Rightarrow \quad a=\left(\frac{1+K}{2+K}\right)$
 Now 

\[
\begin{aligned} a s^{2}+(1+K)&=0 \\ -a \omega^{2}+(1+K)&=0 \\ \text{given }\quad \omega&=2 \mathrm{rad} / \mathrm{sec} \\ -4 a+(1+K)&=0 \\ -4 \frac{(1+K)}{(2+K)}+(1+K)&=0 \\ -4(1+K)+(2+K)(1+K)&=0 \\ (1+K)[(2+K)-4]&=0 \\ K=-1,2 \end{aligned}
\]

 but K=-1 is not possible as system will not  oscillate for this as a=0 
 So, K=2 
$a=\frac{1+K}{2+K}=\frac{3}{4}=0.75$
\end{tcolorbox}

\vspace{1.5em}
\hrule
\vspace{1.5em}

\section*{Question 119: Frequency Response Analysis (GATE EC 2012)}
\addcontentsline{toc}{section}{Question 119: Frequency Response Analysis (GATE EC 2012)}
\noindent \textbf{\color{primary}MCQ} \hfill \textbf{\color{secondary}1 Mark}


\noindent A system with transfer function
 $G(s)=\frac{(s^{2}+9)(s+2)}{(s+1)(s+3)(s+4)}$ 
is excited by sin($\omega$ t). The steady-state output of the system is zero at


\begin{enumerate}[label=(\Alph*)]
    \item $\omega$=1 rad/s
    \item $\omega$=2 rad/s
    \item \textbf{(Correct)} $\omega$=3 rad/s
    \item $\omega$=4 rad/s
\end{enumerate}

\noindent \textbf{Correct Answer:} (C)

\begin{tcolorbox}[solbox, title=Detailed Solution -- Question 119]
Transfer function is

\[
\begin{aligned} G(s)&=\frac{\left(s^{2}+9\right)(s+2)}{(s+1)(s+3)(s+4)} \\ \text{Input }\quad X(s)&=L \cdot T[\sin \omega t]\\ X(s)&=\frac{\omega}{s^{2}+\omega^{2}} \end{aligned}
\]

 So output

\[
Y(s)=\frac{\left(s^{2}+9\right)(s+2)}{(s+1)(s+3)(s+4)} \cdot \frac{\omega}{\left(s^{2}+\omega^{2}\right)}
\]

 Now

\[
\lim _{t \rightarrow \infty} y(t)=\lim _{s \rightarrow 0} S Y(s)=\lim _{s \rightarrow 0} s \frac{\left(s^{2}+9\right)(s+2)}{(s+1)(s+3)(s+4)} \cdot \frac{\omega}{\left(s^{2}+\omega^{2}\right)}
\]

 this will be zero if $s^{2}+\omega^{2}$ will cancel $\left(s^{2}+9\right)$ term as only then final value theorem will be applicable. 
 So at $\omega^{2}=9 \Rightarrow \omega=3 \mathrm{rad} / \mathrm{sec},$ steady state output will be zero.
\end{tcolorbox}

\vspace{1.5em}
\hrule
\vspace{1.5em}

\section*{Question 120: Basics, Block Diagrams \& SFGs (GATE EC 2011)}
\addcontentsline{toc}{section}{Question 120: Basics, Block Diagrams \& SFGs (GATE EC 2011)}
\noindent \textbf{\color{primary}MCQ} \hfill \textbf{\color{secondary}1 Mark}


\noindent The input-output transfer function of a plant $H(s)=\frac{100}{s(s+10)^2}$.
The plant is
placed in a unity negative feedback configuration as shown in the figure below.

\begin{center}\includegraphics[max width=0.85\linewidth]{images/img\_34.jpg}\end{center}

\begin{center}\includegraphics[max width=0.85\linewidth]{images/img\_34.jpg}\end{center}

The signal flow graph that DOES NOT model the plant transfer function H(S) is

\begin{center}\includegraphics[max width=0.85\linewidth]{images/img\_30.jpg}\end{center}

\begin{center}\includegraphics[max width=0.85\linewidth]{images/img\_30.jpg}\end{center}


\begin{enumerate}[label=(\Alph*)]
    \item A
    \item B
    \item C
    \item \textbf{(Correct)} D
\end{enumerate}

\noindent \textbf{Correct Answer:} (D)

\begin{tcolorbox}[solbox, title=Detailed Solution -- Question 120]
For option (d),
$\frac{Y(s)}{U(s)}=\frac{100 / s^{3}}{1+\frac{100}{s^{2}}}$

Which is not transfer function of H(s) .
\end{tcolorbox}

\vspace{1.5em}
\hrule
\vspace{1.5em}

\section*{Question 121: Frequency Response Analysis (GATE EC 2011)}
\addcontentsline{toc}{section}{Question 121: Frequency Response Analysis (GATE EC 2011)}
\noindent \textbf{\color{primary}MCQ} \hfill \textbf{\color{secondary}1 Mark}


\noindent The input-output transfer function of a plant $H(s)=\frac{100}{s(s+10)^2}$.
The plant is
placed in a unity negative feedback configuration as shown in the figure below.

\begin{center}\includegraphics[max width=0.85\linewidth]{images/img\_34.jpg}\end{center}

\begin{center}\includegraphics[max width=0.85\linewidth]{images/img\_34.jpg}\end{center}

The gain margin of the system under closed loop unity negative feedback is


\begin{enumerate}[label=(\Alph*)]
    \item 0 dB
    \item 20 dB
    \item \textbf{(Correct)} 26 dB
    \item 46 dB
\end{enumerate}

\noindent \textbf{Correct Answer:} (C)

\begin{tcolorbox}[solbox, title=Detailed Solution -- Question 121]
\[
\begin{aligned} G(s) H(s) &=\frac{100}{s(s+10)^{2}} \\ \phi &=-90^{\circ}-2 \tan ^{-1}(\omega / 10) \end{aligned}
\]

 For phase cross-over frequency, $\phi=-180^{\circ}$

\[
\begin{aligned} \therefore \quad-180 &=-90^{\circ}-2 \tan ^{-1}(\omega / 10) \\ \Rightarrow & \omega=10 \mathrm{rad} / \mathrm{sec} \\ s &=j \omega \\ \text { Put, } \quad G(j \omega) H(j \omega) &=\frac{100}{j \omega(j \omega+100)^{2}} \\ |G(j \omega) H(j \omega)|_{\omega=10} &=\frac{100}{\omega\left(\omega^{2}+100\right)} \\ &=\frac{100}{10(200)}=\frac{1}{20} \\ \text { G.M. } &=\frac{1}{\mid G(j \omega) H(j \omega)}=20 \end{aligned}
\]

 G.M. in $\mathrm{dB}=20 \log 20=26 \mathrm{dB}$
\end{tcolorbox}

\vspace{1.5em}
\hrule
\vspace{1.5em}

\section*{Question 122: State Space Analysis (GATE EC 2011)}
\addcontentsline{toc}{section}{Question 122: State Space Analysis (GATE EC 2011)}
\noindent \textbf{\color{primary}MCQ} \hfill \textbf{\color{secondary}1 Mark}


\noindent The block diagram of a system with one input it and two outputs $y_1 \; and \; y_2$ is
given below. 

\begin{center}\includegraphics[max width=0.85\linewidth]{images/img\_55.jpg}\end{center}

\begin{center}\includegraphics[max width=0.85\linewidth]{images/img\_55.jpg}\end{center}
 
A state space model of the above system in terms of the state vector $\underline{x}$ and the output vector $\underline{y}= [y_{1}    \; y_{2}]^{T}$ is


\begin{enumerate}[label=(\Alph*)]
    \item $\dot{\underline{x}}=[2]\underline{x} + [1]u; \underline{y}=[1 \; 2]\underline{x}$
    \item \textbf{(Correct)} \[
\dot{\underline{x}}=[-2]\underline{x} + [1]u; \underline{y}=\begin{bmatrix} 1\\ 2 \end{bmatrix}\underline{x}
\]
    \item \[
\dot{\underline{x}}=\begin{bmatrix} -2 & 0\\ 0&-2 \end{bmatrix}\underline{x} +\begin{bmatrix} 1\\ 1 \end{bmatrix}u; \underline{y}=[1 \; 2]\underline{x}
\]
    \item \[
\dot{\underline{x}}=\begin{bmatrix} 2 & 0\\ 0&2 \end{bmatrix}\underline{x} + \begin{bmatrix} 1\\ 1 \end{bmatrix}u; \underline{y}=\begin{bmatrix} 1\\ 2 \end{bmatrix}\underline{x}
\]
\end{enumerate}

\noindent \textbf{Correct Answer:} (B)

\begin{tcolorbox}[solbox, title=Detailed Solution -- Question 122]
\[
\begin{aligned} \frac{Y_{1}(s)}{U(s)} &=\frac{1}{s+2} \\ \frac{Y_{2}(s)}{U(s)} &=\frac{2}{s+2} \\ \frac{Y_{1}(s)}{x_{1}(s)} \frac{x_{1}(s)}{U(s)} &=\frac{1}{s+2} \\ \frac{r_{2}(s)}{x_{2}(s)} \frac{x_{2}(s)}{U(s)} &=\frac{2}{s+2} \end{aligned}
\]

\[
\begin{aligned} \frac{X_{1}(s)}{U(s)} &=\frac{1}{s+2} \text { and } \frac{Y_{1}(s)}{X_{1}(s)}=1 \\ \frac{x_{2}(s)}{U(s)} &=\frac{1}{s+2} \text { and } \frac{Y_{2}(s)}{U(s)}=2 \\ s X_{1}(s)+2 X_{1}(s) &=U(s) \\ \text{and} \quad Y_{1}(s) &=x_{1}(s) \\ s X_{2}(s)+2 X_{2}(s) &=U(s) \\ \text{and }\quad Y_{2}(s)&=2 X_{2}(s) \\ \dot{x}_{1}(t)+2 x_{1}(t)&=U(t) \\ \text{and }y_{1}(t)&=x_{1}(t) \\ \dot{x}_{2}(t)+2 x_{2}(t)&=U(t) \\ y_{2}(t)&=2 x_{2}(t) \end{aligned}
\]

 From Question,

\[
\begin{aligned} y &=\left[y_{1} y_{2}\right]^{\top}=[12]^{\top}=\left[\begin{array}{l}1 \\2\end{array}\right] \\ \dot{x}_{1}(t) &=-2 x_{1}(t)+u(t) \\ \dot{x}_{2}(t) &=-2 x_{2}(t)+u(t) \\ \left[\begin{array}{c}\dot{x}_{1} \\\dot{x}_{2}\end{array}\right]&=\left[\begin{array}{cc}-2 & 0 \\-2 & 0\end{array}\right]\left[\begin{array}{l}x_{1} \\x_{2} \end{array}\right]+\left[\begin{array}{l}1 \\1\end{array}\right] u(t) \\ \dot{x} &=[-2] x+[1] u \end{aligned}
\]

 only option (B) is satisfied.
\end{tcolorbox}

\vspace{1.5em}
\hrule
\vspace{1.5em}

\section*{Question 123: Frequency Response Analysis (GATE EC 2011)}
\addcontentsline{toc}{section}{Question 123: Frequency Response Analysis (GATE EC 2011)}
\noindent \textbf{\color{primary}MCQ} \hfill \textbf{\color{secondary}1 Mark}


\noindent For the transfer function $G(j\omega)=5+j\omega$, the corresponding Nyquist plot for
positive frequency has the form 

\begin{center}\includegraphics[max width=0.85\linewidth]{images/img\_47.jpg}\end{center}

\begin{center}\includegraphics[max width=0.85\linewidth]{images/img\_47.jpg}\end{center}


\begin{enumerate}[label=(\Alph*)]
    \item \textbf{(Correct)} A
    \item B
    \item C
    \item D
\end{enumerate}

\noindent \textbf{Correct Answer:} (A)

\begin{tcolorbox}[solbox, title=Detailed Solution -- Question 123]
\[
\begin{aligned} G(j \omega) &=5+j \omega \\|G(j \omega)| &=\sqrt{25+\omega^{2}} \\ \text { At } \quad\omega&=0\\ At |G(0)|&=\sqrt{25+0}=5 \\ \text { At } \quad\omega&=\infty \\ |G(\infty)| &=\infty \end{aligned}
\]
\end{tcolorbox}

\vspace{1.5em}
\hrule
\vspace{1.5em}

\section*{Question 124: Time Response Analysis (GATE EC 2011)}
\addcontentsline{toc}{section}{Question 124: Time Response Analysis (GATE EC 2011)}
\noindent \textbf{\color{primary}MCQ} \hfill \textbf{\color{secondary}1 Mark}


\noindent The differential equation $100\frac{d^{2}y}{dt^{2}}-20\frac{dy}{dt}+y=x(t)$ describes a system with an in put x(t) and an output y(t). The system, which is initially relaxed, is excited by a unit step input. The output y(t) can be represented by the waveform 

\begin{center}\includegraphics[max width=0.85\linewidth]{images/img\_111.jpg}\end{center}

\begin{center}\includegraphics[max width=0.85\linewidth]{images/img\_111.jpg}\end{center}


\begin{enumerate}[label=(\Alph*)]
    \item \textbf{(Correct)} A
    \item B
    \item C
    \item D
\end{enumerate}

\noindent \textbf{Correct Answer:} (A)

\begin{tcolorbox}[solbox, title=Detailed Solution -- Question 124]
$100 \frac{d^{2} y}{d t^{2}}-20 \frac{d y}{d t}+y=x(t)$
 Taking Laplace transform of both sides
$100 s^{2} Y(s)-20 s Y(s)+Y(s)=\frac{X(s)}{s}$

\[
\Rightarrow \frac{Y(s)}{X(s)}=\frac{1 / s}{\left(100 s^{2}-20 s+1\right)}=\frac{1}{s(10 s-1)^{2}}
\]

 Poles are at $s=\frac{1}{10}, \frac{1}{10}, 0$
 As poles are on the right-hand side of s-plane so given system is unstable system.
 Only option (A) represents unstable system.
\end{tcolorbox}

\vspace{1.5em}
\hrule
\vspace{1.5em}

\section*{Question 125: Root Locus Technique (GATE EC 2011)}
\addcontentsline{toc}{section}{Question 125: Root Locus Technique (GATE EC 2011)}
\noindent \textbf{\color{primary}MCQ} \hfill \textbf{\color{secondary}1 Mark}


\noindent The root locus plot for a system is given below. The open loop transfer function
corresponding to this plot is given by 

\begin{center}\includegraphics[max width=0.85\linewidth]{images/img\_137.jpg}\end{center}

\begin{center}\includegraphics[max width=0.85\linewidth]{images/img\_137.jpg}\end{center}


\begin{enumerate}[label=(\Alph*)]
    \item $G(S)H(S)=k\frac{s(s+1)}{(s+2)(s+3)}$
    \item \textbf{(Correct)} $G(S)H(S)=k\frac{(s+1)}{(s+2)(s+3)^{2}}$
    \item $G(S)H(S)=k\frac{1}{s(s-1)(s+2)(s+3)}$
    \item $G(S)H(S)=k\frac{(s+1)}{(s+2)(s+3)}$
\end{enumerate}

\noindent \textbf{Correct Answer:} (B)

\begin{tcolorbox}[solbox, title=Detailed Solution -- Question 125]
From plot we can observe that one pole terminates
at one zero at position -1 and three poles
terminates to $\infty$. It means there are four poles and
1 zero. Pole at -3 goes on both sides. It means
there are two poles at -3.
\end{tcolorbox}

\vspace{1.5em}
\hrule
\vspace{1.5em}

\section*{Question 126: State Space Analysis (GATE EC 2010)}
\addcontentsline{toc}{section}{Question 126: State Space Analysis (GATE EC 2010)}
\noindent \textbf{\color{primary}MCQ} \hfill \textbf{\color{secondary}1 Mark}


\noindent The signal flow graph of a system is shown below: 

\begin{center}\includegraphics[max width=0.85\linewidth]{images/img\_112.jpg}\end{center}

\begin{center}\includegraphics[max width=0.85\linewidth]{images/img\_112.jpg}\end{center}

The transfer function of the system is


\begin{enumerate}[label=(\Alph*)]
    \item $\frac{s+1}{s^{2}+1}$
    \item $\frac{s-1}{s^{2}+1}$
    \item \textbf{(Correct)} $\frac{s+1}{s^{2}+s+1}$
    \item $\frac{s-1}{s^{2}+s+1}$
\end{enumerate}

\noindent \textbf{Correct Answer:} (C)

\begin{tcolorbox}[solbox, title=Detailed Solution -- Question 126]
Forward path gain,

\[
\begin{aligned} P_{1} &=2\left(\frac{1}{s}\right)\left(\frac{1}{s}\right)(0.5)=\frac{1}{s^{2}} \\ P_{2} &=2\left(\frac{1}{s}\right)(1)(0.5)=\frac{1}{s} \\ \Delta_{1} &=1 \\ \Delta_{2} &=1 \\ \Delta &=1-{-\frac{1}{s}-\frac{1}{s^{2}}} \\ &=1+\frac{1}{s}+\frac{1}{s^{2}} \end{aligned}
\]

 Transfer funct ion of the system,

\[
\begin{aligned} \frac{Y(s)}{U(s)} &=\frac{P_{1} \Delta_{1}+P_{2} \Delta_{2}}{\Delta} \\ &=-\frac{1}{1+\frac{s^{2}}{s}+\frac{1}{s}}=\frac{s+1}{s^{2}+s+1} \end{aligned}
\]
\end{tcolorbox}

\vspace{1.5em}
\hrule
\vspace{1.5em}

\section*{Question 127: State Space Analysis (GATE EC 2010)}
\addcontentsline{toc}{section}{Question 127: State Space Analysis (GATE EC 2010)}
\noindent \textbf{\color{primary}MCQ} \hfill \textbf{\color{secondary}1 Mark}


\noindent The signal flow graph of a system is shown below: 

\begin{center}\includegraphics[max width=0.85\linewidth]{images/img\_112.jpg}\end{center}

\begin{center}\includegraphics[max width=0.85\linewidth]{images/img\_112.jpg}\end{center}

The state variable representation of the system can be


\begin{enumerate}[label=(\Alph*)]
    \item \[
\dot{x}=\begin{vmatrix} 1 & 1\\ -1&0 \end{vmatrix}x+\begin{bmatrix} 0\\ 2 \end{bmatrix}u
\]
 

\[
y=\begin{bmatrix} 0 &0.5 \end{bmatrix}x
\]
    \item \[
\dot{x}=\begin{vmatrix} -1 & 1\\ -1&0 \end{vmatrix}x+\begin{bmatrix} 0\\ 2 \end{bmatrix}u
\]

 
\[
y=\begin{bmatrix} 0 &0.5 \end{bmatrix}x
\]
    \item \[
\dot{x}=\begin{vmatrix} 1 & 1\\ -1&0 \end{vmatrix}x+\begin{bmatrix} 0\\ 2 \end{bmatrix}u
\]
 

\[
y=\begin{bmatrix} 0.5 &0.5 \end{bmatrix}x
\]
    \item \textbf{(Correct)} \[
\dot{x}=\begin{vmatrix} -1 & 1\\ -1&0 \end{vmatrix}x+\begin{bmatrix} 0\\ 2 \end{bmatrix}u
\]
 

\[
y=\begin{bmatrix} 0.5 &0.5 \end{bmatrix}x
\]
\end{enumerate}

\noindent \textbf{Correct Answer:} (D)


\vspace{1.5em}
\hrule
\vspace{1.5em}

\section*{Question 128: Compensators \& Controllers (GATE EC 2010)}
\addcontentsline{toc}{section}{Question 128: Compensators \& Controllers (GATE EC 2010)}
\noindent \textbf{\color{primary}MCQ} \hfill \textbf{\color{secondary}1 Mark}


\noindent A unity negative feedback closed loop system has a plant with the transfer function $G(s)=\frac{1}{s^{2}+2s+2}$ and a controller $G_{c}(s)$ in the
feed forward path. For a unit set input, the transfer function of the controller that
gives minimum steady state error is


\begin{enumerate}[label=(\Alph*)]
    \item $G_{c}(s)=\frac{s+1}{s+2}$
    \item $G_{c}(s)=\frac{(s+1)(s+4)}{(s+2)(s+3)}$
    \item $G_{c}(s)=\frac{s+2}{s+1}$
    \item \textbf{(Correct)} $G_{c}(s)=1+\frac{2}{s}+3s$
\end{enumerate}

\noindent \textbf{Correct Answer:} (D)

\begin{tcolorbox}[solbox, title=Detailed Solution -- Question 128]
Steady state error,

\[
\begin{aligned} e_{s s} &=\lim _{s \rightarrow 0} \frac{s R(s)}{1+G(s) G_{c}(s)} \\ r(t) &=u(t) \\ R(s) &=\frac{1}{s} \\ e_{ss} &=\lim _{s \rightarrow 0} \frac{s \cdot \frac{1}{s}}{1+G(s) G_{c}(s)} \\ e_{ss}&=\lim _{s \rightarrow 0} \frac{1}{1+G(s) G_{c}(s)} \\ \text { Taking, } \quad G_{c}(s) &=\frac{s+1}{s+2}, e_{s s}=\frac{2}{3} \\ \text { Taking, } \quad G_{c}(s) &=\frac{s+2}{s+1}, e_{s s}=\frac{1}{3} \\ \text { Taking, } \quad G_{c}(s) &=\frac{(s+1)(s+4)}{(s+2)(s+3)}, e_{s s}=\frac{3}{5} \\ \text { Taking, } \quad & G_{c}(s)=1+\frac{2}{s}+3 s, e_{s s}=0 \end{aligned}
\]
\end{tcolorbox}

\vspace{1.5em}
\hrule
\vspace{1.5em}

\section*{Question 129: Frequency Response Analysis (GATE EC 2010)}
\addcontentsline{toc}{section}{Question 129: Frequency Response Analysis (GATE EC 2010)}
\noindent \textbf{\color{primary}MCQ} \hfill \textbf{\color{secondary}1 Mark}


\noindent For the asymptotic Bode magnitude plot shown below, the system transfer function can be 

\begin{center}\includegraphics[max width=0.85\linewidth]{images/img\_62.jpg}\end{center}

\begin{center}\includegraphics[max width=0.85\linewidth]{images/img\_62.jpg}\end{center}


\begin{enumerate}[label=(\Alph*)]
    \item \textbf{(Correct)} $\frac{10s+1}{0.1s+1}$
    \item $\frac{100s+1}{0.1s+1}$
    \item $\frac{100s}{10s+1}$
    \item $\frac{0.1s+1}{10s+1}$
\end{enumerate}

\noindent \textbf{Correct Answer:} (A)

\begin{tcolorbox}[solbox, title=Detailed Solution -- Question 129]
\[
\begin{array}{l} \text { system transter function, } \\ \qquad G(s) H(s)=\frac{K\left(1+\frac{s}{0.1}\right)}{\left(1+\frac{s}{10}\right)} \\ \text { Here, } 20 \log K=0 \\ \Rightarrow \quad K=1 \\ \text { Therefore, } G(s) H(s)=\frac{10 s+1}{0.1 s+1} \end{array}
\]
\end{tcolorbox}

\vspace{1.5em}
\hrule
\vspace{1.5em}

\section*{Question 130: Frequency Response Analysis (GATE EC 2010)}
\addcontentsline{toc}{section}{Question 130: Frequency Response Analysis (GATE EC 2010)}
\noindent \textbf{\color{primary}MCQ} \hfill \textbf{\color{secondary}1 Mark}


\noindent A system with transfer function $\frac{Y(s)}{X(s)}=\frac{s}{s+p}$ has an output $y(t)=cos(2t-\frac{\pi }{3})$ for the input signal $x(t)=pcos(2t-\frac{\pi }{2})$. Then, the system parameter p is


\begin{enumerate}[label=(\Alph*)]
    \item $\sqrt{3}$
    \item \textbf{(Correct)} $\frac{2}{\sqrt{3}}$
    \item 1
    \item $\frac{\sqrt{3}}{2}$
\end{enumerate}

\noindent \textbf{Correct Answer:} (B)

\begin{tcolorbox}[solbox, title=Detailed Solution -- Question 130]
\[
\begin{aligned} \phi&=-\frac{\pi}{3}-\left(-\frac{\pi}{2}\right)=\frac{\pi}{6}=30^{\circ}\\ \text{and }\quad\omega&=2 \mathrm{rad} / \mathrm{s} \end{aligned}
\]

 From the transfer function,

\[
\begin{aligned} \phi &=90^{\circ}-\tan ^{-1} \frac{\omega}{p} \\ 90^{\circ}-\tan ^{-1} \frac{2}{p} &=30^{\circ} \\ \tan ^{-1} \frac{2}{p} &=60^{\circ} \\ \frac{2}{p} &=\tan 60^{\circ}=\sqrt{3} \\ p &=\frac{2}{\sqrt{3}} \end{aligned}
\]
\end{tcolorbox}

\vspace{1.5em}
\hrule
\vspace{1.5em}

\section*{Question 131: Basics, Block Diagrams \& SFGs (GATE EC 2010)}
\addcontentsline{toc}{section}{Question 131: Basics, Block Diagrams \& SFGs (GATE EC 2010)}
\noindent \textbf{\color{primary}MCQ} \hfill \textbf{\color{secondary}1 Mark}


\noindent The transfer function Y(s)/R(s) of the system shown is 

\begin{center}\includegraphics[max width=0.85\linewidth]{images/img\_35.jpg}\end{center}

\begin{center}\includegraphics[max width=0.85\linewidth]{images/img\_35.jpg}\end{center}


\begin{enumerate}[label=(\Alph*)]
    \item 0
    \item \textbf{(Correct)} $\frac{1}{s+1}$
    \item $\frac{2}{s+1}$
    \item $\frac{2}{s+3}$
\end{enumerate}

\noindent \textbf{Correct Answer:} (B)

\begin{tcolorbox}[solbox, title=Detailed Solution -- Question 131]
\begin{center}\includegraphics[max width=0.85\linewidth]{images/img\_38.jpg}\end{center}

\begin{center}\includegraphics[max width=0.85\linewidth]{images/img\_38.jpg}\end{center}

\[
\begin{aligned} Q(s) &=P(s)\left[\frac{1}{s+1}-\frac{1}{s+1}\right] \\ &=0 \\ \text{So,}\quad P(s) &=R(s)-0=R(s)\\ Y(s)&=\frac{P(s)}{s+1}=\frac{R(s)}{s+1}\\ \text{Therefore, }\quad\frac{Y(s)}{R(s)}&=\frac{1}{s+1} \end{aligned}\\
\]
\end{tcolorbox}

\vspace{1.5em}
\hrule
\vspace{1.5em}

\section*{Question 132: Frequency Response Analysis (GATE EC 2009)}
\addcontentsline{toc}{section}{Question 132: Frequency Response Analysis (GATE EC 2009)}
\noindent \textbf{\color{primary}MCQ} \hfill \textbf{\color{secondary}1 Mark}


\noindent The Nyquist plot of a stable transfer function G(s) is shown in the figure are
interested in the stability of the closed loop system in the feedback configuration
shown.

\begin{center}\includegraphics[max width=0.85\linewidth]{images/img\_54.jpg}\end{center}

\begin{center}\includegraphics[max width=0.85\linewidth]{images/img\_54.jpg}\end{center}

The gain and phase margins of G(s) for closed loop stability are


\begin{enumerate}[label=(\Alph*)]
    \item 6 dB and $180^{\circ}$
    \item 3 dB and $180^{\circ}$
    \item \textbf{(Correct)} 6 dB and $90^{\circ}$
    \item 3 dB and $90^{\circ}$
\end{enumerate}

\noindent \textbf{Correct Answer:} (C)

\begin{tcolorbox}[solbox, title=Detailed Solution -- Question 132]
\[
\begin{aligned} \text{Gain Margin }&= 20 \log\frac{1}{X}\\ &=20 \log \frac{1}{0.5} \\ &=20 \log 2=6 \mathrm{dB} \\ \text { and phase margin} &=90^{\circ} \end{aligned}
\]
\end{tcolorbox}

\vspace{1.5em}
\hrule
\vspace{1.5em}

\section*{Question 133: Frequency Response Analysis (GATE EC 2009)}
\addcontentsline{toc}{section}{Question 133: Frequency Response Analysis (GATE EC 2009)}
\noindent \textbf{\color{primary}MCQ} \hfill \textbf{\color{secondary}1 Mark}


\noindent The Nyquist plot of a stable transfer function G(s) is shown in the figure are
interested in the stability of the closed loop system in the feedback configuration
shown.

\begin{center}\includegraphics[max width=0.85\linewidth]{images/img\_54.jpg}\end{center}

\begin{center}\includegraphics[max width=0.85\linewidth]{images/img\_54.jpg}\end{center}

Which of the following statements is true ?


\begin{enumerate}[label=(\Alph*)]
    \item G(s) is an all-pass filter
    \item \textbf{(Correct)} G(s) has a zero in the right-half plane
    \item G(s) is the impedance of a passive network
    \item G(s) is marginally stable
\end{enumerate}

\noindent \textbf{Correct Answer:} (B)

\begin{tcolorbox}[solbox, title=Detailed Solution -- Question 133]
Nyquist plot will encirtce the origin as many
number of times as the difference between the
number of right side poles and zeros of open loop
transfer function.

i.e. $\quad N=P-Z$

From the plot,

\[
\begin{aligned}
N&=-1\\
\text{and  }P&=0 \quad  (Given)  \\
\therefore \quad Z&=1
\end{aligned}
\]

Thus, one open loop zero in right half of s-Plane
\end{tcolorbox}

\vspace{1.5em}
\hrule
\vspace{1.5em}

\section*{Question 134: Time Response Analysis (GATE EC 2009)}
\addcontentsline{toc}{section}{Question 134: Time Response Analysis (GATE EC 2009)}
\noindent \textbf{\color{primary}MCQ} \hfill \textbf{\color{secondary}1 Mark}


\noindent The unit step response of an under-damped second order system has steady state value of -2. Which one of the following transfer functions has theses properties ?


\begin{enumerate}[label=(\Alph*)]
    \item $\frac{-2.24}{s^{2}+2.59s+1.12}$
    \item \textbf{(Correct)} $\frac{-3.82}{s^{2}+1.91s+1.91}$
    \item $\frac{-2.24}{s^{2}-2.59s+1.12}$
    \item $\frac{-3.82}{s^{2}-1.91s+1.91}$
\end{enumerate}

\noindent \textbf{Correct Answer:} (B)

\begin{tcolorbox}[solbox, title=Detailed Solution -- Question 134]
Steady state Value =-2 
 denominator:

\[
\begin{aligned} 2 \xi \omega_{n} &=1.91, \omega_{n}^{2}=1.91 \\ \Rightarrow \quad \omega_{n} & \cong 1.4 \\ \xi &=\frac{1.91}{2 \omega_{n}}=\frac{1.91}{2.8} < 1\\ &...\text{underdamped} \end{aligned}
\]
\end{tcolorbox}

\vspace{1.5em}
\hrule
\vspace{1.5em}

\section*{Question 135: Root Locus Technique (GATE EC 2009)}
\addcontentsline{toc}{section}{Question 135: Root Locus Technique (GATE EC 2009)}
\noindent \textbf{\color{primary}MCQ} \hfill \textbf{\color{secondary}1 Mark}


\noindent The feedback configuration and the pole-zero locations of $G(s)=\frac{s^{2}-2s+2}{s^{2}+2s+2}$ are shown below. The root locus for negative values of k , i.e. for $-\infty < k < 0$, has breakaway/break-in points and angle of departure at pole P (with respect to the positive real axis) equal to 

\begin{center}\includegraphics[max width=0.85\linewidth]{images/img\_51.jpg}\end{center}

\begin{center}\includegraphics[max width=0.85\linewidth]{images/img\_51.jpg}\end{center}


\begin{enumerate}[label=(\Alph*)]
    \item $\pm \sqrt{2} \; and \; 0^{\circ}$
    \item \textbf{(Correct)} $\pm \sqrt{2} \; and \;  45^{\circ}$
    \item $\pm \sqrt{3} \; and \;  0^{\circ}$
    \item $\pm \sqrt{3} \; and \;  45^{\circ}$
\end{enumerate}

\noindent \textbf{Correct Answer:} (B)

\begin{tcolorbox}[solbox, title=Detailed Solution -- Question 135]
\begin{center}\includegraphics[max width=0.85\linewidth]{images/img\_136.jpg}\end{center}

\begin{center}\includegraphics[max width=0.85\linewidth]{images/img\_136.jpg}\end{center}

\[
\begin{aligned} 1+G(s) H(s) &=0 \\ 1+\frac{K\left(s^{2}-2 s+2\right)}{s^{2}+2 s+2} &=0 \\ k &=-\frac{s^{2}+2 s+2}{s^{2}-2 s+2} \\ \text{put }\quad \frac{\partial K}{\partial s}&=0\text{ we have}\\ \left(s^{2}-2 s+2\right)(s+1) &-\left(s^{2}+2 s+2\right) \\ (s-1) =0 \\ 2 s^{2}-4 s^{2}+4 &=0 \\ 2 s^{2} &=+4 \\ s &=\pm \sqrt{2} \\ \text{Angle of departure is}\\ \phi_{D}=\phi\\ \text { where } &=\Sigma \phi_{2}-\Sigma \phi_{p} \\ &=225^{\circ} \\ \phi_{D} &=\pm 225^{\circ} \end{aligned} \\ \text{No option is matching.}
\]
\end{tcolorbox}

\vspace{1.5em}
\hrule
\vspace{1.5em}

\section*{Question 136: State Space Analysis (GATE EC 2009)}
\addcontentsline{toc}{section}{Question 136: State Space Analysis (GATE EC 2009)}
\noindent \textbf{\color{primary}MCQ} \hfill \textbf{\color{secondary}1 Mark}


\noindent Consider the system $\frac{dx}{dt}=Ax+Bu$ with 
\[
A=\begin{bmatrix} 1 & 0\\ 0&1 \end{bmatrix}
\]
 and 
\[
B=\begin{bmatrix} p\\ q \end{bmatrix}
\]
 where p and q are arbitrary real numbers. Which of the following statements about the controllability of the system is true ?


\begin{enumerate}[label=(\Alph*)]
    \item The system is completely state controllable for any nonzero values of p and q
    \item Only p = 0 and q = 0 result in controllability
    \item \textbf{(Correct)} The system is uncontrollable for all values of p and q
    \item We cannot conclude about controllability from the given data
\end{enumerate}

\noindent \textbf{Correct Answer:} (C)

\begin{tcolorbox}[solbox, title=Detailed Solution -- Question 136]
\[
\begin{aligned} A&=\left[\begin{array}{ll}1 & 0 \\0 & 1\end{array}\right] \\ B&=\left[\begin{array}{l}p \\q\end{array}\right] \end{aligned}
\]

 For controllabil itycondition is

\[
\begin{aligned} Q_{C}&=\left[B, A B, \ldots A^{n-1} B\right] \neq 0\\ A \cdot B&=\left[\begin{array}{ll}1 & 0 \\0 & 1\end{array}\right]\left[\begin{array}{l}p \\q\end{array}\right]=\left[\begin{array}{l}p+0 \\0+q\end{array}\right]=\left[\begin{array}{l}p \\q\end{array}\right]\\ \text{So}\quad Q_{c}&=\left[\begin{array}{ll}p & p \\ q & q\end{array}\right]=0 \end{aligned}
\]

 So, the system is uncontrollable for all values of p and q
\end{tcolorbox}

\vspace{1.5em}
\hrule
\vspace{1.5em}

\section*{Question 137: Compensators \& Controllers (GATE EC 2009)}
\addcontentsline{toc}{section}{Question 137: Compensators \& Controllers (GATE EC 2009)}
\noindent \textbf{\color{primary}MCQ} \hfill \textbf{\color{secondary}1 Mark}


\noindent The magnitude plot of a rational transfer function G(s) with real coefficients is
shown below. Which of the following compensators has such a magnitude plot ? 

\begin{center}\includegraphics[max width=0.85\linewidth]{images/img\_71.jpg}\end{center}

\begin{center}\includegraphics[max width=0.85\linewidth]{images/img\_71.jpg}\end{center}


\begin{enumerate}[label=(\Alph*)]
    \item Lead compensator
    \item Lag compensator
    \item \textbf{(Correct)} PID compensator
    \item Lead-lag compensator
\end{enumerate}

\noindent \textbf{Correct Answer:} (C)


\vspace{1.5em}
\hrule
\vspace{1.5em}

\section*{Question 138: Compensators \& Controllers (GATE EC 2008)}
\addcontentsline{toc}{section}{Question 138: Compensators \& Controllers (GATE EC 2008)}
\noindent \textbf{\color{primary}MCQ} \hfill \textbf{\color{secondary}1 Mark}


\noindent Group I gives two possible choices for the impedance Z in the diagram. The
circuit elements in Z satisfy the conditions $R_{2}C_{2} > R_{1}C_{1}$. The transfer functions $\frac{V_{o}}{V_{i}}$ represents a kind of controller. Match the impedances in Group I with the type of controllers in Group II 

\begin{center}\includegraphics[max width=0.85\linewidth]{images/img\_49.jpg}\end{center}

\begin{center}\includegraphics[max width=0.85\linewidth]{images/img\_49.jpg}\end{center}

\begin{center}\includegraphics[max width=0.85\linewidth]{images/img\_91.jpg}\end{center}

\begin{center}\includegraphics[max width=0.85\linewidth]{images/img\_91.jpg}\end{center}


\begin{enumerate}[label=(\Alph*)]
    \item Q - 1,R - 2
    \item \textbf{(Correct)} Q - 1,R - 3
    \item Q - 2,R - 3
    \item Q - 3,R - 2
\end{enumerate}

\noindent \textbf{Correct Answer:} (B)

\begin{tcolorbox}[solbox, title=Detailed Solution -- Question 138]
\[
\begin{array}{c} \frac{V_{i}\left(R_{1} C_{1} S+1\right) 1}{R_{1}}=-\frac{V_{0}}{Z} \\ \Rightarrow \frac{V_{0}}{V_{i}}=-\frac{Z\left(R_{1} C_{1} S+1\right)}{R_{1}} \end{array}
\]

 In case of Q, $\quad Z=\frac{R_{2} C_{2} s+1}{C_{2} S+1}$
 In case of R, $\quad Z=\frac{R_{2}}{R_{2} C_{2} s+1}$
 Considering Q 

\[
\frac{V_{0}}{V_{i}}=-\frac{\left(R_{1} C_{1} S+1\right)}{R_{1}} \cdot \frac{\left(R_{2} C_{2} S+1\right)}{C_{2} S}
\]

 Considering R ,

\[
\frac{V_{0}}{V_{i}}=-\frac{\left(R_{1} C_{1} S+1\right)}{R_{1}} \cdot \frac{R_{2}}{\left(R_{2} C_{2} s+1\right)}
\]

$\because$ Given that $R_{2} C_{2}>R_{1} C_{1}$
$\therefore \quad$ Considering R, controller is lag compensator. 
 and Considering Q , Controller is PID controller.
\end{tcolorbox}

\vspace{1.5em}
\hrule
\vspace{1.5em}

\section*{Question 139: Frequency Response Analysis (GATE EC 2008)}
\addcontentsline{toc}{section}{Question 139: Frequency Response Analysis (GATE EC 2008)}
\noindent \textbf{\color{primary}MCQ} \hfill \textbf{\color{secondary}1 Mark}


\noindent The magnitude of frequency responses of an underdamped second order system
is 5 at 0 rad/sec and peaks to $\frac{10}{\sqrt{3}}$ at 5$\sqrt{3}$
rad/sec. The transfer function of the system is


\begin{enumerate}[label=(\Alph*)]
    \item \textbf{(Correct)} $\frac{500}{s^{2}+10s+100}$
    \item $\frac{375}{s^{2}+5s+75}$
    \item $\frac{720}{s^{2}+12s+144}$
    \item $\frac{1125}{s^{2}+5s+225}$
\end{enumerate}

\noindent \textbf{Correct Answer:} (A)

\begin{tcolorbox}[solbox, title=Detailed Solution -- Question 139]
\[
\begin{aligned} M_{r}&=\frac{5}{2 \xi \sqrt{1-\xi^{2}}}=\frac{10}{\sqrt{3}} \\ \text{or}\quad\xi&=\frac{1}{2}\\ \text{Also, }\omega_{r}&=\omega_{n} \sqrt{1-2 z^{2}} \\ 5 \sqrt{2} &=\omega_{n} \sqrt{1-2 x\left(\frac{1}{4}\right)} \\ \text{or}\quad\omega_{n} &=10 \mathrm{rad} / \mathrm{sec} \end{aligned}
\]

 Only option (A) satisfies the conditions.
\end{tcolorbox}

\vspace{1.5em}
\hrule
\vspace{1.5em}

\section*{Question 140: Stability Analysis \& Routh-Hurwitz (GATE EC 2008)}
\addcontentsline{toc}{section}{Question 140: Stability Analysis \& Routh-Hurwitz (GATE EC 2008)}
\noindent \textbf{\color{primary}MCQ} \hfill \textbf{\color{secondary}1 Mark}


\noindent The number of open right half plane of 
$G(s)=\frac{10}{s^{5}+2s^{4}+3s^{3}+6s^{2}+5s+3}$ is


\begin{enumerate}[label=(\Alph*)]
    \item 0
    \item 1
    \item \textbf{(Correct)} 2
    \item 3
\end{enumerate}

\noindent \textbf{Correct Answer:} (C)

\begin{tcolorbox}[solbox, title=Detailed Solution -- Question 140]
Characteristic equation
$q(s)=s^{5}+2 s^{4}+3 s^{3}+6 s^{2}+5 s+3$
 Putting $s=\frac{1}{z}$
$q^{\prime}(z)=3 z^{5}+5 z^{4}+6 z^{3}+3 z^{2}+2 z+1$
 Routh array is

\[
\begin{array}{l|lll} z^{5} & 3 & 6 & 2 \\ z^{4} & 5 & 3 & 1 \\ z^{3} & 21 / 5 & 7 / 5 & \\ z^{2} & \frac{63}{51} =\frac{4}{3} & 1 &\\ z^{1} & \frac{\frac{28}{15}-\frac{21}{5}}{4 / 3} =\frac{-35 / 15}{4 / 3}=\frac{-7}{4} & &\\ z^{0} & 1 & & \end{array}
\]

 Sinc\_e sign changes twice in Routh-array, therefore, there are two poles on right half plane.
\end{tcolorbox}

\vspace{1.5em}
\hrule
\vspace{1.5em}

\section*{Question 141: State Space Analysis (GATE EC 2008)}
\addcontentsline{toc}{section}{Question 141: State Space Analysis (GATE EC 2008)}
\noindent \textbf{\color{primary}MCQ} \hfill \textbf{\color{secondary}1 Mark}


\noindent A signal flow graph of a system is given below. The set of equalities that corresponds to this signal flow graph is

\begin{center}\includegraphics[max width=0.85\linewidth]{images/img\_41.jpg}\end{center}

\begin{center}\includegraphics[max width=0.85\linewidth]{images/img\_41.jpg}\end{center}

\begin{center}\includegraphics[max width=0.85\linewidth]{images/img\_109.jpg}\end{center}

\begin{center}\includegraphics[max width=0.85\linewidth]{images/img\_109.jpg}\end{center}


\begin{enumerate}[label=(\Alph*)]
    \item A
    \item B
    \item C
    \item \textbf{(Correct)} D
\end{enumerate}

\noindent \textbf{Correct Answer:} (D)

\begin{tcolorbox}[solbox, title=Detailed Solution -- Question 141]
\begin{center}\includegraphics[max width=0.85\linewidth]{images/img\_75.jpg}\end{center}

\begin{center}\includegraphics[max width=0.85\linewidth]{images/img\_75.jpg}\end{center}

By solving this we get 

\[
\frac{d}{d t}\left(\begin{array}{l}x_{1} \\ x_{2} \\ x_{3}\end{array}\right)=\left[\begin{array}{ccc}-\gamma & 0 & \beta \\ \gamma & 0 & \alpha \\ -\beta & 0 & -\alpha\end{array}\right]\left(\begin{array}{c}x_{1} \\ x_{2} \\ x_{3}\end{array}\right)+\left[\begin{array}{cc}0 & 1 \\ 0 & 0 \\ 1 & 0\end{array}\right]\left(\begin{array}{c}u_{1} \\ u_{2}\end{array}\right)
\]
\end{tcolorbox}

\vspace{1.5em}
\hrule
\vspace{1.5em}

\section*{Question 142: Stability Analysis \& Routh-Hurwitz (GATE EC 2008)}
\addcontentsline{toc}{section}{Question 142: Stability Analysis \& Routh-Hurwitz (GATE EC 2008)}
\noindent \textbf{\color{primary}MCQ} \hfill \textbf{\color{secondary}1 Mark}


\noindent A certain system has transfer function $G(s)=\frac{s+8}{s^{2}+\alpha s-4}$, where $\alpha$ is a parameter. Consider the standard negative unity feedback configuration as shown below 

\begin{center}\includegraphics[max width=0.85\linewidth]{images/img\_105.jpg}\end{center}

\begin{center}\includegraphics[max width=0.85\linewidth]{images/img\_105.jpg}\end{center}
 
 Which of the following statements is true?


\begin{enumerate}[label=(\Alph*)]
    \item The closed loop systems is never stable for any value of $\alpha$
    \item For some positive value of $\alpha$, the closed loop system is stable, but not for all positive values.
    \item \textbf{(Correct)} For all positive values of $\alpha$, the closed loop system is stable.
    \item The closed loop system stable for all values of $\alpha$, both positive and negative.
\end{enumerate}

\noindent \textbf{Correct Answer:} (C)

\begin{tcolorbox}[solbox, title=Detailed Solution -- Question 142]
Closed loop gain is
$\frac{G(s)}{1+G(s)}=\frac{s+8}{s^{2}+\alpha s-4+s+8}$

Characteristic equation 
$q(s)=s^{2}+(\alpha+1) s+4$

Closed loop system is stable only for $\alpha > -1.$. Therefore, for all positive of $\alpha$, the closed
loop system is stable.
\end{tcolorbox}

\vspace{1.5em}
\hrule
\vspace{1.5em}

\section*{Question 143: Time Response Analysis (GATE EC 2008)}
\addcontentsline{toc}{section}{Question 143: Time Response Analysis (GATE EC 2008)}
\noindent \textbf{\color{primary}MCQ} \hfill \textbf{\color{secondary}1 Mark}


\noindent Group I lists a set of four transfer functions. Group II gives a list of possible
step response y(t). Match the step responses with the corresponding transfer
functions.

\begin{center}\includegraphics[max width=0.85\linewidth]{images/img\_127.jpg}\end{center}

\begin{center}\includegraphics[max width=0.85\linewidth]{images/img\_127.jpg}\end{center}


\begin{enumerate}[label=(\Alph*)]
    \item P - 3,Q - 1,R - 4,S - 2
    \item P - 3,Q - 2,R - 4,S - 1
    \item P - 2,Q - 1,R - 4,S - 2
    \item \textbf{(Correct)} P - 3,Q - 4,R - 1,S - 2
\end{enumerate}

\noindent \textbf{Correct Answer:} (D)

\begin{tcolorbox}[solbox, title=Detailed Solution -- Question 143]
Comparing the given transfer functions with

\[
\begin{array}{c} \frac{\omega_{n}^{2}}{s^{2}+2 \xi \omega_{n} s+\omega_{n}^{2}} \\ P=\frac{25}{s^{2}+25} \Rightarrow \xi=0 \end{array}
\]

 So, P is undamped.
$Q=\frac{36}{s^{2}+20 s+36} \Rightarrow \xi=\frac{20}{2 \times 6}=1.67$
 So, Q is overdamped.
$R=\frac{36}{s^{2}+12 s+36} \Rightarrow \xi=\frac{12}{2 \times 6}=1$
 So, R is critically damped.
$S=\frac{49}{s^{2}+7 s+49} \Rightarrow \xi=\frac{7}{2 \times 7}=0.5$
 So, S is underdamped.
\end{tcolorbox}

\vspace{1.5em}
\hrule
\vspace{1.5em}

\section*{Question 144: Time Response Analysis (GATE EC 2008)}
\addcontentsline{toc}{section}{Question 144: Time Response Analysis (GATE EC 2008)}
\noindent \textbf{\color{primary}MCQ} \hfill \textbf{\color{secondary}1 Mark}


\noindent A linear, time - invariant, causal continuous time system has a rational transfer
function with simple poles at s =- 2 and s =- 4 and one simple zero at s =- 1
. A unit step u(t) is applied at the input of the system. At steady state, the
output has constant value of 1. The impulse response of this system is


\begin{enumerate}[label=(\Alph*)]
    \item [exp(- 2t) + exp(- 4t)]u(t)
    \item [- 4 exp(- 2t) - 12 exp(- 4t) - exp(- t)]u(t)
    \item \textbf{(Correct)} [- 4 exp(- 2t) + 12 exp(- 4t)]u(t)
    \item [- 0.5 exp(- 2t) + 1.5 exp(- 4t)]u(t)
\end{enumerate}

\noindent \textbf{Correct Answer:} (C)

\begin{tcolorbox}[solbox, title=Detailed Solution -- Question 144]
Transfer functions can be written as,

\[
\begin{aligned} H(s)&=\frac{K(s+1)}{(s+2)(s+4)} \\ \text{Input,}\quad R(s)&=\frac{1}{s}\\ C(s)&=R(s) H(s) \\ \text{Given:} \quad \lim _{s \rightarrow 0} s C(s)&=1 \\ \text{or,} \quad\lim _{s \rightarrow 0} \frac{s \cdot K(s+1)}{s(s+2)(s+4)}&=1 \\ \text{Or }\quad \frac{K}{8}&=1 \\ \Rightarrow \quad K &=8 \\ H(s) &=\frac{8(s+1)}{(s+2)(s+4)} \\ &=-\frac{4}{s+2}+\frac{12}{s+4} \\ h(t) &=\left(-4 e^{-2 t}+12 e^{-4 t}\right) u(t) \end{aligned}
\]

Which is also the required impulse response of the system.
\end{tcolorbox}

\vspace{1.5em}
\hrule
\vspace{1.5em}

\section*{Question 145: State Space Analysis (GATE EC 2008)}
\addcontentsline{toc}{section}{Question 145: State Space Analysis (GATE EC 2008)}
\noindent \textbf{\color{primary}MCQ} \hfill \textbf{\color{secondary}1 Mark}


\noindent Consider the matrix 
\[
P=\begin{bmatrix} 0 & 1\\ -2&-3 \end{bmatrix}
\]
. Thr value of $e^{P}$ is


\begin{enumerate}[label=(\Alph*)]
    \item \[
\begin{bmatrix} 2 e^{-2}-3e^{-1}&e^{-1}-e^{-2} \\ 2e^{-2}-2e^{-1}& 5e^{-2}-e^{-1} \end{bmatrix}
\]
    \item \[
\begin{bmatrix}  e^{-2}-e^{-1}&2e^{-1}-e^{-2} \\ 2e^{-2}-4e^{-1}& 3e^{-2}-2e^{-1} \end{bmatrix}
\]
    \item \[
\begin{bmatrix} 5 e^{-2}-e^{-1}&3e^{-1}-e^{-2} \\ 2e^{-2}-6e^{-1}& 4e^{-2}-e^{-1} \end{bmatrix}
\]
    \item \textbf{(Correct)} \[
\begin{bmatrix} 2 e^{-1}-e^{-2}& e^{-1}-e^{-2} \\ -2e^{-1}+2e^{-2}& -e^{-1}+2e^{-2} \end{bmatrix}
\]
\end{enumerate}

\noindent \textbf{Correct Answer:} (D)


\vspace{1.5em}
\hrule
\vspace{1.5em}

\section*{Question 146: Time Response Analysis (GATE EC 2008)}
\addcontentsline{toc}{section}{Question 146: Time Response Analysis (GATE EC 2008)}
\noindent \textbf{\color{primary}MCQ} \hfill \textbf{\color{secondary}1 Mark}


\noindent Step responses of a set of three second-order underdamped systems all have the same percentage overshoot. Which of the following diagrams represents the poles of the three systems ?

\begin{center}\includegraphics[max width=0.85\linewidth]{images/img\_80.jpg}\end{center}

\begin{center}\includegraphics[max width=0.85\linewidth]{images/img\_80.jpg}\end{center}


\begin{enumerate}[label=(\Alph*)]
    \item A
    \item B
    \item \textbf{(Correct)} C
    \item D
\end{enumerate}

\noindent \textbf{Correct Answer:} (C)

\begin{tcolorbox}[solbox, title=Detailed Solution -- Question 146]
Peak overshoot depends on $\xi$ as

\[
\begin{aligned} M_{\rho} &=e^{-\xi \pi / \sqrt{1-\xi^{2}}} \\ \text{where}\quad\xi &=\cos ^{-1} \theta \end{aligned}
\]

 where $\theta$ is the angle made by pole from negative real axis. 
 To make $M_{P}$ same, $\theta$ should be the same.
\end{tcolorbox}

\vspace{1.5em}
\hrule
\vspace{1.5em}

\section*{Question 147: State Space Analysis (GATE EC 2007)}
\addcontentsline{toc}{section}{Question 147: State Space Analysis (GATE EC 2007)}
\noindent \textbf{\color{primary}MCQ} \hfill \textbf{\color{secondary}1 Mark}


\noindent Consider a linear system whose state space representation is $\dot{x}(t)=Ax(t)$. If the initial state vector of the system is 
\[
x(0)=\begin{bmatrix} 1\\ -2 \end{bmatrix}
\]
, then the system response is  
\[
x(t)=\begin{bmatrix} e^{-2t}\\ -2e^{-2t} \end{bmatrix}
\]

 If the initial state vector of the system changes to 
\[
x(0)=\begin{bmatrix} 1\\ -1 \end{bmatrix}
\]
, then the system response becomes  
\[
x(t)=\begin{bmatrix} e^{-t}\\ -e^{-t} \end{bmatrix}
\]
 

The system matrix A is


\begin{enumerate}[label=(\Alph*)]
    \item \[
\begin{bmatrix} 0 & 1\\ -1&1 \end{bmatrix}
\]
    \item \[
\begin{bmatrix} 1 & 1\\ -1&-2 \end{bmatrix}
\]
    \item \[
\begin{bmatrix} 2 & 1\\ -1&-1 \end{bmatrix}
\]
    \item \textbf{(Correct)} \[
\begin{bmatrix} 0 & 1\\ -2&-3 \end{bmatrix}
\]
\end{enumerate}

\noindent \textbf{Correct Answer:} (D)

\begin{tcolorbox}[solbox, title=Detailed Solution -- Question 147]
Multiplication of the eigen value= determinant of
the matrix
$\therefore$  from options it seems determinant should be $\pm  2.$  Only option (d) satisfies as det = 2.
\end{tcolorbox}

\vspace{1.5em}
\hrule
\vspace{1.5em}

\section*{Question 148: State Space Analysis (GATE EC 2007)}
\addcontentsline{toc}{section}{Question 148: State Space Analysis (GATE EC 2007)}
\noindent \textbf{\color{primary}MCQ} \hfill \textbf{\color{secondary}1 Mark}


\noindent Consider a linear system whose state space representation is $x(t)=Ax(t)$. If the initial state vector of the system is 
\[
\dot{x}(0)=\begin{bmatrix} 1\\ -2 \end{bmatrix}
\]
, then the system response is  
\[
x(t)=\begin{bmatrix} e^{-2t}\\ -2e^{-2t} \end{bmatrix}
\]
 
If the initial state vector of the system changes to 
\[
x(0)=\begin{bmatrix} 1\\ -1 \end{bmatrix}
\]
, then the system response becomes  
\[
x(t)=\begin{bmatrix} e^{-t}\\ -e^{-t} \end{bmatrix}
\]
. 

The eigenvalue and eigenvector pairs ($\lambda _{i},v_{i}$) for the system are


\begin{enumerate}[label=(\Alph*)]
    \item \textbf{(Correct)} \[
(-1, \begin{bmatrix} 1\\ -1 \end{bmatrix})\; and \;(-2,\begin{bmatrix} 1\\ -2 \end{bmatrix})
\]
    \item \[
(-2, \begin{bmatrix} 1\\ -1 \end{bmatrix})\; and \;(-1,\begin{bmatrix} 1\\ -2 \end{bmatrix})
\]
    \item \[
(-1, \begin{bmatrix} 1\\ -1 \end{bmatrix})\; and \; (2, \begin{bmatrix} 1\\ -2 \end{bmatrix})
\]
    \item \[
(-2, \begin{bmatrix} 1\\ -1 \end{bmatrix})\; and \;(1,\begin{bmatrix} 1\\ -2 \end{bmatrix})
\]
\end{enumerate}

\noindent \textbf{Correct Answer:} (A)

\begin{tcolorbox}[solbox, title=Detailed Solution -- Question 148]
Sum of the eigen value = Trace of the principle
diagonal matrix

Sum= -3. Only option (A) satisfies both conditions
\end{tcolorbox}

\vspace{1.5em}
\hrule
\vspace{1.5em}

\section*{Question 149: State Space Analysis (GATE EC 2007)}
\addcontentsline{toc}{section}{Question 149: State Space Analysis (GATE EC 2007)}
\noindent \textbf{\color{primary}MCQ} \hfill \textbf{\color{secondary}1 Mark}


\noindent The state space representation of a separately excited DC servo motor dynamics
is given as 

\[
\begin{bmatrix} \frac{d\omega }{dt}\\ \frac{di_{a}}{dt} \end{bmatrix}=\begin{bmatrix} -1 & 1\\ -1&-10 \end{bmatrix}\begin{bmatrix} \omega \\ i_{a} \end{bmatrix}+\begin{bmatrix} 0\\ 10 \end{bmatrix}u
\]

  where $\omega$ is the speed of the motor, $i_{a}$ is the armature current and u is the armature voltage. The transfer function $\frac{\omega (s)}{U(s)}$ of the motor is


\begin{enumerate}[label=(\Alph*)]
    \item \textbf{(Correct)} $\frac{10}{s^{2}+11s+11}$
    \item $\frac{1}{s^{2}+11s+11}$
    \item $\frac{10s+10}{s^{2}+11s+11}$
    \item $\frac{1}{s^{2}+s+11}$
\end{enumerate}

\noindent \textbf{Correct Answer:} (A)

\begin{tcolorbox}[solbox, title=Detailed Solution -- Question 149]
\[
\begin{aligned} \left[\begin{array}{c}\frac{d \omega}{d t} \\\frac{d i_{a}}{d t}\end{array}\right]&=\left[\begin{array}{cc} -1 & 1 \\ -1 & -10 \end{array}\right]\left[\begin{array}{l} \omega \\ i_{a} \end{array}\right]+\left[\begin{array}{c} 0 \\ 10 \end{array}\right] u\\ \Rightarrow \frac{d \omega}{d t}&=-\omega+i_{\mathrm{a}} &\ldots(i)\\ \Rightarrow \frac{d i_{\mathrm{a}}}{d t}&=-\omega-10 i_{\mathrm{a}}+10 u &\ldots(ii) \end{aligned}
\]

 Taking Laplace transform (i) \& (ii) 

\[
\begin{array}{l} \Rightarrow s \omega(s)=-\omega(s)+I_{a}(s) \\ \Rightarrow(s+1) \omega(s)=I_{a}(s) \quad\ldots(iii)\\ \Rightarrow s I_{e}(s)=-\omega(s)-10 I_{a}(s)+10 U(s) \\ \Rightarrow \omega(s)=(-10-s) I_{a}(s)+10 U(s) \\ =(-10-s)(s+1) \omega(s)+10 U(s) \\ =-\left[s^{2}+11 s+10\right] \omega(s)+10 U(s) \\ \Rightarrow\left[s^{2}+11 s+11\right] \omega(s)=10 U(s) \\ \Rightarrow \frac{\omega(s)}{U(s)}=\frac{10}{\left(s^{2}+11 s+11\right)} \end{array}
\]
\end{tcolorbox}

\vspace{1.5em}
\hrule
\vspace{1.5em}

\section*{Question 150: Frequency Response Analysis (GATE EC 2007)}
\addcontentsline{toc}{section}{Question 150: Frequency Response Analysis (GATE EC 2007)}
\noindent \textbf{\color{primary}MCQ} \hfill \textbf{\color{secondary}1 Mark}


\noindent The asymptotic Bode plot of a transfer function is as shown in the figure. The
transfer function G(s) corresponding to this Bode plot is 

\begin{center}\includegraphics[max width=0.85\linewidth]{images/img\_97.jpg}\end{center}

\begin{center}\includegraphics[max width=0.85\linewidth]{images/img\_97.jpg}\end{center}


\begin{enumerate}[label=(\Alph*)]
    \item $\frac{1}{(s+1)(s+20)}$
    \item $\frac{1}{s(s+1)(s+20)}$
    \item $\frac{100}{s(s+1)(s+20)}$
    \item \textbf{(Correct)} $\frac{100}{s(s+1)(1+0.05s)}$
\end{enumerate}

\noindent \textbf{Correct Answer:} (D)

\begin{tcolorbox}[solbox, title=Detailed Solution -- Question 150]
$G(s)=\frac{K}{s(s+1)(s+20)}=\frac{K / 20}{s(1+s)\left(1+\frac{s}{20}\right)}$
 Bode plot is in (1+sT) form 

\[
\begin{aligned} -20 \log \omega+20 \log K&=|60 d B|_{\omega=0.1} \\ -20 \log 0.1+20 \log K&=60 \mathrm{dB} \\ 20 \log K&=40 \mathrm{dB} \\ \Rightarrow \quad K&=100 \\ \therefore \quad G(s)&=\frac{100}{s(s+1)(1+.05 s)} \end{aligned}
\]
\end{tcolorbox}

\vspace{1.5em}
\hrule
\vspace{1.5em}

\section*{Question 151: Root Locus Technique (GATE EC 2007)}
\addcontentsline{toc}{section}{Question 151: Root Locus Technique (GATE EC 2007)}
\noindent \textbf{\color{primary}MCQ} \hfill \textbf{\color{secondary}1 Mark}


\noindent A unity feedback control system has an open-loop transfer function $G(s)=\frac{K}{s(s^{2}+7s+12)}$
The gain K for which $s=-1+j1$  will lie on the root locus of this system is


\begin{enumerate}[label=(\Alph*)]
    \item 4
    \item 5.5
    \item 6.5
    \item \textbf{(Correct)} 10
\end{enumerate}

\noindent \textbf{Correct Answer:} (D)

\begin{tcolorbox}[solbox, title=Detailed Solution -- Question 151]
\[
\begin{aligned} T \cdot F \cdot &=\frac{G(s)}{1+G(s)} \\ A s H(s) &=1 \end{aligned}
\]

 For the point s=-1+j 1 to lie on root locus 

\[
\begin{array} {l} 1+G(s)=0 \\ \Rightarrow 1+\frac{K}{s\left(s^{2}+7 s+12\right)}=0 \\ \text{Putting }\left.s^{2}+7 s+12\right)+K=0 \\ \Rightarrow \quad(-1+j)(1-2 j-1-7+7 j+12)+K=0 \\ \Rightarrow \quad K=+10 \end{array}
\]
\end{tcolorbox}

\vspace{1.5em}
\hrule
\vspace{1.5em}

\section*{Question 152: Compensators \& Controllers (GATE EC 2007)}
\addcontentsline{toc}{section}{Question 152: Compensators \& Controllers (GATE EC 2007)}
\noindent \textbf{\color{primary}MCQ} \hfill \textbf{\color{secondary}1 Mark}


\noindent The open-loop transfer function of a plant is given as $G(s)=\frac{1}{s^{2}-1}$ . If the plant is operated in a unity feedback configuration, then the lead compensator that an stabilize this control system is


\begin{enumerate}[label=(\Alph*)]
    \item \textbf{(Correct)} $\frac{10(s-1)}{s+2}$
    \item $\frac{10(s+4)}{s+2}$
    \item $\frac{10(s+2)}{s+10}$
    \item $\frac{2(s+2)}{s+10}$
\end{enumerate}

\noindent \textbf{Correct Answer:} (A)

\begin{tcolorbox}[solbox, title=Detailed Solution -- Question 152]
Lead compensator is required for stability.

Controllers given in option (A) and (B) are not lead
compensator. From option (C) and (d), with
option (D) system becomes unstable, hence
option (C) is correct.
\end{tcolorbox}

\vspace{1.5em}
\hrule
\vspace{1.5em}

\section*{Question 153: Time Response Analysis (GATE EC 2007)}
\addcontentsline{toc}{section}{Question 153: Time Response Analysis (GATE EC 2007)}
\noindent \textbf{\color{primary}MCQ} \hfill \textbf{\color{secondary}1 Mark}


\noindent The transfer function of a plant is $T(s)=\frac{5}{(s+5)(s^{2}+s+1)}$ The second-order approximation of T(s)  using dominant pole concept is


\begin{enumerate}[label=(\Alph*)]
    \item $\frac{1}{(s+5)(s+1)}$
    \item $\frac{5}{(s+5)(s+1)}$
    \item $\frac{5}{s^{2}+s+1}$
    \item \textbf{(Correct)} $\frac{1}{s^{2}+s+1}$
\end{enumerate}

\noindent \textbf{Correct Answer:} (D)

\begin{tcolorbox}[solbox, title=Detailed Solution -- Question 153]
In dominant pole concept, the factor that has to 
 be eliminated should be in time constant form.

\[
\begin{aligned} \frac{5}{(s+5)\left(s^{2}+s+1\right)}=& \frac{5}{5\left(1+\frac{s}{5}\right)\left(s^{2}+s+1\right)} \\ =& \frac{1}{s^{2}+s+1} \end{aligned}
\]
\end{tcolorbox}

\vspace{1.5em}
\hrule
\vspace{1.5em}

\section*{Question 154: Compensators \& Controllers (GATE EC 2007)}
\addcontentsline{toc}{section}{Question 154: Compensators \& Controllers (GATE EC 2007)}
\noindent \textbf{\color{primary}MCQ} \hfill \textbf{\color{secondary}1 Mark}


\noindent A control system with PD controller is shown in the figure. If the velocity error
constant $K_{V}$=1000 and the damping ratio $\varsigma$=0.5, then the value of $K_{P}$ and $K_{D}$ are 

\begin{center}\includegraphics[max width=0.85\linewidth]{images/img\_23.jpg}\end{center}

\begin{center}\includegraphics[max width=0.85\linewidth]{images/img\_23.jpg}\end{center}


\begin{enumerate}[label=(\Alph*)]
    \item $K_{P}=100,K_{D}=0.09$
    \item \textbf{(Correct)} $K_{P}=100,K_{D}=0.9$
    \item $K_{P}=10,K_{D}=0.09$
    \item $K_{P}=10,K_{D}=0.9$
\end{enumerate}

\noindent \textbf{Correct Answer:} (B)

\begin{tcolorbox}[solbox, title=Detailed Solution -- Question 154]
\[
\begin{aligned} K_{v} &=\lim _{s \rightarrow 0} s G(s) H(s) \\ 1000 &=\lim _{s \rightarrow 0} s \times \frac{\left(K_{P}+K_{D} s\right) 100}{s(s+10)} \\ \Rightarrow \quad K_{P} &=100 \end{aligned}
\]

 Now characteristics eq. $(1+G(s) H(s))=0$

\[
\begin{aligned} 1+\frac{\left(K_{P}+K_{D}\right) 100}{s(s+10)}&=0 \\ \text { Putting } \quad K_{P}&=100 \\ s^{2}+10 s+10^{4}+100 K_{0} s&=0 \\ s^{2}+\left(10+100 K_{0}\right) s+10^{4}&=0 \end{aligned}
\]

 Comparing with standard second order eq.

\[
\begin{array}{l} \text{i.e.}\quad s^{2}+2 \xi \omega_{n} s+\omega_{n}^{2}=0\\ \text{So,} \quad \omega_{n}=100 ; 2 \xi \omega_{n}=10+100 K_{0} \\ \text{Given,}\quad \xi=0.5 \\ 2 \times 0.5 \times 100=10+100 K_{0} \\ K_{D}=0.9 \end{array}
\]
\end{tcolorbox}

\vspace{1.5em}
\hrule
\vspace{1.5em}

\section*{Question 155: Frequency Response Analysis (GATE EC 2007)}
\addcontentsline{toc}{section}{Question 155: Frequency Response Analysis (GATE EC 2007)}
\noindent \textbf{\color{primary}MCQ} \hfill \textbf{\color{secondary}1 Mark}


\noindent The frequency response of a linear, time-invariant system is given by $H(f)=\frac{5}{1+j10 \pi f}$. The step response of the system is


\begin{enumerate}[label=(\Alph*)]
    \item $5(1-e^{-5t})u(t)$
    \item \textbf{(Correct)} $5(1-e^{-\frac{t}{5}})u(t)$
    \item $1/5(1-e^{-5t})u(t)$
    \item $1/5(1-e^{-t/5})u(t)$
\end{enumerate}

\noindent \textbf{Correct Answer:} (B)

\begin{tcolorbox}[solbox, title=Detailed Solution -- Question 155]
\[
\begin{aligned} H(f) &=\frac{5}{1+j 10 \pi f} \\ H(s) &=\frac{5}{1+5 s} \\ &=\frac{5}{5\left(s+\frac{1}{5}\right)}=\frac{1}{s+\frac{1}{5}} \\ \text { Step response } &=\frac{1}{S} \frac{1}{\left(s+\frac{1}{5}\right)} \\ \frac{1}{s} \times \frac{1}{\left(s+\frac{1}{5}\right)} &=\frac{A}{s}+\frac{B}{s+\frac{1}{5}} \\ \Rightarrow A\left(s+\frac{1}{5}\right)&+B s=1 \end{aligned}
\]

 when $s=0, A=5$

\[
\begin{aligned} \text{when }\quad s&=\frac{-1}{5}, B=-5 \\ Y(s) &=\frac{5}{s}-\frac{5}{s+\frac{1}{5}} \\ \Rightarrow y(t)&=5\left[1-e^{-t / 5}\right] u(t) \end{aligned}
\]
\end{tcolorbox}

\vspace{1.5em}
\hrule
\vspace{1.5em}

\section*{Question 156: Stability Analysis \& Routh-Hurwitz (GATE EC 2007)}
\addcontentsline{toc}{section}{Question 156: Stability Analysis \& Routh-Hurwitz (GATE EC 2007)}
\noindent \textbf{\color{primary}MCQ} \hfill \textbf{\color{secondary}1 Mark}


\noindent If the closed-loop transfer function of a control system is given as $T(s)=\frac{s-5}{(s+2)(s+3)}$, then it is


\begin{enumerate}[label=(\Alph*)]
    \item an unstable system
    \item an uncontrollable system
    \item a minimum phase system
    \item \textbf{(Correct)} a non-minimum phase system
\end{enumerate}

\noindent \textbf{Correct Answer:} (D)

\begin{tcolorbox}[solbox, title=Detailed Solution -- Question 156]
As there is a right half zero, the system is a non-minimum phase system.
\end{tcolorbox}

\vspace{1.5em}
\hrule
\vspace{1.5em}

\section*{Question 157: Stability Analysis \& Routh-Hurwitz (GATE EC 2006)}
\addcontentsline{toc}{section}{Question 157: Stability Analysis \& Routh-Hurwitz (GATE EC 2006)}
\noindent \textbf{\color{primary}MCQ} \hfill \textbf{\color{secondary}1 Mark}


\noindent Consider a unity-gain feedback control system whose open-loop transfer function is: 
$G(s)=\frac{as+1}{s^{2}}$ 

With the value of "a" set for a phase - margin of $\frac{\pi }{4}$, the value of unit-impulse response of the open-loop system at t=1 second is equal to


\begin{enumerate}[label=(\Alph*)]
    \item 3.4
    \item 2.4
    \item \textbf{(Correct)} 1.84
    \item 1.74
\end{enumerate}

\noindent \textbf{Correct Answer:} (C)

\begin{tcolorbox}[solbox, title=Detailed Solution -- Question 157]
\[
\begin{aligned} G(s)&=\frac{0.84 s+1}{s^{2}}\\ H(s) &=1, R(s)=1 \\ c(s) &=G(s) \cdot R(s) \\ \therefore \quad C(s) &=\frac{0.84 s+1}{s^{2}} \\ c(t) &=L^{-1}\left[\frac{1+0.84 s}{s^{2}}\right] \\ &=L^{-1}\left[\frac{1}{s^{2}}+\frac{0.84}{s}\right] \\ c(t) &=[t+0.84] U(t) \\ \text{At}\quad t &=1, c(t)=1+0.84=1.84 \end{aligned}
\]
\end{tcolorbox}

\vspace{1.5em}
\hrule
\vspace{1.5em}

\section*{Question 158: Stability Analysis \& Routh-Hurwitz (GATE EC 2006)}
\addcontentsline{toc}{section}{Question 158: Stability Analysis \& Routh-Hurwitz (GATE EC 2006)}
\noindent \textbf{\color{primary}MCQ} \hfill \textbf{\color{secondary}1 Mark}


\noindent Consider a unity-gain feedback control system whose open-loop transfer function is: 
$G(s)=\frac{as+1}{s^{2}}$ 

The value of "a" so that the system has a phase-margin equal to $\frac{\pi }{4}$ is approximately equal to


\begin{enumerate}[label=(\Alph*)]
    \item 2.4
    \item 1.4
    \item \textbf{(Correct)} 0.84
    \item 0.74
\end{enumerate}

\noindent \textbf{Correct Answer:} (C)

\begin{tcolorbox}[solbox, title=Detailed Solution -- Question 158]
\[
\begin{aligned} \mathrm{PM}&=\frac{\pi}{4} \\ \Rightarrow 180+\tan ^{-1} a \omega-180^{\circ}&=\frac{\pi}{4} \\ \tan \frac{\pi}{4}&=\mathrm{a} \omega \\ \Rightarrow\quad a \omega&=1 \end{aligned}
\]

 Now for gain crossover frequency,

\[
\begin{aligned} |G(s)|&=1 \\ \Rightarrow \quad \frac{\sqrt{1+a^{2} \omega^{2}}}{\omega^{2}}&=1 \\ \sqrt{1+1}&=\omega^{2} \quad(\text { as } a \omega=1) \\ f^{2}&=\sqrt{2} \\ \omega&=(2)^{1 / 4} \\ a&=\frac{1}{2^{1 / 4}}=0.84 \end{aligned}
\]
\end{tcolorbox}

\vspace{1.5em}
\hrule
\vspace{1.5em}

\section*{Question 159: State Space Analysis (GATE EC 2006)}
\addcontentsline{toc}{section}{Question 159: State Space Analysis (GATE EC 2006)}
\noindent \textbf{\color{primary}MCQ} \hfill \textbf{\color{secondary}1 Mark}


\noindent A linear system is described by the following state equation 

\[
\dot{X}(t)=AX(t)+BU(t),A=\begin{bmatrix} 0 & 1\\ -1&0 \end{bmatrix}
\]

 The state transition matrix of the system is


\begin{enumerate}[label=(\Alph*)]
    \item \textbf{(Correct)} \[
\begin{bmatrix} cost & sint\\ -sint & cost \end{bmatrix}
\]
    \item \[
\begin{bmatrix} -cost & sint\\ -sint & -cost \end{bmatrix}
\]
    \item \[
\begin{bmatrix} -cost & -sint\\ -sint & cost \end{bmatrix}
\]
    \item \[
\begin{bmatrix} cost & -sint\\ cost & sint \end{bmatrix}
\]
\end{enumerate}

\noindent \textbf{Correct Answer:} (A)

\begin{tcolorbox}[solbox, title=Detailed Solution -- Question 159]
\[
\begin{array}{l} \phi(t)=L^{-1}[S I-A]^{-1} \\ =L^{-1}\left[\left[\begin{array}{ll} s & 0 \\ 0 & s \end{array}\right]-\left[\begin{array}{cc} 0 & 1 \\ -1 & 0 \end{array}\right]\right]^{-1}=L^{-1}\left[\begin{array}{cc} s & -1 \\ 1 & s \end{array}\right]^{-1} \\ =L^{-1}\left[\begin{array}{cc} \frac{s}{s^{2}+1} & \frac{1}{s^{2}+1} \\ \frac{-1}{s^{2}+1} & \frac{s}{s^{2}+1} \end{array}\right]\left[\begin{array}{c} D=s^{2}+1 \\ \end{array}\right] \\ =\left[\begin{array}{cc} \cos t & \sin t \\ -\sin t & \cos t \end{array}\right] \end{array}
\]
\end{tcolorbox}

\vspace{1.5em}
\hrule
\vspace{1.5em}

\section*{Question 160: Compensators \& Controllers (GATE EC 2006)}
\addcontentsline{toc}{section}{Question 160: Compensators \& Controllers (GATE EC 2006)}
\noindent \textbf{\color{primary}MCQ} \hfill \textbf{\color{secondary}1 Mark}


\noindent The transfer function of a phase lead compensator is given by  
$G_{c}(s)=\frac{1+3Ts}{1+Ts}$ where $T > 0$ 
The maximum phase shift provide by such a compensator is


\begin{enumerate}[label=(\Alph*)]
    \item $\frac{\pi }{2}$
    \item $\frac{\pi }{3}$
    \item $\frac{\pi }{4}$
    \item \textbf{(Correct)} $\frac{\pi }{6}$
\end{enumerate}

\noindent \textbf{Correct Answer:} (D)

\begin{tcolorbox}[solbox, title=Detailed Solution -- Question 160]
Maximum phase shift

\[
\begin{aligned} \phi_{m} &=\angle G_{\theta}(s) \\ \phi &=\tan ^{-1} 3 \omega T-\tan ^{-1} \omega T \end{aligned}
\]

 For maximum phase shift

\[
\begin{aligned} \frac{d \phi}{d \omega} &=0 \\ \Rightarrow \frac{3 T}{1+(3 T \omega)^{2}} &=\frac{T}{1+(T \omega)^{2}} \\ 3\left[1+(T \omega)^{2}\right] &=1+(3 T \omega)^{2} \\ 3+3 T^{2} \omega^{2} &=1+9 T^{2} \omega^{2} \\ 2 &=6(\omega T)^{2} \\ (\omega T)^{2} &=\frac{1}{3} \\ \omega T=& \frac{1}{\sqrt{3}} \\\phi_{\max } &=\tan ^{-1} 3 \times \frac{1}{\sqrt{3}}-\tan ^{-1} \frac{1}{\sqrt{3}} \\
&=\frac{\pi}{3}-\frac{\pi}{6}=\frac{\pi}{6}\end{aligned}
\]
\end{tcolorbox}

\vspace{1.5em}
\hrule
\vspace{1.5em}

\section*{Question 161: Time Response Analysis (GATE EC 2006)}
\addcontentsline{toc}{section}{Question 161: Time Response Analysis (GATE EC 2006)}
\noindent \textbf{\color{primary}MCQ} \hfill \textbf{\color{secondary}1 Mark}


\noindent The unit impulse response of a system is $h(t) =  e^{-t},t\geq 0$. For this system the steady-state value of the output for unit step input is equal to


\begin{enumerate}[label=(\Alph*)]
    \item -1
    \item 0
    \item \textbf{(Correct)} 1
    \item $\infty$
\end{enumerate}

\noindent \textbf{Correct Answer:} (C)

\begin{tcolorbox}[solbox, title=Detailed Solution -- Question 161]
\[
\begin{aligned} h(t) &=e^{-t} \\ H(s) &=\frac{1}{s+1} \\ R(s) &=\frac{1}{s} \\ \text { Output } &=H(s) \cdot R(s)=\frac{1}{(s+1)} \cdot \frac{1}{s} \\ \frac{1}{s(s+1)} &=\frac{A}{s}+\frac{B}{s+1} \\ 1 &=A(s+1)+B s \\ s &=0,1=A \\ s &=-1,1=-B \\ \therefore \quad \text{ Output }&=\frac{1}{s}-\frac{1}{s+1} \\ & =\left(1-e^{-t}\right) u(t) \end{aligned}
\]

 When $t=\infty$ at steady state. 
 Output = 1
\end{tcolorbox}

\vspace{1.5em}
\hrule
\vspace{1.5em}

\section*{Question 162: Stability Analysis \& Routh-Hurwitz (GATE EC 2006)}
\addcontentsline{toc}{section}{Question 162: Stability Analysis \& Routh-Hurwitz (GATE EC 2006)}
\noindent \textbf{\color{primary}MCQ} \hfill \textbf{\color{secondary}1 Mark}


\noindent The positive values of "K" and "a"  so that the system shown in the figures below
oscillates at a frequency of 2 rad/sec respectively are 

\begin{center}\includegraphics[max width=0.85\linewidth]{images/img\_119.jpg}\end{center}

\begin{center}\includegraphics[max width=0.85\linewidth]{images/img\_119.jpg}\end{center}


\begin{enumerate}[label=(\Alph*)]
    \item 1, 0.75
    \item \textbf{(Correct)} 2, 0.75
    \item 1, 1
    \item 2, 2
\end{enumerate}

\noindent \textbf{Correct Answer:} (B)

\begin{tcolorbox}[solbox, title=Detailed Solution -- Question 162]
\[
\begin{array}{l} \text { 1+ } G(s) H(s)=1+\frac{K(s+1)}{s^{3}+a s^{2}+2 s+1}=0 \\ \Rightarrow \frac{s^{3}+a s^{2}+2 s+1+K s+K}{s^{3}+a s^{2}+2 s+1}=0 \end{array}
\]

\[
\begin{array}{c|cc} s^{3}& 1 &2+K\\ s^{2}&a & K+1\\ s&\frac{a(2+K)-(K+1)}{a} \end{array}
\]

\[
\begin{array}{l} \text { For oscillation, } \frac{a(2+K)-(K+1)}{a}=0 \\ \qquad \begin{array}{l} a=\frac{K+1}{K+2} \\ a s^{2}+K+1=0 \\ s=j \omega, s^{2}=-\omega^{2}=-4 \\ \Rightarrow-4 a+K+1=0 \\ a=\frac{K+1}{4} \Rightarrow \frac{K+1}{4}=\frac{K+1}{K+2} \Rightarrow K=2 \\ a=0.75 \end{array} \end{array}
\]
\end{tcolorbox}

\vspace{1.5em}
\hrule
\vspace{1.5em}

\section*{Question 163: Frequency Response Analysis (GATE EC 2006)}
\addcontentsline{toc}{section}{Question 163: Frequency Response Analysis (GATE EC 2006)}
\noindent \textbf{\color{primary}MCQ} \hfill \textbf{\color{secondary}1 Mark}


\noindent The Nyquist plot of $G(j\omega )$ $H(j\omega )$ for a closed loop control system, passes through
(- 1, j0) point in the GH plane. The gain margin of the system in dB is equal to


\begin{enumerate}[label=(\Alph*)]
    \item infinite
    \item greater than zero
    \item less than zero
    \item \textbf{(Correct)} zero
\end{enumerate}

\noindent \textbf{Correct Answer:} (D)

\begin{tcolorbox}[solbox, title=Detailed Solution -- Question 163]
\[
\begin{array}{l} \mathrm{G.M}=20 \log \frac{1}{a} \mathrm{dB} \\ a=1 \quad \therefore \mathrm{G} \cdot \mathrm{M}=0 \end{array}
\]
\end{tcolorbox}

\vspace{1.5em}
\hrule
\vspace{1.5em}

\section*{Question 164: Time Response Analysis (GATE EC 2006)}
\addcontentsline{toc}{section}{Question 164: Time Response Analysis (GATE EC 2006)}
\noindent \textbf{\color{primary}MCQ} \hfill \textbf{\color{secondary}1 Mark}


\noindent The unit step response of a system starting from rest is given by 
 $c(t)=1-e^{-2t} \; for \;  t\geq 0$ 
The transfer function of the system is


\begin{enumerate}[label=(\Alph*)]
    \item $\frac{1}{1+2s}$
    \item \textbf{(Correct)} $\frac{2}{2+s}$
    \item $\frac{1}{2+s}$
    \item $\frac{2s}{1+2s}$
\end{enumerate}

\noindent \textbf{Correct Answer:} (B)

\begin{tcolorbox}[solbox, title=Detailed Solution -- Question 164]
\[
\begin{aligned} C(s)&=\frac{1}{s}-\frac{1}{s+2}=\frac{2}{s(s+2)} \\ R(s)&=\frac{1}{s} \quad H(s)=\frac{C(s)}{R(s)}=\frac{2}{s+2} \end{aligned}
\]
\end{tcolorbox}

\vspace{1.5em}
\hrule
\vspace{1.5em}

\section*{Question 165: Frequency Response Analysis (GATE EC 2006)}
\addcontentsline{toc}{section}{Question 165: Frequency Response Analysis (GATE EC 2006)}
\noindent \textbf{\color{primary}MCQ} \hfill \textbf{\color{secondary}1 Mark}


\noindent In the system shown below, x(t) = (sint)u(t). In steady-state, the response y(t)
will be 

\begin{center}\includegraphics[max width=0.85\linewidth]{images/img\_114.jpg}\end{center}

\begin{center}\includegraphics[max width=0.85\linewidth]{images/img\_114.jpg}\end{center}


\begin{enumerate}[label=(\Alph*)]
    \item \textbf{(Correct)} $\frac{1}{\sqrt{2}}sin(t-\frac{\pi }{4})$
    \item $\frac{1}{\sqrt{2}}sin(t+\frac{\pi }{4})$
    \item $\frac{1}{\sqrt{2}}e^{-t}sint$
    \item $sint-cost$
\end{enumerate}

\noindent \textbf{Correct Answer:} (A)

\begin{tcolorbox}[solbox, title=Detailed Solution -- Question 165]
\[
\begin{aligned} y(t) &=x(t) * h(t) \\ Y(s) &=X(s) \cdot H(s) \\ H(j \omega) &=\frac{1}{s+1}=\frac{1}{\sqrt{2}} \angle-45^{\circ} \\ &=\frac{1}{\sqrt{2}} \angle-\frac{\pi}{4} \\ x(t) &=\sin (t) u(t) \\ \therefore \quad y(t) &=\frac{1}{\sqrt{2}} \sin \left(t-\frac{\pi}{4}\right) \end{aligned}
\]
\end{tcolorbox}

\vspace{1.5em}
\hrule
\vspace{1.5em}

\section*{Question 166: Frequency Response Analysis (GATE EC 2006)}
\addcontentsline{toc}{section}{Question 166: Frequency Response Analysis (GATE EC 2006)}
\noindent \textbf{\color{primary}MCQ} \hfill \textbf{\color{secondary}1 Mark}


\noindent The open-loop function of a unity-gain feedback control system is given by  
$G(s)=\frac{K}{(s+1)(s+2)}$ 
 The gain margin of the system in dB is given by


\begin{enumerate}[label=(\Alph*)]
    \item 0
    \item 1
    \item 20
    \item \textbf{(Correct)} $\infty$
\end{enumerate}

\noindent \textbf{Correct Answer:} (D)

\begin{tcolorbox}[solbox, title=Detailed Solution -- Question 166]
For 2nd order system G.M. $=\infty$
\end{tcolorbox}

\vspace{1.5em}
\hrule
\vspace{1.5em}

\section*{Question 167: Frequency Response Analysis (GATE EC 2005)}
\addcontentsline{toc}{section}{Question 167: Frequency Response Analysis (GATE EC 2005)}
\noindent \textbf{\color{primary}MCQ} \hfill \textbf{\color{secondary}1 Mark}


\noindent The open loop transfer function of a unity feedback is given by $G(s)=\frac{3e^{-2s}}{s(s+2)}$. 
  The gain and phase margins of the system will be


\begin{enumerate}[label=(\Alph*)]
    \item -7.09 dB  and  87.5$^{\circ}$
    \item 7.09 dB and 87.5$^{\circ}$
    \item 7.09 dB and -87.5$^{\circ}$
    \item \textbf{(Correct)} -7.09dB and  -87.5$^{\circ}$
\end{enumerate}

\noindent \textbf{Correct Answer:} (D)

\begin{tcolorbox}[solbox, title=Detailed Solution -- Question 167]
\[
\begin{aligned} \text{G.M. at }\omega_{\phi}&=\frac{3}{0.632\left(0.632^{2}+4\right)^{1 / 2}} \\ a &=2.26 \\ G . M &=2010 g \frac{1}{a} \\ \Rightarrow 20 \log \frac{1}{2.26} &=-7.09 \end{aligned}
\]

 since, G.M. is negative system is unstable. 
$\therefore$ P.M. is also negative. 
$\mathrm{P} \mathrm{M} .=-87.5^{\circ}$
\end{tcolorbox}

\vspace{1.5em}
\hrule
\vspace{1.5em}

\section*{Question 168: Frequency Response Analysis (GATE EC 2005)}
\addcontentsline{toc}{section}{Question 168: Frequency Response Analysis (GATE EC 2005)}
\noindent \textbf{\color{primary}MCQ} \hfill \textbf{\color{secondary}1 Mark}


\noindent The open loop transfer function of a unity feedback is given by $G(s)=\frac{3e^{-2s}}{s(s+2)}$. 
 The gain and phase crossover frequencies in rad/sec are, respectively


\begin{enumerate}[label=(\Alph*)]
    \item 0.632 and 1.26
    \item 0.632 and 0.485
    \item 0.485 and 0.632
    \item \textbf{(Correct)} 1.26 and 0.632
\end{enumerate}

\noindent \textbf{Correct Answer:} (D)

\begin{tcolorbox}[solbox, title=Detailed Solution -- Question 168]
Gain cross over frequency where gain is 1. 

\[
\begin{aligned} |G(s)| &=1 \\ \Rightarrow \frac{3}{\omega\left(\omega^{2}+4\right)^{12}} &=1 \\ \Rightarrow \frac{9}{\omega^{2}\left(\omega^{2}+4\right)} &=1 \\ \Rightarrow \omega_{g}&=1.26 \end{aligned}
\]

 Phase cross over frequency 

\[
\begin{aligned} \text{Where }\angle G H&=-180^{\circ} \\ \Rightarrow \quad \omega_{p c}&=0.632\left(\frac{\omega_{\phi}}{2}=\cot 2 \omega_{\phi}\right) \end{aligned}
\]
\end{tcolorbox}

\vspace{1.5em}
\hrule
\vspace{1.5em}

\section*{Question 169: Root Locus Technique (GATE EC 2005)}
\addcontentsline{toc}{section}{Question 169: Root Locus Technique (GATE EC 2005)}
\noindent \textbf{\color{primary}MCQ} \hfill \textbf{\color{secondary}1 Mark}


\noindent An unity feedback system is given as
 $G(s)=\frac{K(1-s)}{s(s+3)}$
 Indicate the correct root locus diagram.

\begin{center}\includegraphics[max width=0.85\linewidth]{images/img\_107.jpg}\end{center}

\begin{center}\includegraphics[max width=0.85\linewidth]{images/img\_107.jpg}\end{center}


\begin{enumerate}[label=(\Alph*)]
    \item \textbf{(Correct)} A
    \item B
    \item C
    \item D
\end{enumerate}

\noindent \textbf{Correct Answer:} (A)

\begin{tcolorbox}[solbox, title=Detailed Solution -- Question 169]
\[
\begin{aligned} 1+G(s) H(s) &=0 \\ K &=\frac{s^{2}+3 s}{1-s} \end{aligned}
\]

 For break away and break in point

\[
\begin{array}{l} \quad \frac{d K}{d s}=(1-s)(2 s+3)+s^{2}+3 s=0 \\ =-s^{2}+2 s+3=0 \\ \Rightarrow s^{2}-2 s-3=0 \\ \Rightarrow(s-3)(s+1)=0 \\ s=3,-1 \end{array}
\]

 -1 is the breakaway point and 3 is the break in point.
\end{tcolorbox}

\vspace{1.5em}
\hrule
\vspace{1.5em}

\section*{Question 170: Compensators \& Controllers (GATE EC 2005)}
\addcontentsline{toc}{section}{Question 170: Compensators \& Controllers (GATE EC 2005)}
\noindent \textbf{\color{primary}MCQ} \hfill \textbf{\color{secondary}1 Mark}


\noindent A double integrator plant G(s) = K/$s^{2}$, H(s) = 1 is to be compensated to
achieve the damping ratio 
\[
\begin{aligned}
\\\zeta
\end{aligned}
\]
=0.5, and an undamped natural frequency, $\omega _{n}$=5 rad/sec.  Which one of the following compensator $G_{c}(s)$ will be suitable ?


\begin{enumerate}[label=(\Alph*)]
    \item \textbf{(Correct)} $\frac{s+3}{s+9.9}$
    \item $\frac{s+9.9}{s+3}$
    \item $\frac{s-6}{s+8.33}$
    \item $\frac{s+6}{s}$
\end{enumerate}

\noindent \textbf{Correct Answer:} (A)

\begin{tcolorbox}[solbox, title=Detailed Solution -- Question 170]
\begin{center}\includegraphics[max width=0.85\linewidth]{images/img\_17.jpg}\end{center}

\begin{center}\includegraphics[max width=0.85\linewidth]{images/img\_17.jpg}\end{center}

\[
\begin{aligned} & \xi=0.5 \\ \cos ^{-1} 0.5 &=60^{\circ} \\ \therefore \quad \theta &=60^{\circ}\\ \angle G(s) &=.\frac{K}{s^{2}}|_{s=-2.5+4.33 j} \\ &=-2 \tan ^{-1} \frac{4.33}{-2.5} ; 120^{\circ} \end{aligned}
\]

 $\therefore$ For compensated system
$\angle=180-120; \; 60^{\circ}$
 (b) \& (d) are lag network \& for compensating lag network  $K /s^{2}$, a lead network is required
$\therefore$ Putting $s=-2.5+4.33 j$ in (a) gives

\[
\frac{K(s+3)}{s^{2}(s+9.9)}=\frac{0.5+j 4.33}{7.4+j 4.33}=53^{\circ} ; 60^{\circ}
\]

$\therefore$ Option (A) is the correct answer.
\end{tcolorbox}

\vspace{1.5em}
\hrule
\vspace{1.5em}

\section*{Question 171: Time Response Analysis (GATE EC 2005)}
\addcontentsline{toc}{section}{Question 171: Time Response Analysis (GATE EC 2005)}
\noindent \textbf{\color{primary}MCQ} \hfill \textbf{\color{secondary}1 Mark}


\noindent A ramp input applied to an unity feedback system results in 5% steady state
error. The type number and zero frequency gain of the system are respectively


\begin{enumerate}[label=(\Alph*)]
    \item \textbf{(Correct)} 1 and 20
    \item 0 and 20
    \item $0 \; and\; \frac{1}{20}$
    \item $1 \; and \; \frac{1}{20}$
\end{enumerate}

\noindent \textbf{Correct Answer:} (A)

\begin{tcolorbox}[solbox, title=Detailed Solution -- Question 171]
\[
\begin{aligned} e_{s s} &=\lim _{s \rightarrow 0} s E(s), R(s)=\frac{1}{s^{2}} \\ &=\lim _{s \rightarrow 0} s \frac{R(s)}{1+G(s)} \\ &=\lim _{s \rightarrow 0} \frac{1}{s+s G(s)}=\text { finite (given) } \\ K_{v} &=\lim _{s \rightarrow 0} s G(s) \\ e_{s s} &=\lim _{s \rightarrow 0} \frac{1}{s G(s)}=\text { finite }=5 \%=\frac{1}{20} \\ \therefore \quad K &=20 \end{aligned}
\]

$K_{v}$ is finite for type 1 system having ramp input.
\end{tcolorbox}

\vspace{1.5em}
\hrule
\vspace{1.5em}

\section*{Question 172: Time Response Analysis (GATE EC 2005)}
\addcontentsline{toc}{section}{Question 172: Time Response Analysis (GATE EC 2005)}
\noindent \textbf{\color{primary}MCQ} \hfill \textbf{\color{secondary}1 Mark}


\noindent In the derivation of expression for peak percent overshoot, 
 $M_{P}=exp(\frac{-\pi \xi }{\sqrt{1-\xi ^{2}}}) \times 100$% 
Which one of the following conditions is NOT required ?


\begin{enumerate}[label=(\Alph*)]
    \item System is linear and time invariant
    \item The system transfer function has a pair of complex conjugate poles and no zeroes.
    \item \textbf{(Correct)} There is no transportation delay in the system
    \item The system has zero initial conditions.
\end{enumerate}

\noindent \textbf{Correct Answer:} (C)


\vspace{1.5em}
\hrule
\vspace{1.5em}

\section*{Question 173: Frequency Response Analysis (GATE EC 2005)}
\addcontentsline{toc}{section}{Question 173: Frequency Response Analysis (GATE EC 2005)}
\noindent \textbf{\color{primary}MCQ} \hfill \textbf{\color{secondary}1 Mark}


\noindent The polar diagram of a conditionally stable system for open loop gain K = 1 is
shown in the figure. The open loop transfer function of the system is known to
be stable. The closed loop system is stable for 

\begin{center}\includegraphics[max width=0.85\linewidth]{images/img\_128.jpg}\end{center}

\begin{center}\includegraphics[max width=0.85\linewidth]{images/img\_128.jpg}\end{center}


\begin{enumerate}[label=(\Alph*)]
    \item $K < 5   \; and \;  \frac{1}{2} < K < \frac{1}{8}$
    \item \textbf{(Correct)} $K < \frac{1}{8}   \; and \;  \frac{1}{2} < K < 5$
    \item $K < \frac{1}{8}   \; and \;  5 < K$
    \item $K>\frac{1}{8}   \; and \;  5>K$
\end{enumerate}

\noindent \textbf{Correct Answer:} (B)

\begin{tcolorbox}[solbox, title=Detailed Solution -- Question 173]
\begin{center}\includegraphics[max width=0.85\linewidth]{images/img\_132.jpg}\end{center}

\begin{center}\includegraphics[max width=0.85\linewidth]{images/img\_132.jpg}\end{center}

System is stable in region -0.2 to -2 \& on the left
side of -8 as number of encirclement there is zero.

\[
\begin{aligned}
0.2 K < 1 &\Rightarrow K < 5 \\
2 K > 1 &\Rightarrow K >  0.5 \\
\therefore 0.5 < K& < 5\\
1 > 8 &K
\end{aligned}
\]

$K < \frac{1}{8}$ 
 (negative sign only shows that it is on negative axis)
\end{tcolorbox}

\vspace{1.5em}
\hrule
\vspace{1.5em}

\section*{Question 174: Time Response Analysis (GATE EC 2005)}
\addcontentsline{toc}{section}{Question 174: Time Response Analysis (GATE EC 2005)}
\noindent \textbf{\color{primary}MCQ} \hfill \textbf{\color{secondary}1 Mark}


\noindent Despite the presence of negative feedback, control systems still have problems of instability because the


\begin{enumerate}[label=(\Alph*)]
    \item \textbf{(Correct)} Components used have non- linearities
    \item Dynamic equations of the subsystem are not known exactly.
    \item Mathematical analysis involves approximations.
    \item System has large negative phase angle at high frequencies.
\end{enumerate}

\noindent \textbf{Correct Answer:} (A)


\vspace{1.5em}
\hrule
\vspace{1.5em}

\section*{Question 175: Frequency Response Analysis (GATE EC 2005)}
\addcontentsline{toc}{section}{Question 175: Frequency Response Analysis (GATE EC 2005)}
\noindent \textbf{\color{primary}MCQ} \hfill \textbf{\color{secondary}1 Mark}


\noindent Which one of the following polar diagrams corresponds to a lag network ?

\begin{center}\includegraphics[max width=0.85\linewidth]{images/img\_106.jpg}\end{center}

\begin{center}\includegraphics[max width=0.85\linewidth]{images/img\_106.jpg}\end{center}


\begin{enumerate}[label=(\Alph*)]
    \item A
    \item B
    \item C
    \item \textbf{(Correct)} D
\end{enumerate}

\noindent \textbf{Correct Answer:} (D)

\begin{tcolorbox}[solbox, title=Detailed Solution -- Question 175]
Let $\frac{1}{s+1}$ be a lag network
 At $\omega=0, \text{Mag}=\infty, \angle=-\tan ^{-1} 0=0$
$\omega=\infty, \mathrm{Mag}=0, \angle=-\tan ^{-1} \infty=-90^{\circ}$
 If in the direction of $\omega$ increasing phase shift is decreasing system is lag network.
\end{tcolorbox}

\vspace{1.5em}
\hrule
\vspace{1.5em}

\section*{Question 176: State Space Analysis (GATE EC 2005)}
\addcontentsline{toc}{section}{Question 176: State Space Analysis (GATE EC 2005)}
\noindent \textbf{\color{primary}MCQ} \hfill \textbf{\color{secondary}1 Mark}


\noindent A linear system is equivalently represented by two sets of state equations : 
$\dot{X}=AX+BU \; and \;  W=CW+DU$ 
The eigenvalues of the representations are also computed as $[\lambda]$ and $[\mu]$. Which one of the following statements is true ?


\begin{enumerate}[label=(\Alph*)]
    \item $[\lambda ]=[\mu ]  \; and \;  X=W$
    \item \textbf{(Correct)} $[\lambda ]=[\mu ]  \; and \;  X\neq W$
    \item $[\lambda ]\neq [\mu ]  \; and \;  X= W$
    \item $[\lambda ]= [\mu ]  \; and \;  X\neq W$
\end{enumerate}

\noindent \textbf{Correct Answer:} (B)

\begin{tcolorbox}[solbox, title=Detailed Solution -- Question 176]
Eigen vlues of $A=[\lambda]$

Eigen values of  $W=[\mu]$

The eigen values of a system are always unique. 
$\mathrm{So},[\lambda]=[\mu]$

But a system can be represented by different state
models having different set of state variables.

\[
\begin{array}{l}
X=W \\
X \neq W
\end{array}
\]

Both are possible conditions.
\end{tcolorbox}

\vspace{1.5em}
\hrule
\vspace{1.5em}

\section*{Question 177: State Space Analysis (GATE EC 2004)}
\addcontentsline{toc}{section}{Question 177: State Space Analysis (GATE EC 2004)}
\noindent \textbf{\color{primary}MCQ} \hfill \textbf{\color{secondary}1 Mark}


\noindent Given 
\[
A=\begin{bmatrix} 1 & 0\\ 0&1 \end{bmatrix}
\]
, the state transition matrix $e^{At}$ is given by


\begin{enumerate}[label=(\Alph*)]
    \item \[
\begin{bmatrix} 0 &e^{-t} \\ e^{-t}&0 \end{bmatrix}
\]
    \item \textbf{(Correct)} \[
\begin{bmatrix} e^{t} &0 \\ 0 & e^{t} \end{bmatrix}
\]
    \item \[
\begin{bmatrix} e^{-t} &0 \\ 0 & e^{-t} \end{bmatrix}
\]
    \item \[
\begin{bmatrix} 0 &e^{t} \\ e^{t}&0 \end{bmatrix}
\]
\end{enumerate}

\noindent \textbf{Correct Answer:} (B)

\begin{tcolorbox}[solbox, title=Detailed Solution -- Question 177]
\[
\begin{aligned} [s I-A] &=\left[\begin{array}{ll}s & 0 \\0 & s\end{array}\right]-\left[\begin{array}{ll}1 & 0 \\0 & 1\end{array}\right] \\ s I-A] & =\left[\begin{array}{cc}s-1 & 0 \\0 & s-1\end{array}\right] \\ e^{A t} & =[s I-A]^{-1}\\ &=\left[\begin{array}{cc}-\frac{1}{s-1} & 0 \\ 0 & \frac{1}{s-1}\end{array}\right]=\left[\begin{array}{ll}e^{t} & 0 \\ 0 & e^{t}\end{array}\right] \end{aligned}
\]
\end{tcolorbox}

\vspace{1.5em}
\hrule
\vspace{1.5em}

\section*{Question 178: State Space Analysis (GATE EC 2004)}
\addcontentsline{toc}{section}{Question 178: State Space Analysis (GATE EC 2004)}
\noindent \textbf{\color{primary}MCQ} \hfill \textbf{\color{secondary}1 Mark}


\noindent The state variable equations of a system are : 
 1.  $\dot{x_{1}}=-3x_{1}-x_{2}+u$

  2. $\dot{x_{2}}=2x_{1}$

  $y=x_{1}+u$

 The system is


\begin{enumerate}[label=(\Alph*)]
    \item controllable but not observable
    \item observable but not controllable
    \item neither controllable nor observable
    \item \textbf{(Correct)} controllable and observable
\end{enumerate}

\noindent \textbf{Correct Answer:} (D)

\begin{tcolorbox}[solbox, title=Detailed Solution -- Question 178]
\[
\begin{aligned} \left[\begin{array}{c}\dot{x}_{1} \\\dot{x}_{2}\end{array}\right]&=\left[\begin{array}{cc}-3 & -1 \\2 & 0\end{array}\right]\left[\begin{array}{l}x_{1} \\x_{2}\end{array}\right]+\left[\begin{array}{l}1 \\0\end{array}\right] u \\ y&=\left[\begin{array}{ll}1 & 0\end{array}\right]\left[\begin{array}{l}x_{1} \\x_{2}\end{array}\right]+\left[\begin{array}{l}1 \\0\end{array}\right] u \\ Q_{C}&=\left[\begin{array}{ll}B & A B\end{array}\right]\\ \end{aligned}
\]

\[
\begin{aligned} A B&=\left[\begin{array}{c}-3 \\ 2\end{array}\right] \\ \therefore \quad a_{c}&=\left[\begin{array}{cc}1 & -3 \\ 0 & 2\end{array}\right] \neq 0 \\ \therefore \text{Controllable } \\ Q_{0}&=\left[C^{\top} \quad A^{\top} C^{\prime}\right] \\ C^{T}&=\left[\begin{array}{l}1 \\ 0\end{array}\right], \quad A^{T}&=\left[\begin{array}{ll}-3 & 2 \\ -1 & 0\end{array}\right] \\ A^{\top} C^{T}&=\left[\begin{array}{l}-3 \\ -1\end{array}\right] \\ \therefore \quad Q_{0}&=\left[\begin{array}{ll}1 & -3 \\ 0 & -1\end{array}\right] \neq 0 \\ \therefore \text{ Observable} \end{aligned}
\]
\end{tcolorbox}

\vspace{1.5em}
\hrule
\vspace{1.5em}

\section*{Question 179: Stability Analysis \& Routh-Hurwitz (GATE EC 2004)}
\addcontentsline{toc}{section}{Question 179: Stability Analysis \& Routh-Hurwitz (GATE EC 2004)}
\noindent \textbf{\color{primary}MCQ} \hfill \textbf{\color{secondary}1 Mark}


\noindent For the polynomial $P(s)=s^{5}+s^{4}+2s^{3}+2s^{2}+3s+15$, the number of roots
which lie in the right half of the s -plane is


\begin{enumerate}[label=(\Alph*)]
    \item 4
    \item \textbf{(Correct)} 2
    \item 3
    \item 1
\end{enumerate}

\noindent \textbf{Correct Answer:} (B)

\begin{tcolorbox}[solbox, title=Detailed Solution -- Question 179]
$P(s)=s^{5}+s^{4}+2 s^{3}+2 s^{2}+3 s+15$

\[
\begin{array}{c|ccc} s^{5}& 1 & 2 & 3 \\ s^{4} & 1 & 2 & 15 \\ s^{3} & 0(\epsilon) & -12 & 0 \\ s^{2} & \frac{2 \epsilon+12}{\epsilon} & 15 & 0 \\ s&\frac{\frac{\frac{-12(2 \epsilon+12)}{\epsilon}}{\epsilon }-15 \epsilon }{\frac{\epsilon}{\left(\frac{2 \epsilon+12}{\epsilon}\right)}}&& \end{array}
\]

$\frac{2 \epsilon+12}{\epsilon}=t$
 Let $\in$ be a small positive no. 
$\frac{-12 t-15 \epsilon}{t}$
$s \Rightarrow-12-\frac{15 \epsilon}{t}$
$s ^{0}\quad 15$
$\therefore$ Two sign change from $s^{2}$ to s and s to $s^{0}$
$\therefore$ 2 roots on RHS of s plane.
\end{tcolorbox}

\vspace{1.5em}
\hrule
\vspace{1.5em}

\section*{Question 180: Stability Analysis \& Routh-Hurwitz (GATE EC 2004)}
\addcontentsline{toc}{section}{Question 180: Stability Analysis \& Routh-Hurwitz (GATE EC 2004)}
\noindent \textbf{\color{primary}MCQ} \hfill \textbf{\color{secondary}1 Mark}


\noindent The open-loop transfer function of a unity feedback system is $G(s)=\frac{K}{s(s^{2}+s+2)(s+3)}$. The range of K for which the system is stable is


\begin{enumerate}[label=(\Alph*)]
    \item \textbf{(Correct)} $\frac{21}{4} > K > 0$
    \item $13 > K > 0$
    \item $\frac{21}{4} < K < \infty$
    \item $-6 < K < \infty$
\end{enumerate}

\noindent \textbf{Correct Answer:} (A)

\begin{tcolorbox}[solbox, title=Detailed Solution -- Question 180]
\[
\begin{array}{l} G(s)=\frac{K}{s\left(s^{2}+s+2\right)(s+3)} \\ H(s)=1\\ 1+G(s) H(s)=1+\frac{K}{s\left(s^{3}+3 s^{2}+s^{2}+3 s+2 s+6\right)} \\ =\frac{s^{4}+4 s^{3}+5 s^{2}+6 s+K}{s\left(s^{3}+4 s^{2}+5 s+6\right)} \\ 1+G(s) H(s)=0 \end{array}
\]

\[
\begin{array}{l|ll}s^{4} & 1 & 5\;K \\ s^{3} & 4 & 6 \\ s^{2} & \frac{7}{2} & k \\ s & \frac{\frac{7}{2} \times 6-4 K}{\frac{7}{2}} & 0 \\ s^{0} & k & \frac{7}{2}\end{array}
\]

 For system to be stable, $K > 0$ 
$(21-4 K) \cdot \frac{2}{7}> 0$
$\frac{21}{4} > K \Rightarrow K< \frac{21}{4}$
$\therefore \quad \frac{21}{4} > K > 0$
\end{tcolorbox}

\vspace{1.5em}
\hrule
\vspace{1.5em}

\section*{Question 181: Basics, Block Diagrams \& SFGs (GATE EC 2004)}
\addcontentsline{toc}{section}{Question 181: Basics, Block Diagrams \& SFGs (GATE EC 2004)}
\noindent \textbf{\color{primary}MCQ} \hfill \textbf{\color{secondary}1 Mark}


\noindent Consider the signal flow graph shown in Fig. The gain $\frac{X_{5}}{X_{1}}$ is

\begin{center}\includegraphics[max width=0.85\linewidth]{images/img\_29.jpg}\end{center}

\begin{center}\includegraphics[max width=0.85\linewidth]{images/img\_29.jpg}\end{center}


\begin{enumerate}[label=(\Alph*)]
    \item $\frac{1-(be+cf+dg)}{abcd}$
    \item $\frac{bedg}{1-(be+cf+dg)}$
    \item \textbf{(Correct)} $\frac{abcd}{1-(be+cf+dg)+bedg}$
    \item $\frac{1-(be+cf+dg)+bedg}{abcd}$
\end{enumerate}

\noindent \textbf{Correct Answer:} (C)

\begin{tcolorbox}[solbox, title=Detailed Solution -- Question 181]
\[
\begin{array}{l} P_{1}=a b c d, \quad \Delta_{1}=1 \\ L_{1}=b e, L_{2}=c f, L_{3}=d g \\ \text { Non-touching loops are } L_{1} \& L_{3}=b e d g \\ \therefore \quad \frac{X_{5}}{X_{1}}=\frac{a b c d}{1-(b e+c f+d g)+b e d g} \end{array}
\]
\end{tcolorbox}

\vspace{1.5em}
\hrule
\vspace{1.5em}

\section*{Question 182: Frequency Response Analysis (GATE EC 2004)}
\addcontentsline{toc}{section}{Question 182: Frequency Response Analysis (GATE EC 2004)}
\noindent \textbf{\color{primary}MCQ} \hfill \textbf{\color{secondary}1 Mark}


\noindent A system has poles at 0.1 Hz, 1 Hz and 80 Hz; zeros at 5 Hz, 100 Hz and 200
Hz. The approximate phase of the system response at 20 Hz is


\begin{enumerate}[label=(\Alph*)]
    \item \textbf{(Correct)} $-90^{\circ}$
    \item $0^{\circ}$
    \item $90^{\circ}$
    \item $-180^{\circ}$
\end{enumerate}

\noindent \textbf{Correct Answer:} (A)

\begin{tcolorbox}[solbox, title=Detailed Solution -- Question 182]
Pole at 0.01 and 1 Hz gives $-180^{\circ}$ phase.

Zero at 5 Hz gives $+90^{\circ}$ phase
$\therefore \quad$at $20Hz-90^{\circ}$ phase shift is provided.
\end{tcolorbox}

\vspace{1.5em}
\hrule
\vspace{1.5em}

\section*{Question 183: Time Response Analysis (GATE EC 2004)}
\addcontentsline{toc}{section}{Question 183: Time Response Analysis (GATE EC 2004)}
\noindent \textbf{\color{primary}MCQ} \hfill \textbf{\color{secondary}1 Mark}


\noindent A causal system having the transfer function H(s)=1/(s+2) is excited with
10u(t). The time at which the output reaches 99% of its steady state value is


\begin{enumerate}[label=(\Alph*)]
    \item 2.7 sec
    \item 2.5 sec
    \item \textbf{(Correct)} 2.3 sec
    \item 2.1 sec
\end{enumerate}

\noindent \textbf{Correct Answer:} (C)

\begin{tcolorbox}[solbox, title=Detailed Solution -- Question 183]
\[
\begin{aligned} H(s) &=\frac{1}{s+2} \\ r(t) &=10 u(t) \\ R(s) &=\frac{10}{s} \\ C(s) &=H(s) \cdot R(s)=\frac{1}{s+2} \cdot \frac{10}{s} \\ \Rightarrow \quad \frac{10}{s(s+2)} &=\frac{A}{s}+\frac{B}{s+2} \\ 10 &=A(s+2)+B s \\ s &=0,10=2 A \\ \Rightarrow \quad A &=5 \\ s &=-2,10=-2 B \\ \Rightarrow \quad B &=-5 \\ \therefore \quad C(s) &=\frac{5}{s}-\frac{5}{s+2} \\ c(t) &=5\left[1-e^{2 t}\right] u(t) \end{aligned}
\]

 Steady state value when t=0 is 5.99% of steady
 state value reaches at

\[
\begin{aligned} & 5\left[1-e^{2 t}\right]=0.99 \times 5 \\ \Rightarrow \quad & 1-e^{-2 t}=0.99 \\ & e^{2 t}=0.1 \\ \Rightarrow \quad &-2 t=\ln 0.1 \\ \Rightarrow \quad & t=2.3 \mathrm{sec} \end{aligned}
\]
\end{tcolorbox}

\vspace{1.5em}
\hrule
\vspace{1.5em}

\section*{Question 184: Time Response Analysis (GATE EC 2004)}
\addcontentsline{toc}{section}{Question 184: Time Response Analysis (GATE EC 2004)}
\noindent \textbf{\color{primary}MCQ} \hfill \textbf{\color{secondary}1 Mark}


\noindent A system described by the following differential equation $\frac{d^{2}y}{dt^{2}}+3\frac{dy}{dt}+2y=x(t)$ is initially at rest. For input x(t)=2u(t), the output y(t) is


\begin{enumerate}[label=(\Alph*)]
    \item \textbf{(Correct)} $(1-2e^{-t}+e^{-2t})u(t)$
    \item $(1+2e^{-t}-2e^{-2t})u(t)$
    \item $(0.5+2e^{-t}+1.5e^{-2t})u(t)$
    \item $(0.5+2e^{-t}+2e^{-2t})u(t)$
\end{enumerate}

\noindent \textbf{Correct Answer:} (A)

\begin{tcolorbox}[solbox, title=Detailed Solution -- Question 184]
\[
\begin{aligned} \frac{d^{2} y}{d t^{2}}+\frac{3 d y}{d t}+2 y &=X(t) \\ s^{2} Y(s)+3 s Y(s)+2 Y(s) &=X(s) \\ x(t) &=2 u(t) \\ X(s) &=\frac{2}{s} \\ \therefore \quad\left(s^{2}+3 s+2\right) Y(s) &=\frac{2}{s} \\ r(s) &=\frac{2}{s(s+2)(s+1)} \\ \frac{2}{s(s+2)(s+1)} &=\frac{A}{s}+\frac{B}{s+2}+\frac{C}{s+1} \end{aligned}
\]

\[
\begin{aligned} 2&=A(s+2)(s+1)+B s(s+1)+C(s+2) s \\ s&=0,2=2 A \Rightarrow A=1 \\ s&=-1,2=-C \Rightarrow C=-2 \\ s&=-2,2=2 B \Rightarrow B=1 \\ Y(s)&=\frac{1}{s}+\frac{1}{s+2}-\frac{2}{s+1} \\ y(t)&=(1+e^{-2 t}-2 e^{-t}) u(t)  \end{aligned}
\]
\end{tcolorbox}

\vspace{1.5em}
\hrule
\vspace{1.5em}

\section*{Question 185: Frequency Response Analysis (GATE EC 2004)}
\addcontentsline{toc}{section}{Question 185: Frequency Response Analysis (GATE EC 2004)}
\noindent \textbf{\color{primary}MCQ} \hfill \textbf{\color{secondary}1 Mark}


\noindent Consider the Bode magnitude plot shown in the fig. The transfer function H(s) is

\begin{center}\includegraphics[max width=0.85\linewidth]{images/img\_39.jpg}\end{center}

\begin{center}\includegraphics[max width=0.85\linewidth]{images/img\_39.jpg}\end{center}


\begin{enumerate}[label=(\Alph*)]
    \item $\frac{(s+10)}{(s+1)(s+100)}$
    \item $\frac{10(s+10)}{(s+10)(s+100)}$
    \item \textbf{(Correct)} $\frac{10^{2}(s+1)}{(s+10)(s+100)}$
    \item $\frac{10^{3}(s+100)}{(s+1)(s+10)}$
\end{enumerate}

\noindent \textbf{Correct Answer:} (C)

\begin{tcolorbox}[solbox, title=Detailed Solution -- Question 185]
\[
\begin{aligned} 20 \log K &=-20 \mathrm{dB} \\ \Rightarrow \quad K &=10^{-1}=0.1 \\ \text { Zero at } \omega=1 & \& \text { poles at } \omega=10,100 \\ H(s) &=\frac{K(s+1)}{\left(\frac{s}{10}+1\right)\left(\frac{s}{100}+1\right)} \\ &=\frac{0.1(s+1) \times 10 \times 100}{(s+10)(s+100)} \\ &=\frac{10^2 (s+1)}{(s+10)(s+100)} \end{aligned}
\]
\end{tcolorbox}

\vspace{1.5em}
\hrule
\vspace{1.5em}

\section*{Question 186: Root Locus Technique (GATE EC 2004)}
\addcontentsline{toc}{section}{Question 186: Root Locus Technique (GATE EC 2004)}
\noindent \textbf{\color{primary}MCQ} \hfill \textbf{\color{secondary}1 Mark}


\noindent Given $G(s)H(s)=\frac{K}{s(s+1)(s+3)}$,  the point of intersection of the asymptotes of
the root loci with the real axis is


\begin{enumerate}[label=(\Alph*)]
    \item -4
    \item 1.33
    \item \textbf{(Correct)} -1.33
    \item 4
\end{enumerate}

\noindent \textbf{Correct Answer:} (C)

\begin{tcolorbox}[solbox, title=Detailed Solution -- Question 186]
$\text { Centroid }=\frac{\Sigma P-\Sigma Z}{P-Z}=\frac{-1-3}{3}=-1.33$
\end{tcolorbox}

\vspace{1.5em}
\hrule
\vspace{1.5em}

\section*{Question 187: Stability Analysis \& Routh-Hurwitz (GATE EC 2004)}
\addcontentsline{toc}{section}{Question 187: Stability Analysis \& Routh-Hurwitz (GATE EC 2004)}
\noindent \textbf{\color{primary}MCQ} \hfill \textbf{\color{secondary}1 Mark}


\noindent The gain margin for the system with open-loop transfer function $G(s)H(s)=\frac{2(1+s)}{s^{2}}$, is


\begin{enumerate}[label=(\Alph*)]
    \item \textbf{(Correct)} $\infty$
    \item 0
    \item 1
    \item $-\infty$
\end{enumerate}

\noindent \textbf{Correct Answer:} (A)

\begin{tcolorbox}[solbox, title=Detailed Solution -- Question 187]
\[
\begin{aligned} \angle G(s) H(s) &=-180^{\circ}+\tan ^{-1} \omega \\ \text{For}\quad\omega_{\theta} &=-180^{\circ}+\tan ^{-1} \omega=-180^{\circ} \\ \omega &=0 \\ |G(s) H(s)| &=\frac{2 \sqrt{1+\omega^{2}}}{\omega^{2}}=\infty \\ G \cdot M \cdot &=\frac{1}{\infty}=0 \\ \ln \quad d B G \cdot M &=\infty \end{aligned}
\]
\end{tcolorbox}

\vspace{1.5em}
\hrule
\vspace{1.5em}

\section*{Question 188: State Space Analysis (GATE EC 2003)}
\addcontentsline{toc}{section}{Question 188: State Space Analysis (GATE EC 2003)}
\noindent \textbf{\color{primary}MCQ} \hfill \textbf{\color{secondary}1 Mark}


\noindent The zero-input response of a system given by the state-space equation 

\[
\begin{bmatrix} \dot{X_1}\\\dot{X_2} \end{bmatrix}=\begin{bmatrix} 1 &0 \\ 1& 1 \end{bmatrix}\begin{bmatrix} X_1\\X_2 \end{bmatrix}
\]
 and 
\[
\begin{bmatrix} X_1(0)\\X_2(0) \end{bmatrix}=\begin{bmatrix} 1\\0 \end{bmatrix}
\]
 is


\begin{enumerate}[label=(\Alph*)]
    \item \[
\begin{bmatrix} te^{t}\\t \end{bmatrix}
\]
    \item \[
\begin{bmatrix} e^{t}\\t \end{bmatrix}
\]
    \item \textbf{(Correct)} \[
\begin{bmatrix} e^{t}\\t e^{t} \end{bmatrix}
\]
    \item \[
\begin{bmatrix} t\\t e^{t} \end{bmatrix}
\]
\end{enumerate}

\noindent \textbf{Correct Answer:} (C)

\begin{tcolorbox}[solbox, title=Detailed Solution -- Question 188]
\[
\begin{aligned} (s I-A) &=\left[\begin{array}{ll} s & 0 \\ 0 & s \end{array}\right]-\left[\begin{array}{ll} 1 & 0 \\ 1 & 1 \end{array}\right] \\ &=\left[\begin{array}{cc} s-1 & 0 \\ -1 & s-1 \end{array}\right] \\ (s I-A)^{-1} &=\frac{\left[\begin{array}{cc} (s-1) & 0 \\ +1 & (s-1) \end{array}\right]}{(s-1)^{2}}\\ &=\left[\begin{array}{cc}\frac{1}{s-1} & 0 \\ \frac{+1}{(s-1)^{2}} & \frac{1}{s-1}\end{array}\right] \\ L^{-1}[s I-A]^{-1}&=e^{A t}=\left[\begin{array}{cc}e^{t} & 0 \\ t e^{t} & e^{t}\end{array}\right] \\ &=\left[\begin{array}{cc}e^{t} & 0 \\ t e^{t} & e^{t}\end{array}\right]\left[\begin{array}{l}1 \\ 0\end{array}\right]=\left[\begin{array}{c}e^{t} \\ t e^{t}\end{array}\right] \end{aligned}
\]
\end{tcolorbox}

\vspace{1.5em}
\hrule
\vspace{1.5em}

\section*{Question 189: Frequency Response Analysis (GATE EC 2003)}
\addcontentsline{toc}{section}{Question 189: Frequency Response Analysis (GATE EC 2003)}
\noindent \textbf{\color{primary}MCQ} \hfill \textbf{\color{secondary}1 Mark}


\noindent The gain margin and the phase margin of feedback system with $G(s)H(s)=\frac{s}{(s+100)^{3}}$ are


\begin{enumerate}[label=(\Alph*)]
    \item $0 dB,0^{\circ}$
    \item \textbf{(Correct)} $\infty ,\infty$
    \item $\infty ,0^{\circ}$
    \item 88.5 dB,$\infty$
\end{enumerate}

\noindent \textbf{Correct Answer:} (B)

\begin{tcolorbox}[solbox, title=Detailed Solution -- Question 189]
\[
\begin{aligned} G(s) H(s)&=\frac{s}{(s+100)^{3}}\\ \text{Calculating } \omega_{\mathrm{pc}}, \\ -180^{\circ}&=+90-3 \tan ^{-1} \frac{\omega_{p c}}{100} \\ -270^{\circ}&=-3 \tan ^{-1} \frac{\omega_{p c}}{100}\\ \text{or, }\quad\omega_{p c}&=\infty\\ \therefore \quad\mathrm{G} \cdot \mathrm{M} \cdot&=\infty \end{aligned}
\]

 Hence cannot be determined. 
 Calculating $\omega_{g c}$
$\text { From, }|G(s) H(s)|=1$
$\omega_{g c}=$ negative value
 G.M. and P.M. of the system cannot be determined.
\end{tcolorbox}

\vspace{1.5em}
\hrule
\vspace{1.5em}

\section*{Question 190: Time Response Analysis (GATE EC 2003)}
\addcontentsline{toc}{section}{Question 190: Time Response Analysis (GATE EC 2003)}
\noindent \textbf{\color{primary}MCQ} \hfill \textbf{\color{secondary}1 Mark}


\noindent A second-order system has the transfer function $\frac{C(s)}{R(s)}=\frac{4}{s^{2}+4s+4}$.
 With r(t) as the unit-step function, the response c(t) of the system is represented by 

\begin{center}\includegraphics[max width=0.85\linewidth]{images/img\_79.jpg}\end{center}

\begin{center}\includegraphics[max width=0.85\linewidth]{images/img\_79.jpg}\end{center}


\begin{enumerate}[label=(\Alph*)]
    \item A
    \item \textbf{(Correct)} B
    \item C
    \item D
\end{enumerate}

\noindent \textbf{Correct Answer:} (B)

\begin{tcolorbox}[solbox, title=Detailed Solution -- Question 190]
\[
\begin{aligned} \frac{C(s)}{R(s)} &=\frac{4}{s^{2}+4 s+4} \\ 2 \xi(0) &=4, \omega_{n}=2 \\ \therefore \quad \xi &=1 \quad \text { (Critical damping) } \\ t_{s} &=\frac{4}{\xi \omega_{n}}=\frac{4}{1 \times 2}=2 \end{aligned}
\]
\end{tcolorbox}

\vspace{1.5em}
\hrule
\vspace{1.5em}

\section*{Question 191: Frequency Response Analysis (GATE EC 2003)}
\addcontentsline{toc}{section}{Question 191: Frequency Response Analysis (GATE EC 2003)}
\noindent \textbf{\color{primary}MCQ} \hfill \textbf{\color{secondary}1 Mark}


\noindent The approximate Bode magnitude plot of a minimum phase system is shown in Fig. below. The transfer function of the system is  

\begin{center}\includegraphics[max width=0.85\linewidth]{images/img\_131.jpg}\end{center}

\begin{center}\includegraphics[max width=0.85\linewidth]{images/img\_131.jpg}\end{center}


\begin{enumerate}[label=(\Alph*)]
    \item \textbf{(Correct)} $10^{8}\frac{(s+0.1)^{3}}{(s+10)^{2}(s+100)}$
    \item $10^{7}\frac{(s+0.1)^{3}}{(s+10)(s+100)}$
    \item $10^{8} \frac{(s+0.1)^{2}}{(s+10)^{2}(s+100)}$
    \item $10^{9}\frac{(s+0.1)^{2}}{(s+10)(s+100)^{2}}$
\end{enumerate}

\noindent \textbf{Correct Answer:} (A)

\begin{tcolorbox}[solbox, title=Detailed Solution -- Question 191]
$\omega=0.1 to 10,+120 \mathrm{dB}$ change
$\therefore$3 zeroes at 0.1
$\omega=10 \text { to } 100,-40 \mathrm{dB}$
 [i.e. +60 to +20] change 
 So two poles at $\omega=10$
$\omega=100,-20 \mathrm{dB}$ change.
$\therefore$ one pole at $\omega=100$
$\therefore T(s)=\frac{K(s+0.1)^{3}}{(s+10)^{2}(s+100)}$
$.20 \log (K / \omega)|_{\omega=10}=140$
$.\frac{K}{\omega}|_{\omega=10}=10^{7}$
$\therefore \quad K=10^{8}$
\end{tcolorbox}

\vspace{1.5em}
\hrule
\vspace{1.5em}

\section*{Question 192: Root Locus Technique (GATE EC 2003)}
\addcontentsline{toc}{section}{Question 192: Root Locus Technique (GATE EC 2003)}
\noindent \textbf{\color{primary}MCQ} \hfill \textbf{\color{secondary}1 Mark}


\noindent The root locus of system $G(s)H(s)=\frac{K}{s(s+2)(s+3)}$ has the break-away point
located at


\begin{enumerate}[label=(\Alph*)]
    \item (- 0.5,0)
    \item (- 2.548,0)
    \item (- 4,0)
    \item \textbf{(Correct)} (- 0.784,0)
\end{enumerate}

\noindent \textbf{Correct Answer:} (D)

\begin{tcolorbox}[solbox, title=Detailed Solution -- Question 192]
\[
\begin{array}{l} 1+G(s) H(s)=0 \\ \Rightarrow 1+\frac{K}{s(s+2)(s+3)}=0 \\ \Rightarrow 1+\frac{K}{s\left(s^{2}+5 s+6\right)}=0 \\ \Rightarrow \quad K=-1 \cdot\left(s^{3}+5 s^{2}+6 s\right) \\ \Rightarrow \frac{d K}{d s}=-\left(3 s^{2}+10 s+6\right) \\ \text{For break away point  } \frac{d K}{d s}=0 \\ \Rightarrow 3 s^{2}+10 s+6=0 \\ s=\frac{-10 \pm \sqrt{100-72}}{6} \\ =\frac{-10 \pm 5.3}{6}=-0.784,-2.55 \\ \end{array}
\]

\begin{center}\includegraphics[max width=0.85\linewidth]{images/img\_56.jpg}\end{center}

\begin{center}\includegraphics[max width=0.85\linewidth]{images/img\_56.jpg}\end{center}

Centroid  $=\frac{-2-3}{3}=\frac{-5}{3}=-1.66$

angle of asymtotes  $=(2 K+1) \frac{\pi}{P-Z}$
$=(2 K+1) \frac{\pi}{3}$
$\frac{\pi}{3}, \pi, \frac{5 \pi}{3}$
$\therefore$ Break away point is (-0.784,0)
as -2.55 does not lie on RL.
\end{tcolorbox}

\vspace{1.5em}
\hrule
\vspace{1.5em}

\section*{Question 193: Basics, Block Diagrams \& SFGs (GATE EC 2003)}
\addcontentsline{toc}{section}{Question 193: Basics, Block Diagrams \& SFGs (GATE EC 2003)}
\noindent \textbf{\color{primary}MCQ} \hfill \textbf{\color{secondary}1 Mark}


\noindent The signal flow graph of a system is shown in Fig. below. The transfer function C(s)/R(s) of the system is  

\begin{center}\includegraphics[max width=0.85\linewidth]{images/img\_21.jpg}\end{center}

\begin{center}\includegraphics[max width=0.85\linewidth]{images/img\_21.jpg}\end{center}


\begin{enumerate}[label=(\Alph*)]
    \item $\frac{6}{s^{2}+29s+6}$
    \item $\frac{6s}{s^{2}+29s+6}$
    \item $\frac{s(s+2)}{s^{2}+29s+6}$
    \item \textbf{(Correct)} $\frac{s(s+27)}{s^{2}+29s+6}$
\end{enumerate}

\noindent \textbf{Correct Answer:} (D)

\begin{tcolorbox}[solbox, title=Detailed Solution -- Question 193]
\[
\begin{aligned} P_{1}&=1 \\ \Delta_{1}&=1+\frac{3}{s}+\frac{24}{s}=\frac{s+27}{s}\\ G(s) &=\frac{P_{1} \Delta_{1}}{1-(\text { loop gain })+\text { pair of non- touching loops }} \\ L_{1}&=\frac{-3}{s}, L_{2} =\frac{-24}{s}, L_{3}=\frac{-2}{s} \end{aligned}
\]

$L_{1}$ and $L_{3}$ are non-touching.

\[
\begin{aligned} \therefore G(s)&=\frac{\frac{(s+27)}{s}}{1-\left(\frac{-3}{s}-\frac{24}{s}-\frac{2}{s}\right)+\frac{-2}{s} \times \frac{-3}{s}} \\ &=\frac{\frac{(s+27)}{s}}{1+\frac{29}{s}+\frac{6}{s^{2}}} \\ \frac{C(s)}{R(s)} &=\frac{s(s+27)}{s^{2}+29 s+6} \end{aligned}
\]
\end{tcolorbox}

\vspace{1.5em}
\hrule
\vspace{1.5em}

\section*{Question 194: Compensators \& Controllers (GATE EC 2003)}
\addcontentsline{toc}{section}{Question 194: Compensators \& Controllers (GATE EC 2003)}
\noindent \textbf{\color{primary}MCQ} \hfill \textbf{\color{secondary}1 Mark}


\noindent A PD controller is used to compensate a system. Compared to the uncompensated system, the compensated system has


\begin{enumerate}[label=(\Alph*)]
    \item a higher type number
    \item reduced damping
    \item \textbf{(Correct)} higher noise amplification
    \item larger transient overshoot
\end{enumerate}

\noindent \textbf{Correct Answer:} (C)

\begin{tcolorbox}[solbox, title=Detailed Solution -- Question 194]
A PD controller does not effect the type of the
system.

It reduces the overshoot and increases the
damping. 

It increases the bandwidth therefor SNR
decreases,
System becomes more prone to noise
\end{tcolorbox}

\vspace{1.5em}
\hrule
\vspace{1.5em}

\section*{Question 195: Frequency Response Analysis (GATE EC 2003)}
\addcontentsline{toc}{section}{Question 195: Frequency Response Analysis (GATE EC 2003)}
\noindent \textbf{\color{primary}MCQ} \hfill \textbf{\color{secondary}1 Mark}


\noindent Fig. shows the Nyquist plot of the open-loop transfer function G(s)H(s) of a system. If G(s)H(s) has one right-hand pole, the closed-loop system is 

\begin{center}\includegraphics[max width=0.85\linewidth]{images/img\_129.jpg}\end{center}

\begin{center}\includegraphics[max width=0.85\linewidth]{images/img\_129.jpg}\end{center}


\begin{enumerate}[label=(\Alph*)]
    \item \textbf{(Correct)} always stable
    \item unstable with one closed-loop right hand pole
    \item unstable with two closed-loop right hand poles
    \item unstable with three closed-loop right hand poles
\end{enumerate}

\noindent \textbf{Correct Answer:} (A)

\begin{tcolorbox}[solbox, title=Detailed Solution -- Question 195]
The encirclement of critical point (-1,0) in A.C
direction is once.

\[
\begin{aligned}
\therefore   \quad N&=1, P=1 &\text{(given)}\\
Z&=P-N=O
\end{aligned}
\]

No zero in RH of s-plane So system is stable.
\end{tcolorbox}

\vspace{1.5em}
\hrule
\vspace{1.5em}

\section*{Question 196: Frequency Response Analysis (GATE EC 2002)}
\addcontentsline{toc}{section}{Question 196: Frequency Response Analysis (GATE EC 2002)}
\noindent \textbf{\color{primary}MCQ} \hfill \textbf{\color{secondary}1 Mark}


\noindent The system with the open loop transfer function $G(s)H(s)=\frac{1}{s(s^{2}+s+1)}$ has a gain margin of


\begin{enumerate}[label=(\Alph*)]
    \item - 6 dB
    \item \textbf{(Correct)} 0dB
    \item 3.5dB
    \item 6dB
\end{enumerate}

\noindent \textbf{Correct Answer:} (B)

\begin{tcolorbox}[solbox, title=Detailed Solution -- Question 196]
$G(s) H(s)=\frac{1}{s\left(s^{2}+s+1\right)}$
$\omega$ when $\angle G(s) H(s)=-180^{\circ}$

\[
\begin{aligned} -180 &=-90-\tan ^{-1} \frac{\omega}{1-\omega^{2}} \\ -90 &=-\tan ^{-1} \frac{\omega}{1-\omega^{2}} \\ 1-\omega^{2} &=0 \\ \omega_{\alpha c} &=1 \mathrm{rad} / \mathrm{sec} \end{aligned}
\]

 Value of gain at $\omega_{\rho c}=1$

\[
\begin{aligned} |G(s) H(s)| &=\frac{1}{\sqrt{((1-\omega^{2})^{2}+|j \omega|^{2}}} \\ \therefore \quad G \cdot M \cdot &=-20 \log 1=0 \end{aligned}
\]
\end{tcolorbox}

\vspace{1.5em}
\hrule
\vspace{1.5em}

\section*{Question 197: Stability Analysis \& Routh-Hurwitz (GATE EC 2002)}
\addcontentsline{toc}{section}{Question 197: Stability Analysis \& Routh-Hurwitz (GATE EC 2002)}
\noindent \textbf{\color{primary}MCQ} \hfill \textbf{\color{secondary}1 Mark}


\noindent The characteristic polynomial of a system is 
$q (s) = 2s^{5}+s^{4}+4s^{3}+2s^{2}+2s+1$. 
The system is


\begin{enumerate}[label=(\Alph*)]
    \item stable
    \item marginally stable
    \item \textbf{(Correct)} unstable
    \item oscillatory
\end{enumerate}

\noindent \textbf{Correct Answer:} (C)

\begin{tcolorbox}[solbox, title=Detailed Solution -- Question 197]
Routh table is

\[
\begin{array}{c|ccc} s^{5} & 2(1) & 4(2) & 2(1) \\ s^{4} & 1 & 2 & 1 \\ s^{3} & 0 & 0 & 0 \\ s^{2} & & & \\ s^{1} & & & \\ s^{0} & & & \end{array}
\]

\[
\begin{aligned} \therefore \frac{d}{d s}\left(s^{4}+2 s^{2}+1\right) &=0 \\ 4 s^{3}+4 s &=0 \\ s &=\pm j, s=\pm j \\ \frac{d}{d s}\left(s^{2}+1\right) &=0 \\ 2 s &=0 \\ \Rightarrow s &=0 \end{aligned}
\]

 Double roots on imaginary axis so system is unstable
\end{tcolorbox}

\vspace{1.5em}
\hrule
\vspace{1.5em}

\section*{Question 198: Time Response Analysis (GATE EC 2002)}
\addcontentsline{toc}{section}{Question 198: Time Response Analysis (GATE EC 2002)}
\noindent \textbf{\color{primary}MCQ} \hfill \textbf{\color{secondary}1 Mark}


\noindent The transfer function of a system is $G(s)=\frac{100}{(s+1)(s+100)}$. For a unit step input to the system the approximate settling time for 2% criterion is


\begin{enumerate}[label=(\Alph*)]
    \item 100sec
    \item \textbf{(Correct)} 4sec
    \item 1sec
    \item 0.01sec
\end{enumerate}

\noindent \textbf{Correct Answer:} (B)

\begin{tcolorbox}[solbox, title=Detailed Solution -- Question 198]
$G(s)=\frac{100}{(s+1)(s+100)}$
 Taking dominant pole consideration, 
$s=-100$ pole is not taken.
$\therefore \quad G(s)=\frac{100}{s+1}$
 Now it is 1 st order system 
$\therefore \quad t_{s}=4 T=4 \times 1=4 \mathrm{s}$
\end{tcolorbox}

\vspace{1.5em}
\hrule
\vspace{1.5em}

\section*{Question 199: Stability Analysis \& Routh-Hurwitz (GATE EC 2002)}
\addcontentsline{toc}{section}{Question 199: Stability Analysis \& Routh-Hurwitz (GATE EC 2002)}
\noindent \textbf{\color{primary}MCQ} \hfill \textbf{\color{secondary}1 Mark}


\noindent The system shown in figure remains stable when 

\begin{center}\includegraphics[max width=0.85\linewidth]{images/img\_98.jpg}\end{center}

\begin{center}\includegraphics[max width=0.85\linewidth]{images/img\_98.jpg}\end{center}


\begin{enumerate}[label=(\Alph*)]
    \item $K < -1$
    \item $- 1< K < 1$
    \item $1< K < 3$
    \item \textbf{(Correct)} $K < -3$
\end{enumerate}

\noindent \textbf{Correct Answer:} (D)

\begin{tcolorbox}[solbox, title=Detailed Solution -- Question 199]
\[
\frac{Y(s)}{R(s)}=\frac{\frac{K}{s}}{1-\left(\frac{3}{s}+\frac{K}{s}\right)}=\frac{K}{s-(3+K)}
\]

 For system to be stable,

\[
\begin{aligned} 3+K& < 0 \\ \Rightarrow \quad K& < -3 \end{aligned}
\]
\end{tcolorbox}

\vspace{1.5em}
\hrule
\vspace{1.5em}

\section*{Question 200: State Space Analysis (GATE EC 2002)}
\addcontentsline{toc}{section}{Question 200: State Space Analysis (GATE EC 2002)}
\noindent \textbf{\color{primary}MCQ} \hfill \textbf{\color{secondary}1 Mark}


\noindent The transfer function Y(s)U(s) of a system described by the state equations
x(t) = -2x(t) + 2u(t) and y(t) = 0.5x(t) is


\begin{enumerate}[label=(\Alph*)]
    \item $\frac{0.5}{(s-2)}$
    \item $\frac{1}{(s-2)}$
    \item $\frac{0.5}{(s+2)}$
    \item \textbf{(Correct)} $\frac{1}{(s+2)}$
\end{enumerate}

\noindent \textbf{Correct Answer:} (D)

\begin{tcolorbox}[solbox, title=Detailed Solution -- Question 200]
\[
\begin{aligned} \dot{x}(t)=-2 x(t)+2 u(t) &\ldots(i)\\ y(t)=0.5 x(t)&\ldots(ii) \end{aligned}
\]

 From (i), Taking Laplace transform of (i)

\[
\begin{aligned} s X(s)&=-2 X(s)+2 U(s)\\ X(s)[s+2]&=2 U(s)\\ \Rightarrow \quad X(s)&=\frac{2 U(s)}{(s+2)} \end{aligned}
\]

 Taking Laplace transform of (ii)

\[
\begin{aligned} Y(s)&=0.5 X(s) \\ Y(s)&=\frac{0.5 \times 2 U(s)}{s+2} \\ \therefore \quad \frac{Y(s)}{U(s)}&=\frac{1}{(s+2)} \end{aligned}
\]
\end{tcolorbox}

\vspace{1.5em}
\hrule
\vspace{1.5em}

\section*{Question 201: Frequency Response Analysis (GATE EC 2002)}
\addcontentsline{toc}{section}{Question 201: Frequency Response Analysis (GATE EC 2002)}
\noindent \textbf{\color{primary}MCQ} \hfill \textbf{\color{secondary}1 Mark}


\noindent The phase margin of a system with the open-loop transfer function $G(s)H(s)=\frac{(1-s)}{(s+1)(s+2)}$ is


\begin{enumerate}[label=(\Alph*)]
    \item $0^{\circ}$
    \item $63.4^{\circ}$
    \item $90^{\circ}$
    \item \textbf{(Correct)} $\infty$
\end{enumerate}

\noindent \textbf{Correct Answer:} (D)

\begin{tcolorbox}[solbox, title=Detailed Solution -- Question 201]
\[
\begin{aligned} \mathrm{PM}&=180^{\circ}+.\angle \mathrm{GH}(j \omega)|_{\omega_{9 c}} \\ \text { For } \omega_{\mathrm{gc}^{\prime}} &|G(s) H(s)|_{\omega=\omega_{\infty}}=1\\ |\frac{1-s}{(1+s)(2+s)}|&=\frac{\sqrt{1+\omega^{2}}}{\sqrt{1+\omega^{2} \sqrt{4+\omega^{2}}}}&=1 \\ \sqrt{4+\omega^{2}}&=1 \\ \Rightarrow \quad 4+\omega^{2}&=1 \\ \omega^{2}&=-3 \text{ (imaginary)}\\ \end{aligned}
\]

 So no gain crossover frequency
$\therefore PM=\infty$
\end{tcolorbox}

\vspace{1.5em}
\hrule
\vspace{1.5em}

\section*{Question 202: Root Locus Technique (GATE EC 2002)}
\addcontentsline{toc}{section}{Question 202: Root Locus Technique (GATE EC 2002)}
\noindent \textbf{\color{primary}MCQ} \hfill \textbf{\color{secondary}1 Mark}


\noindent Which of the following points is NOT on the root locus of a system with the open
loop transfer function $G(s)H(s)=\frac{k}{s(s+1)(s+3)}$


\begin{enumerate}[label=(\Alph*)]
    \item $S=-j\sqrt{3}$
    \item \textbf{(Correct)} $S=-1.5$
    \item $S=-3$
    \item $S=-\infty$
\end{enumerate}

\noindent \textbf{Correct Answer:} (B)

\begin{tcolorbox}[solbox, title=Detailed Solution -- Question 202]
\begin{center}\includegraphics[max width=0.85\linewidth]{images/img\_43.jpg}\end{center}

\begin{center}\includegraphics[max width=0.85\linewidth]{images/img\_43.jpg}\end{center}

RL lies where no. of poles and zeroes to the right
of the pole is odd.
$\therefore \; s = -1.5$ doesn't lie on AL
\end{tcolorbox}

\vspace{1.5em}
\hrule
\vspace{1.5em}

\section*{Question 203: Time Response Analysis (GATE EC 2002)}
\addcontentsline{toc}{section}{Question 203: Time Response Analysis (GATE EC 2002)}
\noindent \textbf{\color{primary}MCQ} \hfill \textbf{\color{secondary}1 Mark}


\noindent Consider a system with the transfer function $G(s)=\frac{s+6}{ks^{2}+s+6}$. Its damping ratio will be 0.5 when the value of k is


\begin{enumerate}[label=(\Alph*)]
    \item $\frac{2}{6}$
    \item 3
    \item \textbf{(Correct)} $\frac{1}{6}$
    \item 6
\end{enumerate}

\noindent \textbf{Correct Answer:} (C)

\begin{tcolorbox}[solbox, title=Detailed Solution -- Question 203]
\[
\begin{array}{l} \frac{s+6}{k\left(s^{2}+\frac{s}{k}+\frac{6}{k}\right)} \\ \text { Comparing with } s^{2}+2 \xi \omega_{n}+\omega_{n}^{2} \\ \begin{aligned} \omega_{n}&=\sqrt{\frac{6}{k}} \\ 2 \xi \omega_{n}&=\frac{1}{k} \\ 2 \times 0.5 \times \sqrt{\frac{6}{k}}&=\frac{1}{k} \\ \Rightarrow \qquad \frac{6}{k}&=\frac{1}{k^{2}} \\ k&=\frac{1}{6} \end{aligned} \end{array}
\]
\end{tcolorbox}

\vspace{1.5em}
\hrule
\vspace{1.5em}

\section*{Question 204: Stability Analysis \& Routh-Hurwitz (GATE EC 2001)}
\addcontentsline{toc}{section}{Question 204: Stability Analysis \& Routh-Hurwitz (GATE EC 2001)}
\noindent \textbf{\color{primary}MCQ} \hfill \textbf{\color{secondary}1 Mark}


\noindent The feedback control system in figure is stable

\begin{center}\includegraphics[max width=0.85\linewidth]{images/img\_141.jpg}\end{center}

\begin{center}\includegraphics[max width=0.85\linewidth]{images/img\_141.jpg}\end{center}


\begin{enumerate}[label=(\Alph*)]
    \item For all $K\geq 0$
    \item Only If $K\geq 1$
    \item \textbf{(Correct)} Only If $0\leq K < 1$
    \item Only If $0\leq K\leq 2$
\end{enumerate}

\noindent \textbf{Correct Answer:} (C)

\begin{tcolorbox}[solbox, title=Detailed Solution -- Question 204]
\[
\begin{aligned} T . F . &=\frac{G_{1} G_{2}}{1+G_{1} G_{2} H} \\ &=\frac{\frac{K(s-2)}{(s+2)^{2}}}{1+\frac{K(s-2)(s-2)}{(s+2)^{2}}} \\ &=\frac{K(s-2)}{(s+2)^{2}+K(s-2)^{2}} \end{aligned}
\]

$\therefore$ Characteristic equation
$=s^{2}+4+4 s+K s^{2}-4 K s+4 K$
$(1+K) s^{2}+K(4-4 K) s+4 K+4=0$
 Routh table is

\[
\begin{array}{l|ll} s^{2} & (1+K) & (4 K+4) \\ s^{1} & K(4-4 K) & 0 \\ s^{0} & 4(K+1) & \end{array}
\]

 For system to be stable,

\[
\begin{aligned} 1-K& > 0\\ 1& > K \text{ and } K \geq 0 \quad \text{ (given in question)} \end{aligned}
\]

 From third row,

\[
\begin{aligned} K& > -1 \\ \therefore \quad 0 &\leq K < 1 \end{aligned}
\]
\end{tcolorbox}

\vspace{1.5em}
\hrule
\vspace{1.5em}

\section*{Question 205: Time Response Analysis (GATE EC 2001)}
\addcontentsline{toc}{section}{Question 205: Time Response Analysis (GATE EC 2001)}
\noindent \textbf{\color{primary}MCQ} \hfill \textbf{\color{secondary}1 Mark}


\noindent The open-loop DC gain of a unity negative feedback system with closed-loop
transfer function $\frac{s+4}{s^{2}+7s+13}$ is


\begin{enumerate}[label=(\Alph*)]
    \item $\frac{4}{13}$
    \item \textbf{(Correct)} $\frac{4}{9}$
    \item 4
    \item 13
\end{enumerate}

\noindent \textbf{Correct Answer:} (B)

\begin{tcolorbox}[solbox, title=Detailed Solution -- Question 205]
\[
\begin{aligned} \mathrm{CLTF} &=\frac{G(s)}{1+G(s) H(s)}=\frac{s+4}{s^{2}+7 s+13} \\ \frac{1+G(s) H(s)}{G(s)} &=\frac{s^{2}+7 s+13}{s+4} \end{aligned}
\]

$H(s)=1$ for unity feedback.

\[
\begin{aligned} \frac{1}{G(s)} &=\frac{s^{2}+7 s+13}{s+4}-1 \\ \frac{1}{G(s)} &=\frac{s^{2}+6 s+9}{s+4} \\ \therefore \quad G(s) &=\frac{s+4}{s^{2}+6 s+9} \\ \text { For D.C. } \quad s&=0\\ \therefore \quad G(s) &=\text{open loop gain} \\ &=\frac{4}{9} \end{aligned}
\]
\end{tcolorbox}

\vspace{1.5em}
\hrule
\vspace{1.5em}

\section*{Question 206: Basics, Block Diagrams \& SFGs (GATE EC 2001)}
\addcontentsline{toc}{section}{Question 206: Basics, Block Diagrams \& SFGs (GATE EC 2001)}
\noindent \textbf{\color{primary}MCQ} \hfill \textbf{\color{secondary}1 Mark}


\noindent An electrical system and its signal-flow graph representations are shown in figure
(a) and (b) respectively. The values of $G_{2}$ and H, respectively are

\begin{center}\includegraphics[max width=0.85\linewidth]{images/img\_16.jpg}\end{center}

\begin{center}\includegraphics[max width=0.85\linewidth]{images/img\_16.jpg}\end{center}


\begin{enumerate}[label=(\Alph*)]
    \item $\frac{Z_{3}(s)}{Z_{2}(s)+Z_{3}(s)+Z_{4}(s)},\frac{-Z_{3}(s)}{Z_{1}(s)+Z_{3}(s)}$
    \item $\frac{-Z_{3}(s)}{Z_{2}(s)-Z_{3}(s)+Z_{4}(s)},\frac{-Z_{3}(s)}{Z_{1}(s)+Z_{3}(s)}$
    \item \textbf{(Correct)} $\frac{Z_{3}(s)}{Z_{2}(s)+Z_{3}(s)+Z_{4}(s)},\frac{Z_{3}(s)}{Z_{1}(s)+Z_{3}(s)}$
    \item $\frac{-Z_{3}(s)}{Z_{2}(s)-Z_{3}(s)+Z_{4}(s)},\frac{Z_{3}(s)}{Z_{1}(s)+Z_{3}(s)}$
\end{enumerate}

\noindent \textbf{Correct Answer:} (C)

\begin{tcolorbox}[solbox, title=Detailed Solution -- Question 206]
From KVL in both loops In first loop, 
$V_{i}(s)=I_{1}(s) Z_{1}(s)+\left[I_{1}(s)-I_{2}(s)\right] Z_{3}(s)$
$V_{i}(s)=I_{1}(s)\left[Z_{1}(s)+Z_{3}(s)\right]-I_{2}(s) Z_{3}(s)$

\[
\frac{V_{i}(s)}{Z_{1}(s)+Z_{3}(s)}=I_{1}(s)-\frac{I_{2}(s) \cdot Z_{3}(s)}{Z_{1}(s)+Z_{3}(s)} \; \; ...(i)
\]

 In second loop, 
$\left[I_{2}(s)-I_{1}(s)\right] Z_{3}(s)+I_{2}(s) Z_{2}(s)+I_{2}(s) Z_{4}(s)=0$
$I_{2}(s)\left(Z_{2}(s)+Z_{3}(s)+Z_{4}(s)\right)=I_{1}(s) \cdot Z_{3}(s)$
$G_{2}=\frac{I_{2}(s)}{I_{1}(s)}=\frac{Z_{3}(s)}{Z_{2}(s)+Z_{3}(s)+Z_{4}(s)}$
 From SFG, $I_{1}(s)=V_{i} G_{1}(s)+I_{2}(s) H(s)$

\[
I_{1}(s)=V_{i}(s) \frac{1}{Z_{1}(s)+Z_{3}(s)}+I_{2}(s) \frac{Z_{3}(s)}{Z_{1}(s)+Z_{3}(s)}
\]

 $\therefore \quad H(s)=\frac{Z_{3}(s)}{Z_{1}(s)+Z_{3}(s)}$
 (comparing above two equations)
\end{tcolorbox}

\vspace{1.5em}
\hrule
\vspace{1.5em}

\section*{Question 207: Frequency Response Analysis (GATE EC 2001)}
\addcontentsline{toc}{section}{Question 207: Frequency Response Analysis (GATE EC 2001)}
\noindent \textbf{\color{primary}MCQ} \hfill \textbf{\color{secondary}1 Mark}


\noindent The Nyquist plot for the open-loop transfer function G(s) of a unity negative
feedback system is shown in figure. If G(s) has no pole in the right half of splane, the number of roots of the system characteristic equation in the right half
of s-plane is

\begin{center}\includegraphics[max width=0.85\linewidth]{images/img\_1.jpg}\end{center}

\begin{center}\includegraphics[max width=0.85\linewidth]{images/img\_1.jpg}\end{center}


\begin{enumerate}[label=(\Alph*)]
    \item \textbf{(Correct)} 0
    \item 1
    \item 2
    \item 3
\end{enumerate}

\noindent \textbf{Correct Answer:} (A)

\begin{tcolorbox}[solbox, title=Detailed Solution -- Question 207]
$N=0$
 (1 encirclement in CW direction and other in ACW)
$P=0$(no pole in right half)
 So, $N=P-Z$
$Z=P-N=0$
$\therefore \quad$ No roots on RH of s-plane.
\end{tcolorbox}

\vspace{1.5em}
\hrule
\vspace{1.5em}

\section*{Question 208: Root Locus Technique (GATE EC 2001)}
\addcontentsline{toc}{section}{Question 208: Root Locus Technique (GATE EC 2001)}
\noindent \textbf{\color{primary}MCQ} \hfill \textbf{\color{secondary}1 Mark}


\noindent The root-locus diagram for a closed loop feedback system is shown in figure. The system is overdamped.

\begin{center}\includegraphics[max width=0.85\linewidth]{images/img\_67.jpg}\end{center}

\begin{center}\includegraphics[max width=0.85\linewidth]{images/img\_67.jpg}\end{center}


\begin{enumerate}[label=(\Alph*)]
    \item only if $0\leq K\leq 1$
    \item only if $1 < k < 5$
    \item only if K$>$5
    \item \textbf{(Correct)} if $0\leq K < 1\, or \, K > 5$
\end{enumerate}

\noindent \textbf{Correct Answer:} (D)

\begin{tcolorbox}[solbox, title=Detailed Solution -- Question 208]
For overdamping roots of characteristic equation
should lie on negative axis and be unequal.
\end{tcolorbox}

\vspace{1.5em}
\hrule
\vspace{1.5em}

\section*{Question 209: Time Response Analysis (GATE EC 2001)}
\addcontentsline{toc}{section}{Question 209: Time Response Analysis (GATE EC 2001)}
\noindent \textbf{\color{primary}MCQ} \hfill \textbf{\color{secondary}1 Mark}


\noindent If the characteristic equation of a closed-loop system is $s^{2}+2s+2=0$, then the system is


\begin{enumerate}[label=(\Alph*)]
    \item over damped
    \item critically damped
    \item \textbf{(Correct)} underdamped
    \item undamped
\end{enumerate}

\noindent \textbf{Correct Answer:} (C)

\begin{tcolorbox}[solbox, title=Detailed Solution -- Question 209]
\[
\begin{aligned} s^{2}+2 \xi \omega_{n}+\omega_{n}^{2} &=0 \\ 2 \xi\omega{n} &=2, \xi=\frac{1}{\omega_{n}} \\ \omega_{n} &=\sqrt{2}\\ \therefore \quad \xi&=\frac{1}{\sqrt{2}} \end{aligned}
\]

$\because \xi < 1,$ hence response will be underdamped.
\end{tcolorbox}

\vspace{1.5em}
\hrule
\vspace{1.5em}

\section*{Question 210: Basics, Block Diagrams \& SFGs (GATE EC 2001)}
\addcontentsline{toc}{section}{Question 210: Basics, Block Diagrams \& SFGs (GATE EC 2001)}
\noindent \textbf{\color{primary}MCQ} \hfill \textbf{\color{secondary}1 Mark}


\noindent The equivalent of the block diagram in figure is given in

\begin{center}\includegraphics[max width=0.85\linewidth]{images/img\_46.jpg}\end{center}

\begin{center}\includegraphics[max width=0.85\linewidth]{images/img\_46.jpg}\end{center}

\begin{center}\includegraphics[max width=0.85\linewidth]{images/img\_81.jpg}\end{center}

\begin{center}\includegraphics[max width=0.85\linewidth]{images/img\_81.jpg}\end{center}


\begin{enumerate}[label=(\Alph*)]
    \item a
    \item b
    \item c
    \item \textbf{(Correct)} d
\end{enumerate}

\noindent \textbf{Correct Answer:} (D)

\begin{tcolorbox}[solbox, title=Detailed Solution -- Question 210]
Take off points is moved $G_{2}$ so $H/G_{2}$
\end{tcolorbox}

\vspace{1.5em}
\hrule
\vspace{1.5em}

\end{document}